% ======================================================================
%  Kapustin–Witten 理论导读讲义（讲义体，非解经本）
%  编译：在 lecture_notes/ 下运行 xelatex main.tex 两遍
% ======================================================================
\documentclass[11pt,a4paper,fontset=fandol,twoside]{ctexart}
% ======================================================================
%  Kapustin–Witten 理论导读讲义 — 共享导言区
%  编译：xelatex main.tex（两遍），工作目录 lecture_notes/
% ======================================================================

% ---------- 数学 ----------
\usepackage{amsmath,amssymb,amsthm,mathtools}
\usepackage{mathrsfs}
\usepackage{bm}

% ---------- 排版 ----------
\usepackage[margin=2.35cm,top=2.5cm,bottom=2.5cm]{geometry}
\usepackage{graphicx}
\usepackage{booktabs,array,longtable,multicol,multirow}
\usepackage{enumitem}
\setlist{itemsep=2pt,topsep=3pt,parsep=0pt}
\usepackage{xcolor}
\usepackage{xspace}
\usepackage{fancyhdr}
\usepackage[colorlinks=true,linkcolor=blue!55!black,urlcolor=teal!70!black,citecolor=blue!55!black]{hyperref}

% ---------- 图 ----------
\usepackage{tikz}
\usetikzlibrary{arrows.meta,positioning,calc,shapes.geometric,fit,backgrounds,decorations.pathreplacing}
\usepackage{tikz-cd}

% ---------- 颜色 ----------
\definecolor{kwblue}{RGB}{28,66,120}
\definecolor{kwgray}{RGB}{105,105,105}
\definecolor{kwred}{RGB}{150,40,40}
\definecolor{kwgreen}{RGB}{30,95,60}
\definecolor{kwboxbg}{RGB}{247,246,240}

% ---------- 标题样式 ----------
\usepackage{titlesec}
\titleformat{\section}{\Large\bfseries\color{kwblue}}{第\,\thesection\,讲}{0.6em}{}
% switch to plain appendix headings (A, B, C ... instead of 第 N 讲)
\newcommand{\appendixstyle}{%
  \titleformat{\section}{\Large\bfseries\color{kwblue}}{\thesection}{0.8em}{}%
}
\titleformat{\subsection}{\large\bfseries\color{kwblue!85!black}}{\thesubsection}{0.6em}{}
\titleformat{\subsubsection}{\normalsize\bfseries\color{kwblue!70!black}}{\thesubsubsection}{0.5em}{}
\titlespacing*{\section}{0pt}{2.2ex plus .4ex}{1.2ex}
\titlespacing*{\subsection}{0pt}{1.6ex plus .3ex}{0.8ex}

% ---------- 页眉页脚 ----------
\pagestyle{fancy}
\fancyhf{}
\fancyhead[LE,RO]{\small\color{kwgray}\thepage}
\fancyhead[LO]{\small\color{kwgray}Kapustin--Witten 理论导读}
\fancyhead[RE]{\small\color{kwgray}\nouppercase{\leftmark}}
\renewcommand{\headrulewidth}{0.3pt}
\setlength{\headheight}{14pt}
\setlength{\emergencystretch}{2.2em}

% ---------- 定理类 ----------
\theoremstyle{definition}
\newtheorem{definition}{定义}[section]
\newtheorem{example}[definition]{例}
\newtheorem{remark}[definition]{评注}
\newtheorem{exercise}[definition]{习题}
\theoremstyle{plain}
\newtheorem{theorem}[definition]{定理}
\newtheorem{proposition}[definition]{命题}
\newtheorem{lemma}[definition]{引理}
\newtheorem{corollary}[definition]{推论}
\newtheorem{conjecture}[definition]{猜想}
\renewcommand{\proofname}{证明}
\renewcommand{\thedefinition}{\arabic{section}.\arabic{definition}}

% ---------- 讲义专用盒子 ----------
\usepackage[most]{tcolorbox}
\newtcolorbox{goals}{enhanced,breakable,colback=kwblue!4,colframe=kwblue!55,
  boxrule=0.5pt,left=6pt,right=6pt,top=4pt,bottom=4pt,
  fonttitle=\bfseries\small,title={本讲目标}}
\newtcolorbox{follow}{enhanced,breakable,colback=kwgray!6,colframe=kwgray!45,
  boxrule=0.4pt,left=6pt,right=6pt,top=3pt,bottom=3pt,
  fonttitle=\bfseries\small,title={跟读提示（原文页段）}}
\newtcolorbox{keybox}{enhanced,breakable,colback=yellow!6,colframe=yellow!45!orange,
  boxrule=0.5pt,left=6pt,right=6pt,top=4pt,bottom=4pt,
  fonttitle=\bfseries\small,title={本讲核心}}
\newtcolorbox{checkbox}{enhanced,breakable,colback=kwgreen!4,colframe=kwgreen!45,
  boxrule=0.4pt,left=6pt,right=6pt,top=3pt,bottom=3pt,
  fonttitle=\bfseries\small,title={自检问题}}

% ---------- 数学宏 ----------
\newcommand{\MH}{\mathcal{M}_{H}}
\newcommand{\MC}{\mathcal{M}(G,C)}
\newcommand{\Ldual}{{}^{L}G}
\newcommand{\LdualC}{{}^{L}G_{\mathbb{C}}}
\newcommand{\gC}{\mathfrak{g}_{\mathbb{C}}}
\newcommand{\CC}{\mathbb{C}}
\newcommand{\PP}{\mathbb{P}}
\newcommand{\ZZ}{\mathbb{Z}}
\newcommand{\RR}{\mathbb{R}}
\newcommand{\HH}{\mathbb{H}}
\newcommand{\dd}{\mathrm{d}}
\newcommand{\ii}{\mathrm{i}}
\newcommand{\ee}{\mathrm{e}}
\newcommand{\Tr}{\operatorname{Tr}}
\newcommand{\Dmod}{\mathcal{D}}
\newcommand{\kw}{K\!W}
\newcommand{\src}[1]{\textup{\textsf{[#1]}}}
\newcommand{\term}[2]{#1（#2）}

% ---------- 引用键（附录 C 提供 bibitem） ----------
% \cite{kw06} \cite{w08} \cite{fr09} \cite{gw06} \cite{bz04} \cite{set13} \cite{ht03}


\title{\bfseries Kapustin--Witten 理论：几何朗兰兹的物理入门\\[6pt]
\large 一份十讲的导读讲义}
\author{威腾讨论班}
\date{2026 年 11 月（v1.0）}

\begin{document}

\maketitle
\thispagestyle{empty}

% 前言与导读（第 0 部分，不编号进「讲」）
\section*{写在前面：这份讲义怎么用}
\addcontentsline{toc}{section}{写在前面：这份讲义怎么用}
\markboth{写在前面}{写在前面}

\subsection*{一句话}

Kapustin--Witten 2006 \cite{kw06} 做的是这样一件事：把四维 $\mathcal N=4$ 超对称 Yang--Mills 理论做一次特定的拓扑扭化，
再把它紧化到一条 Riemann 面上；于是一个纯四维的物理问题，在低能下变成一个以 Hitchin 模空间 $\MH$ 为目标的二维 sigma 模型。
在这个二维模型里，四维的电磁对偶（S-对偶）恰好变成二维的镜像对称，而四维的 Wilson/'t~Hooft 线算子恰好变成作用在膜上的算子。
把这些对应一条一条接起来，就得到几何 Langlands 纲领的一份「物理推导」。

这份讲义不替代原文，也不替代我们那本 219 页的解经本。它的目标只有一个：
\textbf{让你能自己翻开 \cite{kw06}，从第一页读到 §12，知道每一节在做什么。}

\subsection*{三条读法}

\begin{center}
\begin{tabular}{@{}llp{8.6cm}@{}}
\toprule
路线 & 时间 & 怎么读 \\
\midrule
快览 & 约 3 小时 & 读第 0 部分 + 第 1、4、6 讲 + 每讲的「本讲核心」与「跟读提示」，跳过全部推导。 \\
标准 & 约 3 天 & 十讲顺序读，每讲后做附录 B 的 3 道习题，再按「跟读提示」回原文对应页段。 \\
深入 & 约 2 周 & 十讲 + 每讲的自检问题全部写答案 + 原文 §2--§12 逐节读 + 第 8 讲讲到的后续文献。 \\
\bottomrule
\end{tabular}
\end{center}

\subsection*{需要什么前置}

\begin{itemize}
\item \textbf{物理}：路径积分与规范场（会写 $\int \mathcal D A\, e^{\ii S}$ 就行）、超对称的基本语言（超荷、多重态）、
      知道 $\mathcal N=4$ 超对称 Yang--Mills 是什么。拓扑扭化的技术细节第 3 讲现讲。
\item \textbf{几何}：复流形与近复结构、向量丛与联络、层与凝聚层的入门概念。
      Hitchin 模空间、超凯勒商这些第 1 讲现讲；D-模在第 7 讲现讲。
\item \textbf{不需要}：不需要先学几何 Langlands 的完整理论，也不需要懂导出代数几何。
\end{itemize}

\subsection*{记号约定}

\begin{center}
\begin{tabular}{@{}ll@{}}
\toprule
记号 & 含义 \\
\midrule
$G$ & 紧单连通单李群（主要例子 $SU(N)$）；$\gC$ 是其复化李代数 \\
$\Ldual$ & Langlands 对偶群（$\LdualC$ 是其复化）；第 2 讲解释它为什么出现在这里 \\
$C$ & 亏格 $g\ge 2$ 的紧 Riemann 面（几何 Langlands 的「曲线」） \\
$\Sigma$ & 紧化用的另一条曲面；四维理论放在 $\Sigma\times C$ 上 \\
$\MH$ & $C$ 上 $G$-Higgs 丛的模空间（Hitchin 模空间），三个复结构 $I,J,K$ \\
$\tau$ & 复耦合 $\tau=\theta/2\pi+4\pi\ii/g_{\rm YM}^2$；$\Psi$ 是扭化后的典范参数 \\
$^L\!R$、$R$ & $\Ldual$ 与 $G$ 的表示；$\mathcal O(x)$、$K_C$ 是 $C$ 上的线丛与典范丛 \\
\src{B p.148} & 原文（\cite{kw06}）的印刷页锚点；本讲义所有断言都尽量给到页 \\
\bottomrule
\end{tabular}
\end{center}

\subsection*{十讲地图}

\begin{center}
\small
\begin{tabular}{@{}clp{5.4cm}l@{}}
\toprule
讲 & 主题 & 一句话 & 原文页段 \\
\midrule
1 & 舞台 & 一条曲线、一个群、一个模空间 $\MH$ & §4, pp.38--56 \\
2 & 四维主角 & $\mathcal N=4$ SYM、S-对偶、$\Ldual$、Table 1 & §1--§2, pp.2--16 \\
3 & 扭化 & 为什么要把超对称扭掉；一族 TFT 与 $\Psi$ & §3, pp.16--38 \\
4 & 紧化 & 四维的桥塌成二维的镜像对称 & §5, pp.56--78 \\
5 & 膜与线算子 & Wilson/'t~Hooft 变成作用在膜上的算子 & §6--§8, pp.78--123 \\
6 & 核心等式 & 't~Hooft 算子 $=$ Hecke 算子 & §9--§10, pp.123--184 \\
7 & 翻译器 & 从 A-膜到 D-模，得到 Hecke 特征层 & §11, pp.184--207 \\
8 & 收获与边界 & 证明了什么、没证明什么、之后读什么 & §12, pp.207--222 \\
\bottomrule
\end{tabular}
\end{center}

\subsection*{三条诚实提醒}

\begin{enumerate}
\item \textbf{这不是数学证明。} 原文是物理论证：给出对应、给证据、给一致性检查，把「为什么应该如此」讲清楚，
      但没有证明几何 Langlands。严格的证明要等到 2024 年前后 Gaitsgory 等人的纯代数几何工作（第 8 讲会说）。
\item \textbf{只做非分歧情形。} 原文只处理 $C$ 上无分歧（unramified）的情形；分歧与野分歧是 Gukov--Witten \cite{gw06} 的主题。
\item \textbf{物理侧的很多步骤是「应该」。} 原文里 ``we expect''、``should be'' 这类措辞不是客套，
      是作者在标明「这一步只有物理直觉，没有逐条验证」。本讲义会用「原文只写应该」这样的话提醒你。
\end{enumerate}

\subsection*{和另外两份材料的关系}

我们还有两份配套材料，读者可以按需取用：
\begin{itemize}
\item 一份 219 页的\textbf{解经本}（对原文逐节深读，含四层：骨架/结构/骨血/细度，附 nLab 注释与 60 条原文问答）；
\item 一次\textbf{讨论班}的八轮报告与 112 条语料（八个人每节一人报告、三席质询、独立复算）。
这份讲义是把那两份材料里「最该先知道的东西」压成十讲——所以它比解经本浅，但更适合第一次读原文。
\end{itemize}


\tableofcontents
\newpage

% ============================================================
%  Lecture 1: The stage (curve, group, moduli space)
% ============================================================

\section{舞台：一条曲线、一个群、一个模空间}
\label{sec:stage}

\begin{goals}
\begin{enumerate}
\item 说清楚这十讲的主角为什么是一个三元组 $(C,G,\MH)$，其中每个字母各管什么。
\item 会写 Hitchin 方程与 Hitchin 对 $(E,\phi)$ 的定义，并能解释 $\phi$ 为什么既是“场”又是“一次形式”。
\item 能说清 $\MH$ 的三个复结构、维数公式，以及 Hitchin 纤维化为什么让二维 sigma 模型变得可解。
\end{enumerate}
\end{goals}

\subsection{为什么是一条曲线和一个群}

先把结论放在这里：后面十讲反复出现的那个对象，是一个二维 sigma 模型的 target，而它由两样东西决定——一条紧 Riemann 面 $C$，和一个紧 Lie 群 $G$。这两样东西分别来自几何与物理，本讲其余部分都在解释它们怎么合到一起。

几何那头，主角是 $C$ 上的 $G$-丛，等价地是 $\pi_1(C)\to G$ 的表示。几何 Langlands 纲领说的是：$C$ 上 $G$ 侧的对象与 $\Ldual$ 侧的对象之间有一一对应。这个陈述本身是纯几何的，与物理无关；\kw{} 做的事情（\cite{kw06}）是给它一个物理来源。

物理那头，四维 $\mathcal N=4$ 超对称 Yang--Mills 理论先做一次特殊的拓扑扭化（第 3 讲详述），再放到乘积 $M=\Sigma\times C$ 上。从 $\Sigma$ 的方向看，低能理论是一个二维 sigma 模型，target 就是 $\MH$。原文的说法是：四维扭化理论 on $\Sigma$ reduces to a sigma-model of maps $\Phi:\Sigma\to \MH$ \src{B p.38}。这句话是整份讲义的起点，你把它记住，理由后面都会补上。

为什么 $C$ 要紧、而且要 $g\ge 2$？有两个理由。第一，$g\ge2$ 时半单 $G$ 上的平坦联络一般不可约，$\MH$ 的一般点也就光滑，可约解给出的奇异点可以被躲开 \src{B p.39}。第二，$C$ 紧才使后面要用的紧性成立：Hitchin 纤维化的纤维是紧的复环面，“纤维紧”这一条在原文里反复出现。反过来，若 $C$ 是球面，$\pi_1$ 平凡；若是环面，$2g-2=0$，维数公式直接退化。所以 $g\ge2$ 不是技术性的偷懒，而是舞台本身的要求。

$\mathcal N=4$ 这个“四”也不是随意的：它给出十六个超荷，正好是四维无引力规范理论能带的最大超对称（第 2 讲细说）。而 $G$ 要半单、要紧：四维规范理论要酉性，规范群就得紧；几何描述里到处用复化 $G_\CC$，所以原文一开头就声明 always assume $G$ to be compact and denote its complexification as $G_\CC$ \src{B p.6}。至于半单——若 $G$ 含 $U(1)$ 因子，每个解都带连续的 $U(1)$ 规范对称，低能描述会变成 sigma 模型乘一个超对称 $U(1)$ 规范理论，原文明确说自己只在半单情形工作 \src{B p.40}。这两条约定你先收下，本讲不再提。

\subsection{丛、Higgs 场，以及 $\phi$ 的两个身份}

定义从一个 $G$-丛 $E\to C$ 开始，$A$ 是它的联络。这里有一个二维特有的魔术：在复维 $1$ 上，$A$ 的 $(0,1)$ 部分 $\bar\partial_A$ 自动满足 $\bar\partial_A^2=0$——因为 $(0,2)$-形式在曲线上不存在。于是任意联络都把 $E$ 变成全纯丛。原文一句话带过：any connection on a smooth $G$-bundle over a Riemann surface automatically turns $E$ into a holomorphic $G$-bundle \src{B p.49}。高维没有这种好事，所以“联络即全纯结构”是这一讲的第一个礼物，也是为什么整套构造只在曲线上成立。

第二个对象是 \term{Higgs field}{Higgs 场} $\phi$：它是 $K_C\otimes\mathrm{ad}(E)$ 的全纯截面，也就是取值在伴随丛里、带一个 $(1,0)$ 一次形式的全纯对象 \src{B p.49}。把 $E$ 连同 $\phi$ 一起叫做 \term{Higgs bundle}{Higgs 丛}。$g\ge2$ 时 $H^0(C,K_C\otimes\mathrm{ad}E)$ 非零，所以这不是空定义。

现在说本讲第一个要点：$\phi$ 有两个身份，缺一个都读不下去。

\textbf{身份一，它是场。}四维理论里有六个实标量，扭化之后其中四个被重新解释为四维流形上的一次形式；紧化到 $C$ 上，它的 $(1,0)$ 部分就是 $\phi$ \src{B p.56}。所以 $\phi$ 不是标量，也不是凭空造出来的几何对象——它来自四维场论。这个身份在第 3 讲扭化时会再出现一次，那时你会知道“为什么恰好是四个标量、为什么恰好是一次形式”。

\textbf{身份二，它是余切方向。}稳定 $G$-丛的模空间记作 $M$（原文的 $\mathcal M$），它在 $E$ 处的切空间是 $H^1(C,\mathrm{ad}E)$；由 Serre 对偶，余切空间是 $H^0(C,K_C\otimes\mathrm{ad}E)$。而 $\phi$ 正好住在这里。于是 $\phi$ 是“$M$ 的一个余切方向”：$\phi=0$ 的切片就是 $M\subset\MH$，更一般地 $\MH$ 包含 $T^*M$ 作为开稠密子集（差一个高余维的坏集）\src{B p.52}。

两个身份合起来，图像是“一个丛，加一个余切方向”。这也预告了下一节：三个复结构里，总有一个把 $\phi$ 当几何数据（Higgs 丛），另一个把 $\phi$ 与 $A$ 混成复联络（平坦联络）。

这里容易卡住：不要把 $\phi$ 当标量。一旦当标量，复结构 $J$ 里“$A+\ii\phi$ 是复联络”的写法立刻失去意义——$A$ 是一次形式，$\phi$ 必须也是一次形式，两者才能相加 \src{B p.51}。

\subsection{三个复结构，以及为什么存在两种商}

把所有的 $(A,\phi)$ 收进一个无穷维仿射空间 $W$。它带一个平坦度规（原文 (4.2)，只需知道它把 $\Tr(\delta A\wedge*\delta A+\delta\phi\wedge*\delta\phi)$ 在 $C$ 上积分）\src{B p.41}。$W$ 上可以定义三个复结构 $I,J,K$；用“哪些线性函数全纯”来定义，比写矩阵清楚：

\begin{itemize}
\item $I$：全纯坐标是 $A^{0,1}$ 与 $\phi^{1,0}$，即 $A$ 的 $(0,1)$ 部分与 $\phi$ 的 $(1,0)$ 部分全纯；
\item $J$：全纯坐标是 $A+\ii\phi$，即 $A$ 与 $\ii\phi$ 的各分量一起；
\item $K$：全纯坐标是 $A_z-\phi_z$ 与 $A_{\bar z}+\phi_{\bar z}$。
\end{itemize}

这三个结构由四元数关系 $IJ=-JI=K$ 等联系在一起。原文把它们推广成一整族复结构 $I_w$，参数 $w\in\CC\PP^1_h$：$w=0$ 是 $I$，$w=\infty$ 是 $-I$，$w=-\ii$ 是 $J$，$w=1$ 是 $K$ \src{B p.43 eq 4.6}。另有一个 $U(1)$ 作用 $\phi\mapsto\lambda\phi$（$|\lambda|=1$）把 $I_w$ 送到 $I_{\lambda^{-1}w}$，并固定 $w=0,\infty$ 两点 \src{B p.43}。请记住：$w$ 只是“$\MH$ 上复结构族”的参数，不是理论参数；理论参数 $\Psi$ 要到第 3 讲才出现，把这两个混起来是最常见的误读。

三个复结构各自配一个辛形式 $\omega_I,\omega_J,\omega_K$，以及 $\Omega_I=\omega_J+\ii\omega_K$ 等 \src{B p.44 eq 4.8}。这里有一个不对称，值得停下来看一眼：$\omega_I$ 与 $\omega_K$ 可以用微分形式整体写出（不含 $C$ 的复结构），而 $\omega_J$ 不行；对 $\Omega$ 来说关系正好反过来——$\Omega_I,\Omega_K$ 依赖 $C$ 的复结构，$\Omega_J$ 不依赖。这个不对称在讲镜像与 S-对偶时是主角，现在只要记住它存在。

现在做商。规范群作用在 $W$ 上，保持这三个辛形式；它给出的三个矩映射，其零点正好是 Hitchin 方程的三条：
\[
\mu_I=\mu_J=\mu_K=0
\quad\Longleftrightarrow\quad
F-\phi\wedge\phi=0,\qquad D\phi=D^*\phi=0 ,
\]
这两边分别是原文的 (4.11) 与 (4.1) \src{B p.46, B p.39}。于是
\[
\MH=W/\!/\!/\,G
\]
是一个 \term{hyper-Kähler quotient}{超凯勒商} \src{B p.46}。这个名词是字面意思：三个矩映射同时为零再商，商出来的流形自动带三个复结构——所以“超凯勒”不是额外假设，而是商法的产物。

这里有一个值得花五分钟想的区别：为什么不能只做一个复商？原文把两种商摆在一起看 \src{B p.49}。设 $\nu_I=\mu_J+\ii\mu_K$ 是复结构 $I$ 下的全纯矩映射，$\nu_I=0$ 就是 $\bar\partial_A\phi=0$（即 $\phi$ 全纯）。在它的零点集上有两个选择：

\begin{enumerate}
\item[(I)] 再要求实矩映射 $\mu_I=0$，只商 $G$-值的规范变换；
\item[(II)] 直接商 $G_\CC$-值的规范变换。
\end{enumerate}

(I) 给出带三个复结构的超凯勒流形；(II) 给出的只是复流形，而且只在当前这个复结构里有意义。两者相差一个高余维的稳定化修正，但在“看哪个复结构”这件事上完全不同：(II) 在 $I$ 下把对象变成稳定 Hitchin 对 $(E,\phi)$，这是 Hitchin 的定理 \src{B p.50}；(II) 在 $J$ 下把对象变成平坦 $G_\CC$ 联络，亦即 $\pi_1(C)\to G_\CC$ 的稳定同态，这是 Corlette--Donaldson 的定理 \src{B p.51}。而所有 $w\ne0,\infty$ 的 $I_w$ 给出的是同一个空间 $Y$ \src{B p.52}。所以“同一个模空间，三种复结构”不是比喻，而是两条定理拼起来的事实。表 \ref{tab:three-I} 把它压成一张表。

\begin{table}[htbp]
\centering
\small
\begin{tabular}{@{}llll@{}}
\toprule
复结构 & 全纯坐标 & 矩映射为零 & 商 $G_\CC$ 后的对象 \\
\midrule
$I$（$w=0$） & $A^{0,1},\ \phi^{1,0}$ & $\bar\partial_A\phi=0$ & 稳定 Hitchin 对 $(E,\phi)$ \\
$J$（$w=-\ii$） & $A+\ii\phi$ & $F(A+\ii\phi)=0$ & 稳定平坦 $G_\CC$ 联络 \\
$K$（$w=1$） & $A_z-\phi_z,\ A_{\bar z}+\phi_{\bar z}$ & 复联络 $(A_z-\phi_z,A_{\bar z}+\phi_{\bar z})$ 平坦 & 同一个 $Y$ \\
\bottomrule
\end{tabular}
\caption{$\MH$ 的三个复结构：同一个点集，三种全纯描述。前两行分别是 Hitchin 定理与 Corlette--Donaldson 定理，出处见 \src{B pp.50--52}。}
\label{tab:three-I}
\end{table}

\subsection{维数、Hitchin 纤维化，以及“为什么可解”}

先记结论：$\MH$ 是复维 $(2g-2)\dim G$ 的复流形 \src{B p.39}。本讲不证。直觉是：每多一个把手，就多 $2\dim G$ 个复方向；$g=2$、$G=SU(2)$ 时是 $6$，而稳定丛模空间 $M$ 是 $3g-3=3$ 维，$T^*M$ 正好 $6$ 维，与上一节“$\MH\supset T^*M$”对得上。原文也只是引用 Hitchin 的结果，不在这里重算。

真正让人兴奋的是 \term{Hitchin fibration}{Hitchin 纤维化}。取 $\phi$ 的 $G$-不变多项式 $P_i$，$i=1,\dots,r$（$r=\mathrm{rank}\,G$），次数 $d_i$ 满足
\[
\sum_{i=1}^{r}(2d_i-1)=\dim G .
\]
这些多项式给出映射
\[
\pi:\ \MH\longrightarrow B:=\bigoplus_{i=1}^{r}H^0\bigl(C,K_C^{\,d_i}\bigr),
\qquad (E,\phi)\longmapsto\bigl(P_1(\phi),\dots,P_r(\phi)\bigr),
\]
公式 (4.24) 见 \src{B p.53}。$SU(2)$ 时 $B=H^0(C,K_C^2)\cong\CC^{3g-3}$，映射就是 $\phi\mapsto\Tr\phi^2$ \src{B p.53}。因为 $\dim H^0(C,K_C^{\,d})=(2d-1)(g-1)$，从上面的恒等式直接算出
\[
\dim B=(g-1)\dim G=\tfrac12\dim \MH ,
\]
所以一般纤维的复维数也是 $(g-1)\dim G$ \src{B p.54 eq (4.25)}。原文还给了一个漂亮的计数：$SU(2)$ 时对固定 $w\in B$，方程 $\Tr\phi^2=w$ 是 $3g-3$ 个二次方程对 $3g-3$ 个未知数，一般有 $2^{3g-3}$ 个解 \src{B p.54}。

接下来是 Hitchin 的两条主定理。第一条：$\MH$ 在复结构 $I$ 下是 \term{completely integrable system}{完全可积系统}——存在 $\tfrac12\dim\MH$ 个 $I$-全纯、代数独立、对 $\Omega_I$ 的 Poisson 括号两两交换的函数。它们可以取成 $B$ 上线性函数的拉回，显式写法是
\[
H_{\alpha,i}=\int_C \alpha\wedge P_i(\phi),\qquad \alpha\in H^1\bigl(C,K_C^{\,1-d_i}\bigr),
\]
见 \src{B p.55 eq (4.26)}。交换性有一个一句话证明：在 $W$ 上 $\phi_z$ 只与 $A_{\bar z}$ 共轭、与 $\phi_z$ 自己交换，而 $H_{\alpha,i}$ 只是 $\phi$ 的函数 \src{B p.55}。第二条：这些对易流给出 $I$-全纯的交换流，而复环面恰好承认交换流；再由全纯辛几何的一般结果与纤维紧性，得到一般纤维是复环面 \src{B p.55}。原文还补了一句本单元最值钱的结论：纤维由 $H_k=h_k$ 定义，维数是一半、并且 $\Omega_I$ 在其上各向同性，所以纤维是 Lagrange 子流形 \src{B p.56}。

顺带说一句：$(2g-2)\dim G$ 与 $(g-1)\dim G$ 这两个数会反复出现，因为 $\MH$ 的“一半”在底 $B$ 上，“一半”在纤维里。以后看到“一半维数”的说法，回到这一页。

最后回答本讲标题里的问题：这件事为什么让二维 sigma 模型可解？当一个 sigma 模型的 target 是“底加紧环面纤维”的纤维化时，在纤维方向做 T-对偶几乎不改变物理，只是把周期取倒数。于是同一个二维理论可以在一个复结构里看成 $B$ 模型、在另一个里看成 $A$ 模型，两者互为镜像；而 $G$ 与 $\Ldual$ 的 Hitchin 系统恰好互为镜像 \cite{w08}。这正是原文第 5 讲“四维电磁对偶降成二维 T-对偶”的机制。

但这里必须诚实，这是本讲最容易过度宣称的地方：原文 §4.3 只造好了 $G$ 这一侧——纤维是复环面、纤维 Lagrange、纤维化存在，这些都是本节原话；而 T-对偶、镜像这些词在 §4 里根本没出现，它们要到 §5.3 才第一次出场 \src{B p.70}。所以准确的说法是：§4 给的是“可解的结构”，而不是“已经解了”。你读原文 §4 时，不要提前把它当成推导已完成。

\begin{figure}[htbp]
\centering
\begin{tikzpicture}[>=Stealth, font=\small]
  % 底 B
  \draw[kwblue,thick] (-4.0,0) -- (4.0,0);
  \node[kwblue,anchor=west] at (4.05,0) {$B$};
  % 一般纤维：环面画成带孔的椭圆
  \foreach \x in {-2.9,-1.3}
  {
    \draw[kwred,thick] (\x,1.85) ellipse (0.6 and 0.4);
    \draw[kwred,thick] (\x,1.85) ellipse (0.2 and 0.13);
    \draw[->,kwgray] (\x,0.12) -- (\x,1.35);
    \fill[kwblue] (\x,0) circle (1.2pt);
  }
  \node[font=\scriptsize,kwred] at (-2.1,2.55) {一般纤维：复环面（Lagrange）};
  % 退化纤维：8 字形
  \draw[kwred,thick] (1.7,2.2) circle (0.34);
  \draw[kwred,thick] (1.7,1.52) circle (0.34);
  \draw[->,kwgray] (1.7,0.12) -- (1.7,1.15);
  \draw[kwred,fill=white] (1.7,0) circle (2pt);
  \node[font=\scriptsize,kwred] at (1.7,2.75) {判别式轨迹};
  % 维数标注
  \node[font=\footnotesize,anchor=west] at (-3.95,-0.55) {底 $B$：复维 $(g-1)\dim G$};
  \node[font=\footnotesize,anchor=west] at (0.75,-0.55) {纤维：复维 $(g-1)\dim G$};
  \node[left=2pt,font=\scriptsize] at (-1.28,0.7) {$\pi$};
\end{tikzpicture}
\caption{Hitchin 纤维化 $\pi:\MH\to B$。一般纤维是复环面，并且在 $\Omega_I$（等价地 $\omega_J,\omega_K$）下是 Lagrange 子流形；在判别式轨迹上纤维退化。底与纤维各占一半复维数，这是后文“可解”的全部几何来源。}
\label{fig:fibration}
\end{figure}

\begin{keybox}
$\MH=W/\!/\!/\,G$ 是无穷维仿射空间按三个矩映射（即 Hitchin 方程）做的超凯勒商；它同时带三个复结构、复维 $(2g-2)\dim G$，并被 Hitchin 纤维化成底为 $B$（复维一半）、纤维为复环面（Lagrange）的可积系统。二维 sigma 模型的可解性就建立在这“一半环面加一半对易哈密顿量”的几何上——而对偶那一半要到 §5 才补上。
\end{keybox}

\begin{follow}
\begin{itemize}
\item \textbf{B §4 开头两页，pp.38--39}：紧化、极小解、Hitchin 方程 (4.1) 与维数公式。这是全讲的出发点，值得先读三遍。
\item \textbf{B §4.1，pp.41--44}：$W$ 的度规 (4.2)、$I,J,K$ 的矩阵定义 (4.3)--(4.5)、$I_w$ 家族 (4.6)。不必记矩阵，记住“哪个量依赖 $C$ 的复结构”。
\item \textbf{B §4.2，pp.49--52}：两种商 (I)/(II) 与两条定理（Hitchin、Corlette--Donaldson）。本讲第 3 节的内容基本都在这一段里。
\item \textbf{B §4.3，pp.52--56}：纤维化 (4.24)--(4.25)、完全可积 (4.26)、纤维是 Lagrange。p.55 的“复环面承认交换流”与 p.56 的 Lagrange 论证，是本节最值钱的两句。
\end{itemize}
\end{follow}

\begin{checkbox}
\begin{itemize}
\item 写出 Hitchin 方程的三条，并说明它们分别是哪三个矩映射的零点。
\item 复结构 $I$ 与 $J$ 下，把全纯矩映射设为零再商 $G_\CC$，分别得到什么几何对象？两条定理各叫什么名字？
\item 为什么 $\dim_\CC\MH=(2g-2)\dim G$ 而 $\dim B=(g-1)\dim G$？两者相差的那一半落在哪里？
\item “纤维是复环面”与“纤维是 Lagrange”各用了什么条件？为什么后者同时是对 $\omega_J$ 与 $\omega_K$ 的陈述？
\item 本讲中哪一句只说明“结构已经造好”，哪一句才是“对偶已经做出来”？§4.3 里有没有 T-对偶？
\end{itemize}
\end{checkbox}

\section{四维主角：$\mathcal{N}=4$ 超对称 Yang--Mills 与 S-对偶}
\label{sec:physics}

\begin{goals}
\begin{enumerate}
\item 记住 $\mathcal N=4$ SYM 的场内容与十六个超荷，知道它为什么是四维最大的无引力规范理论。
\item 会写复耦合 $\tau$，能一句话说清 S-对偶到底断言了什么。
\item 会读 Table 1，能解释 GNO 对偶群为什么正是 $\Ldual$，并分清这一讲里哪些是定理、哪些只是证据。
\end{enumerate}
\end{goals}

\subsection{从十维降下来：场内容与十六个超荷}

上一讲我们把四维量子场论当成一台黑箱用了：它在 $\Sigma\times C$ 上紧化，给出 target 为 $\MH$ 的二维 sigma model。这一讲把黑箱打开，顺便请出全篇第二个主角——S-对偶。

造 $\mathcal N=4$ SYM 最省事的办法是从十维降维。为什么偏偏是十维？Nahm 定理说，十维是超对称 Yang--Mills 能存在的最大维数；再往上，超荷作用会带进自旋 $3/2$ 的粒子，那已经是超引力 \src{B p.6}。十维的场只有两个：规范场 $A$（$G$-丛上的联络）与一个正手征旋量 $\lambda$（$S_+\otimes\mathrm{ad}(E)$ 的截面）\src{B p.7}。约定上 $A$ 与 $\lambda$ 取实李代数、在酉表示中反厄米，即 $A=-\ii A'$，与通常物理约定差一个 $\ii$，这样写是为了迁就数学文献 \src{B p.7 eq (2.2)}。$\Tr$ 是不变负定二次型，归一化到 $(1/8\pi^2)\int_{S^4}\Tr F\wedge F$ 取任意整数 \src{B p.8}。

十维的超对称由十六个守恒超荷 $Q_a$（$a=1,\dots,16$）实现，它们取值在 $S_-$ 里，满足原文 (2.9)：
\[
\{Q_a,Q_b\}=\sum_{I=1}^{10}\Gamma^I_{ab}\,P_I ,
\]
并且 $Q$ 与动量 $P_I$ 交换、按 $S_-$ 变换 \src{B p.9}。降维的手续是让所有场都不依赖 $x^4,\dots,x^9$。于是 $A$ 的四个分量给四维规范场 $A_\mu$，剩下六个分量给六个实标量 $\phi_i=A_{i+4}$（$i=0,\dots,5$）。十维曲率按指标落在哪个方向，自动分裂成 $F_{\mu\nu}$、$D_\mu\phi_i$、$[\phi_i,\phi_j]$ 三类，玻色部分就长成
\[
I_4=\frac{1}{e^2}\int d^4x\ \Tr\Bigl(\tfrac12 F_{\mu\nu}F^{\mu\nu}
+\sum_i D_\mu\phi_i D^\mu\phi_i+\tfrac12\sum_{i,j}[\phi_i,\phi_j]^2\Bigr)
\]
（原文 (2.10)，\src{B p.9}）。三块各有读法：规范场动能；六个标量的动能；以及标量之间的四次势——最后一项是正定平方和，这正是 $\mathcal N=4$ 有平坦势方向、能一直保超对称的原因。若再要求拉氏量在曲率上是二次的，那么给定 $G$ 之后它唯一到几个耦合常数，没有别的选择 \src{B p.6}。

费米子有几个？十维的 $\lambda$ 手征，降维后在 $Spin(1,3)\times Spin(6)$ 下分解。前一个因子是四维洛伦兹群，后一个因子满足 $Spin(6)\cong SU(4)$，这就是 \term{R-symmetry}{R-对称} 群 $SU(4)_R$。手征条件把十六维旋量表示按
\[
(2,1,4)\oplus(1,2,\bar 4)
\]
分解（原文 (2.14)，\src{B p.10}）。两块互为复共轭，加起来实维数仍是十六；四维超代数于是写成原文 (2.15)：
\[
\{Q^{AX},Q^Y_{\dot A}\}=\delta^X_Y\sum_{\mu=0}^{3}\Gamma^\mu_{A\dot A}P_\mu,
\qquad \{Q,Q\}=\{\bar Q,\bar Q\}=0,
\]
其中 $A,\dot A$ 是四维旋量指标，$X,Y=1,\dots,4$ 是 $SU(4)_R$ 指标，超荷总数 $4\times 4=16$ \src{B p.11}。同一个分解也作用在 $\lambda$ 上，所以费米子是四个 Weyl 费米子。场内容这一行就记全了：一个规范场 $A_\mu$、四个 Weyl 费米子、六个实标量。

顺便说一句：如果你觉得 1+4+6 这个组合眼熟，那不是巧合——它是十维 $10=1+4+6$ 的指标分裂，也是 $SU(4)_R$ 的 $\mathbf 4$ 加 $\bar{\mathbf 4}$ 拼出六个实标量的原因：$\mathbf 6=\mathbf 4\wedge\mathbf 4$。

\subsection{复耦合 $\tau$ 与 S-对偶}

$\mathcal N=4$ SYM 有两个经典耦合：规范耦合 $e$ 与 $\theta$ 角。$\theta$ 项是拓扑项，等于 $\theta$ 乘以第二陈数，也就是瞬子数 \src{B p.10}。量子力学只看作用量模 $2\pi$ 的整数倍，而 $\frac{1}{8\pi^2}\int\Tr F\wedge F$ 是整数，所以 $\theta$ 是周期变量 $\theta\sim\theta+2\pi$。把两者合成一个复数
\[
\tau=\frac{\theta}{2\pi}+\frac{4\pi\ii}{e^2},
\]
则因 $\Im\tau>0$，$\tau$ 落在上半平面 $\HH$ 里（原文 (2.12)，\src{B p.10}）。$SL(2,\RR)$ 按 $\tau\mapsto(a\tau+b)/(c\tau+d)$、$ad-bc=1$ 作用。其中 $T:\tau\mapsto\tau+1$ 只是把 $\theta$ 平移 $2\pi$，在任何 $G$ 下都是经典对称 \src{B p.11}。

真正的重头是 $S$。原文的一句话是：\ldots an additional quantum symmetry that inverts $\tau$, exchanges $G$ with $\Ldual$, and exchanges electric and magnetic charges \src{B p.11}。请盯住三个词：\emph{additional}（它不是已有的经典对称）、\emph{quantum}（只在量子理论里成立）、\emph{inverts}（把耦合倒过来）。$S$ 与 $T$ 一起生成 $SL(2,\RR)$ 的一个离散子群 $\Gamma$。$G$ 属 simply-laced 型时，$S$ 就是原文 (2.16)：
\[
S=\begin{pmatrix}0&1\\-1&0\end{pmatrix},\qquad
\tau\mapsto-\frac{1}{\tau},
\]
与 $T$ 生成 $SL(2,\ZZ)$ \src{B p.11}。$G$ 非 simply-laced 时要稍作调整：
\[
\tau\mapsto-\frac{1}{n_g\tau},\qquad
n_g=(|\alpha_{\mathrm{long}}|/|\alpha_{\mathrm{short}}|)^2 ,
\]
即 $n_g=2$（$F_4$）、$n_g=3$（$G_2$），由原文 (2.17) 的矩阵实现 \src{B p.11}。$\tau$ 写成 $4\pi\ii/e^2$ 而不是 $\ii/e^2$ 的归一化，正好让 simply-laced 情形下 $S$ 是干净的对合。

S 到底交换了什么？图~\ref{fig:sdual} 把这句话画成一张图：$\tau$ 与 $G$ 一起被换掉，而电与磁的角色互换。图上写了 $\tau'=-1/(n_g\tau)$，意思是右边的耦合是左边耦合的变换值；右边那个 $\Ldual$ 并不是同一个群，它是下一节的主角。

\begin{figure}[htbp]
\centering
\begin{tikzpicture}[>=Stealth, font=\small,
  side/.style={draw=kwblue,thick,rounded corners=3pt,fill=kwblue!5,
    align=center,inner sep=7pt,text width=4.1cm}]
  \node[side] (g) {$G$ 侧\\[3pt]
    $\tau=\theta/2\pi+4\pi\ii/e^{2}$\\[3pt]
    电激发：$G$ 的权格};
  \node[side,right=2.2cm of g] (dg) {$\Ldual$ 侧\\[3pt]
    $\tau'=-1/(n_g\tau)$\\[3pt]
    磁激发：$\Ldual$ 的权格};
  \draw[<->,very thick,kwred] (g.east) -- node[above=2pt,font=\small]{$S$}
    node[below=2pt,font=\scriptsize,align=center]{电 $\leftrightarrow$ 磁\\$G\leftrightarrow\Ldual$} (dg.west);
\end{tikzpicture}
\caption{S-对偶交换的三件事：把 $\tau$ 换成 $-1/(n_g\tau)$、把 $G$ 换成 $\Ldual$、把电荷与磁荷互换。它是量子对称，不在经典拉氏量里。}
\label{fig:sdual}
\end{figure}

还有一个细节第 3 讲会直接用到：$S$ 在超荷上不是平凡地作用。原文特别说明，这件事不能从经典拉氏量读出来，因为除了 $T$ 生成的子群以外，$\Gamma$ 根本不含经典对称 \src{B p.13}。量子推理给出的答案是：$S$ 在 $\mathbf 4$（即 $Q^{AX}$）上乘 $\ee^{\ii\hat\phi}$、在 $\bar{\mathbf 4}$ 上乘 $\ee^{-\ii\hat\phi}$，而
\[
\ee^{-\ii\hat\phi(\gamma)}=\Bigl(\frac{|c\tau+d|}{c\tau+d}\Bigr)^{1/2}
\]
（原文 (2.25)，\src{B p.15}）；平方根意味着作用在超荷上的是 $\Gamma$ 的双覆盖，要由 $(-1)^F$ 扩一下才对得上。为什么这件事要紧？因为第 3 讲扭化所选的，正是“哪些超荷在扭化后还活着”，所以 $S$ 怎样作用在超荷上，直接决定 S-对偶在扭化理论里变成什么。这一句可以先记下，不必现在算。

\subsection{GNO、Montonen--Olive，以及为什么是 $\Ldual$}

S-对偶不是凭空冒出来的，它有两个来源，而这两个来源正好构成本讲的骨架。

第一个是 GNO（Goddard--Nuyts--Olive，1976）。他们证明：在规范理论里，电荷取规范群的权格，而磁荷取“对偶群”的权格 \src{B p.2}；一般紧李群的磁荷分析更早由 Englert--Windey 给出 \src{B p.2}。为什么磁荷会落到别处？因为 't~Hooft--Polyakov 单极子不是微扰态：它的磁荷由解在无穷远处的绕数决定，而绕数的取值取决于 $G$ 里那个 $U(1)$ 子群怎么嵌入，也就是余权格。于是“能承载这种磁荷”的群，自然是 $G$ 的对偶群。原文在 §2 末的 Coulomb 分支讨论里把这件事写得更细：电、磁荷分别取权格与余权格，中心荷依赖 $\tau$ \src{B pp.14--15}。

第二个是 Montonen--Olive（1977）。他们比 GNO 更进一步，提出一个猜想：以 $G$ 为基础的那个规范理论，与把 $e$ 换成 $1/e$、把群换成对偶群的那个理论，是同一个理论 \src{B p.2}。GNO 说的是“磁荷住在哪个格”，MO 说的是“整个理论等价”；后者强得多，因为理论等价不仅要求谱一致，还要求所有关联函数一致。

两条线在 1977 年底合流。原文写：GNO 对偶群 actually the same as the Langlands dual group $\Ldual$，这是 Atiyah 观察到的，他当时建议第二作者注意 Langlands 纲领与量子场论的联系，并推荐了 GNO 与 MO 两篇文章 \src{B p.3}。之后是三步：超对称规范理论里 MO 对偶更自然 \src{B p.3}；$\mathcal N=4$ 是它最合适的场所；$\theta$ 角把原来的 $\ZZ_2$ 对偶扩成 $SL(2,\ZZ)$。

为什么 GNO 对偶群正是 $\Ldual$？本讲只给格层面的解释，因为原文本身也只给了这句话。Langlands 对偶在群论上就是把根与余根互换，于是权格与余权格互换；而磁荷量子化恰恰要求磁荷落在 $G$ 的余权格上，而“$G$ 的余权格”与“$\Ldual$ 的权格”是同一组格。就这一步，但值得写下来。单连通性的差别则表现为中心商：$SU(N)$ 与 $PSU(N)=SU(N)/\ZZ_N$ 的中心与基本群互换，这正是表~\ref{tab:dual} 第二行。再补一句诚实提醒：$\pi_1(G)\neq0$ 时还要看特征类 $\xi\in H^2(M,\pi_1(G))$，不同 $\xi$ 的配分函数按对偶群的酉表示变换 \src{B p.12}，这里不详谈。

\subsection{Table 1 怎么读，以及为什么这只是“证据”}

Table 1 只有两列：左边 $G$，右边 $\Ldual$。原文的七行照抄如下（只去掉排版）\src{B p.3}：

\begin{table}[htbp]
\centering
\small
\begin{tabular}{@{}ll@{}}
\toprule
$G$ & $\Ldual$ \\
\midrule
$U(N)$ & $U(N)$ \\
$SU(N)$ & $PSU(N)=SU(N)/\ZZ_N$ \\
$Spin(2n)$ & $SO(2n)/\ZZ_2$ \\
$Sp(n)$ & $SO(2n+1)$ \\
$Spin(2n+1)$ & $Sp(n)/\ZZ_2$ \\
$G_2$ & $G_2$ \\
$E_8$ & $E_8$ \\
\bottomrule
\end{tabular}
\caption{Table 1 的简化复制：$G$ 与其 Langlands（GNO）对偶群 $\Ldual$，见 \src{B p.3}。}
\label{tab:dual}
\end{table}

怎么读？三条读法。

第一，右列不是“随便另一个群”，而是由 $G$ 的根数据对偶地定出来的。$G$ 单连通时 $\Ldual$ 是伴随型，所以右列常出现商：$SU(N)$ 对偶到 $PSU(N)$，$Sp(n)$ 对偶到 $SO(2n+1)$，$Spin(2n)$ 对偶到 $SO(2n)$ 再除中心。对偶是对合的：把表倒过来读，$PSU(N)$ 的对偶正是 $SU(N)$。

第二，中心与基本群互换。$SU(N)$ 的中心是 $\ZZ_N$，$PSU(N)$ 的基本群是 $\ZZ_N$；这一行是“磁荷取余权格”最干净的体现。同理，$Sp(n)$ 那一行的奥妙在于：$Sp(n)$ 与 $Spin(2n+1)$ 是仅有的“李代数不同构”的一对对偶群，其余各行两边都是同一个李代数；正因为李代数不同构，$\tau$ 在两边的归一化要重新对齐，$n_g=2$ 就是这么来的。

第三，自对偶行要小心。$G_2$ 与 $E_8$ 的根与余根同构，所以自对偶；$U(N)$ 行也自对偶，但严格说 $U(N)$ 不是半单群，它的 $U(1)$ 因子要单独处理——§1 写这张表时还没有 §2 的“$G$ 紧”假设，读的时候要知道这个差别。还有一处缺席值得注意：表里没有 $F_4$。它作为群是自对偶的，但它非 simply-laced，$S$ 作用在 $\tau$ 上是 $\tau\mapsto-1/(2\tau)$。这个小小的缺席恰好说明：群的对偶与耦合的对偶是两个层次，前者是群论，后者还要看归一化。

最后把这一讲的地位说清楚：以上全是期望与证据，不是定理。原文自己用的词就是 conjecture 与 asserts——$S$ 是被“断言”为一个量子对称的 \src{B p.11}。四十多年来它有一大批互相独立的证据：BPS 谱的匹配、瞬子配分函数的计算、弦论里的实现，以及整本 \cite{kw06} 的前提——如果 S-对偶不成立，几何 Langlands 的物理推导根本起不来。但“证据极多”不等于“已经证明”。到 2026 年为止，四维 $\mathcal N=4$ SYM 的 S-对偶仍然没有数学证明。这也解释了 K--W 的写法：它把 S-对偶当作公理式的输入，推出一整套几何 Langlands 的陈述；而真正给出独立证明的，是 2024 年前后 Gaitsgory 等人的纯代数几何工作，那条路线完全不经过四维场论。物理路线与数学证明是两条独立的路，最后指向同一个答案——这是全篇最让人服气的地方，也是最需要你自己判断的地方。

\begin{keybox}
$\mathcal N=4$ SYM 是 S-对偶最干净的场所：$S$ 把 $\tau$ 换成 $-1/(n_g\tau)$、把 $G$ 换成 $\Ldual$、把电与磁互换。$\Ldual$ 正是 GNO 1976 从磁荷量子化读出的对偶群，也是 Langlands 对偶群；但这一整套是猜想与证据，不是定理。
\end{keybox}

\begin{follow}
\begin{itemize}
\item \textbf{B §1，pp.2--5}：GNO、Montonen--Olive、Atiyah 的观察、Table 1 与后续发展。全篇所有“为什么”的来源，先读它再读技术细节。
\item \textbf{B §2.1，pp.6--11}：场内容、十六超荷、降维、$\tau$ 与 $SU(4)_R$。沿 (2.2)--(2.12) 读，盯住反厄米约定与 $\Tr$ 的归一化。
\item \textbf{B §2.2，pp.11--16}：S-对偶的断言、(2.16)/(2.17)，以及 pp.13--15 超荷上的相位 $\hat\phi$。这一段是第 3 讲扭化的直接前置。
\end{itemize}
\end{follow}

\begin{checkbox}
\begin{itemize}
\item 写出 $\mathcal N=4$ SYM 的场内容：规范场、实标量、Weyl 费米子各几个？它们分别从十维的哪个场来？
\item $\theta$ 为什么是周期变量？$\tau$ 为什么落在上半平面？为什么 $T:\tau\mapsto\tau+1$ 是经典对称而 $S$ 不是？
\item S-对偶断言的三件事是什么？非 simply-laced 时 $\tau\mapsto-1/(n_g\tau)$，$n_g$ 是什么，为什么 $F_4$ 对应 $n_g=2$？
\item GNO 1976 说了什么？为什么单极子的磁荷落在余权格而不是权格？
\item 读 Table 1：$SU(N)$ 一行的中心与基本群怎样互换？为什么表里没有 $F_4$？
\item 这一讲里哪些是定理、哪些是猜想？S-对偶对 $\mathcal N=4$ 被证明了吗？若没有，K--W 凭什么用它？
\end{itemize}
\end{checkbox}

% ======================================================================
%  第 3 讲  扭化：为什么要把超对称「扭」掉
%  页锚以 B 版原文（\cite{kw06}）的印刷页为准
% ======================================================================
\section{扭化：为什么要把超对称「扭」掉}
\label{sec:twist}

\begin{goals}
\begin{enumerate}
\item 说清楚弯曲流形上原来的超荷为什么不够用，拓扑场论到底向超荷要什么。
\item 说清楚 GL 扭化选了哪个子群，为什么恰好是它把六个标量排成一个一形式加两个标量。
\item 说清楚参数 $t\in\PP^1$ 与典范参数 $\Psi$ 的关系，以及为什么只有 $\Psi$ 有物理意义。
\end{enumerate}
\end{goals}

\subsection{弯曲流形上的麻烦}

先回忆平坦 $\RR^4$ 上的情形。$\mathcal N=4$ 超对称 Yang--Mills 理论有 16 个超荷；在超对称场论里，能动张量可以写成
\begin{equation}
T_{\mu\nu}=\{Q,\,G_{\mu\nu}\}
\label{eq:T-exact}
\end{equation}
的形式（$G_{\mu\nu}$ 是某个张量算符，等号在适当的改善项意义下成立）。只要真空保持 $Q$，$T_{\mu\nu}$ 就是 $Q$-恰当的，基态能量为零，配分函数对度量的整体缩放不敏感。这是拓扑场论最朴素的来源：路径积分会「局部化」到 $Q$-不变构型上，而 $Q$-恰当的能动张量保证度量形变不改变关联函数。

麻烦在于，这个论证默认 $Q$ 是平坦空间里的整体守恒荷。把理论放到一般弯曲流形 $M$ 上，超荷不再自动是整体定义的标量：它们按 $M$ 的旋量丛变换，沿 $M$ 平行移动时会互相混合。一般度规的和乐群会把 16 个超荷搅在一起，通常一个也留不下来，更谈不上从中挑出一个满足 \eqref{eq:T-exact} 的 $Q$。换句话说，普通超对称在弯曲流形上\emph{自动}掉了，理论并不因此变成拓扑的。\src{B p.16}

要让理论「拓扑」，就必须换一个在任意 $M$ 上都存在的超荷。做法不是改物理，而是改我们称呼 Lorentz 群的方式：把 $M$ 切丛的结构群 $\mathrm{Spin}(4)$ 对角地嵌入 $\mathrm{Spin}(4)\times\mathrm{Spin}(6)$，其中 $\mathrm{Spin}(6)$ 是 $\mathcal N=4$ 理论的 $R$-对称群。这个操作叫\term{twisting}{扭化}。扭化之后，场的自旋按新的对角群来数，原来按 $\mathrm{Spin}(6)$ 指标变换的场被重新组合成 $M$ 上的普通张量场；更重要的是，16 个超荷里会出现若干在 $M$ 上整体定义的标量分量——它们才是候选的拓扑超荷。在平坦空间上，扭化前后的理论只是同一批场的重新打包，物理内容并没有变。

拓扑场论的判据可以写得很干净：存在一个标量超荷 $Q$，满足
\begin{equation}
Q^2=0,\qquad T_{\mu\nu}=\{Q,\,b_{\mu\nu}\},
\label{eq:top-criterion}
\end{equation}
于是一切度量变分都是 $Q$-恰当的，关联函数与度量无关。这个判据在 \cite{fr09} 里有明确表述（\src{D p.13}）。需要说明的是，\cite{kw06} 并没有从这个判据出发去证明什么；原文的做法是\emph{先构造}：挑一个扭化，验证 $Q^2=0$，再检查得到的理论确实有拓扑性。判据是事后的组织原则，不是原文的证明路线。这一点值得先记住，因为本讲后面几处「为什么恰好是这个扭化」都只有构造性的回答。

在 $\mathcal N=4$ 上，$R$-对称群是 $\mathrm{Spin}(6)\cong SU(4)$，比 $\mathcal N=2$ 的 $U(1)_R$ 大得多，因此可选的扭化不止一种。原文 p.17 说有\emph{三种}不等价的扭化：前两种（其中一种在文献 [45] 里讨论过）与 Donaldson 理论密切相关，给出瞬子型的拓扑不变量；第三种在 [52--57] 里被研究过，却一直没有找到用武之地，而它恰好是与几何 Langlands 有关的那一个，原文称之为 \term{GL twist}{GL 扭化}。Donaldson--Witten 的类比是这样来的：$\mathcal N=2$ 超对称 Yang--Mills 经过一次 $U(1)_R$ 扭化成为拓扑 Yang--Mills 理论，其路径积分局部化在瞬子上，关联函数给出四维流形的 Donaldson 不变量。第三种扭化可以看作这个故事的 $\mathcal N=4$ 版本，只不过它多留下了一整个参数族。

\subsection{GL 扭化：选哪个子群}

把 $\mathcal N=4$ 的场按 $\mathrm{Spin}(4)\times\mathrm{Spin}(6)$ 写出来：规范场 $A_\mu$ 是 $(\mathbf 2,\mathbf 2;\mathbf 1)$，六个标量 $\varphi_I$ 是 $(\mathbf 1,\mathbf 1;\mathbf 6)$，费米子落在 $(\mathbf 2,\mathbf 1;\mathbf 4)\oplus(\mathbf 1,\mathbf 2;\bar{\mathbf 4})$ 里。GL 扭化选取一个同态
\begin{equation}
\kappa:\ \mathrm{Spin}(4)\longrightarrow\mathrm{Spin}(6),
\label{eq:kappa}
\end{equation}
它把 $\mathrm{Spin}(4)$ 的旋量表示 $\mathbf 4$ 送到 $\mathrm{Spin}(6)$ 的旋量表示 $\mathbf 4$，然后令
\begin{equation}
\mathrm{Spin}'(4)=(1\times\kappa)\bigl(\mathrm{Spin}(4)\bigr)\subset \mathrm{Spin}(4)\times\mathrm{Spin}(6)
\label{eq:spinx}
\end{equation}
为新的 Lorentz 群。$\kappa$ 在 $\RR^4$ 上与原来的 $\mathrm{Spin}(4)$ 作用相同，但在规范场上作用不同——这正是「扭」这个字的全部含义。\src{B p.17}

扭化能不能给出拓扑超荷，判据是正手征旋量空间 $S_+$ 里有没有非零的不变向量。对 GL 扭化，答案是肯定的，而且不止一个：$S_+$ 里有两个线性无关的不变旋量，记作 $\varepsilon_\ell,\varepsilon_r$，分别来自 $(\mathbf 2,\mathbf 1;\mathbf 4)$ 与 $(\mathbf 1,\mathbf 2;\bar{\mathbf 4})$。\src{B p.18} 相应地，扭化后的理论保留两个标量超荷 $Q_\ell,Q_r$，任意组合
\begin{equation}
Q=u\,Q_\ell+v\,Q_r,\qquad (u,v)\neq(0,0)
\label{eq:Q-combo}
\end{equation}
都是标量、都在任意 $M$ 上有定义。为什么 $Q^2=0$？原文的论证是：$Q^2$ 作为算符必须与新的 Lorentz 群交换，而 \src{B p.18 eq (2.15)} 的右边不存在 $\mathrm{Spin}'(4)$-不变的算符，唯一可能是零；写成反对易子就是 $Q^2=\tfrac12\{Q,Q\}=0$。\src{B p.18} 这个论证只针对平移部分。把规范部分也算进来，有
\begin{equation}
\delta_T^2=-\ii\,(u^2+v^2)\,L_\sigma,
\label{eq:deltaT2}
\end{equation}
其中 $L_\sigma$ 是场依赖的规范变换（$\sigma$ 是后面要出现的标量）。因此严格平方为零只在 $u^2+v^2=0$ 时成立；对一般的 $t$，只能说 $Q^2$ 是「模掉一个规范变换」为零。\src{B p.31}\src{B p.33} 后文会看到，$u^2+v^2=0$ 恰是 $t=\pm\ii$ 那两个点。

为什么要选这个 $\kappa$，而不是前两种扭化？原文没有给出「非它不可」的证明，这是需要诚实说明的地方：它只构造并指出这一种与几何 Langlands 相关。不过构造本身已经透露了理由。按 $\mathrm{Spin}'(4)$ 分解，六个标量所在的空间 $\wedge^2\mathbf 4$ 拆成
\begin{equation}
\mathbf 6_v=\wedge^2\mathbf 4
=(\mathbf 2,\mathbf 2)_0\oplus(\mathbf 1,\mathbf 1)_{2}\oplus(\mathbf 1,\mathbf 1)_{-2},
\label{eq:6-decomp}
\end{equation}
下标的 $K$ 是 $\mathrm{Spin}(6)$ 里由 $\Gamma_{89}$ 生成的 $U(1)$ 荷。第一项 $(\mathbf 2,\mathbf 2)_0$ 正是一个一形式，后两项是两个标量。也就是说，扭化把「六个标量」自动重组为「一个伴随值一形式加两个复标量」；而一形式 $\varphi=\varphi_\mu\dd x^\mu$ 与规范场 $A_\mu$ 可以合并成复联络 $\mathcal A=A+\ii\varphi$。当 $t=\pm\ii$ 时，后面的 BPS 方程恰好是 $\mathcal A$ 平坦、外加 $D{*}\varphi=0$，这正是 Hitchin 方程；而平坦 $G_{\mathbb C}$ 联络与 Hitchin 模空间正是几何 Langlands 的舞台。前两种扭化给不出这样的「一形式加平坦复联络」结构，这就是 GL 扭化与几何 Langlands 相连的结构性原因。

\subsection{扭化后场长什么样}

把扭化后的场列成表，会看得更清楚。玻色子的内容如表~\ref{tab:bosons} 所示：
\begin{table}[htbp]
\centering
\caption{扭化后玻色场的 $\mathrm{Spin}'(4)$ 表示与 $K=\Gamma_{89}$ 荷。前两行同属 $(\mathbf 2,\mathbf 2)_0$，扭化后在 Lorentz 群下不可区分；区别在来源。}
\label{tab:bosons}
\begin{tabular}{@{}llcl@{}}
\toprule
场 & $\mathrm{Spin}'(4)$ 表示 & $K=\Gamma_{89}$ & 来源 \\
\midrule
$A_\mu$ & $(\mathbf 2,\mathbf 2)_0$ & $0$ & 四维规范场 \\
$\varphi_\mu$ & $(\mathbf 2,\mathbf 2)_0$ & $0$ & $\mathbf 6$ 中的 $\mathbf 4$ 个标量 \\
$\sigma$ & $(\mathbf 1,\mathbf 1)_2$ & $2$ & $\mathbf 6$ 中的单态 \\
$\bar\sigma$ & $(\mathbf 1,\mathbf 1)_{-2}$ & $-2$ & $\mathbf 6$ 中的单态 \\
\bottomrule
\end{tabular}
\end{table}
\src{B p.19} 两个一形式都是 $(\mathbf 2,\mathbf 2)_0$，在扭化后的 Lorentz 群下无法区分，它们的区别在来源：$A_\mu$ 是规范场，$\varphi_\mu$ 来自标量。$\sigma$ 与 $\bar\sigma$ 是复标量；它们的 $K$ 荷为 $\pm2$，这解释了为什么 $\sigma$ 在超荷下不变：超对称变分要把 $\sigma$ 变成 $K=3$ 的费米子，而场里根本没有 $K=3$ 的分量，于是 $\delta_T\sigma=0$。\src{B p.24}

费米子的 16 个分量按 $K$ 分块，只有 $K=-\ii$ 的两个分量与不变超荷对应，其余分量两两配对，在扭化后的理论里不产生额外的标量超荷。\src{B p.23} 一个容易忽略的细节是：两个不变超荷的 $K$ 荷\emph{相同}。这一点直接决定了参数族的形状。若两个超荷的 $K$ 荷不同，$U(1)_K$ 对称性会把其中一个（或某个线性组合）挑出来，$(u,v)$ 的整个 $\PP^1$ 就被压成一个点；正因为 $K$ 荷相同，两个超荷才平权，族才是完整的 $\PP^1$。\src{B p.22}

扭化后的超对称变换把 $A_\mu$ 与 $\varphi_\mu$ 混合，把 $\sigma$ 单独留下。写出来最能说明问题的是 $t=\pm\ii$ 时的 BPS 方程：
\begin{equation}
\bigl(F-\varphi\wedge\varphi+t\,D\varphi\bigr)^+=0,\qquad
\bigl(F-\varphi\wedge\varphi-t^{-1}D\varphi\bigr)^-=0,\qquad
D{*}\varphi=0,
\label{eq:bps}
\end{equation}
其中 $\pm$ 是四维流形上的自对偶/反自对偶投影。当 $t=\pm\ii$ 时，前两条方程等价于
\begin{equation}
F(\mathcal A)=0,\qquad \mathcal A=A+\ii\varphi,
\label{eq:flat}
\end{equation}
即\emph{平坦 $G_{\mathbb C}$ 联络}，再加上 $D{*}\varphi=0$。其解模规范变换，就得到平坦联络模空间的一部分（\src{B p.25}）：
\begin{equation}
Y_M(G)=\operatorname{Hom}\bigl(\pi_1(M),G_{\mathbb C}\bigr).
\label{eq:flatmoduli}
\end{equation}
当 $t$ 为实数时，方程变成椭圆型方程族；$t=0,\infty$ 是退化极限，非平凡瞬子只在极限附近出现。\src{B p.26}

\subsection{一族理论：$t\in\PP^1$ 与典范参数 $\Psi$}

由 \eqref{eq:Q-combo}，只差一个整体标度，两个数 $(u,v)$ 给出一个复参数
\begin{equation}
t=\frac{v}{u}\in\PP^1=\CC\cup\{\infty\}.
\label{eq:t-param}
\end{equation}
所以扭化不是给一个理论，而是给一族理论。族上还有两个「显然」的等价：整体缩放 $(u,v)\mapsto\lambda(u,v)$ 不改变 $t$；$SU(4)_R$ 的中心元素 $I$ 作用为 $I(t)=-t$，把 $t$ 与 $-t$ 等同（\src{B p.22}）。除这两条之外，四维一侧没有别的等价——真正的等价来自 S-对偶，因子 $(c\tau+d)/|c\tau+d|$ 作用在 $t$ 上，使 $t$ 不再只是标签。这正是后面能生出几何 Langlands 的地方。

物理上真正有意义的参数不是 $t$，而是 $t$ 与复耦合 $\tau=\theta/2\pi+4\pi\ii/g_{\mathrm{YM}}^2$ 的组合。原因是：扭化后作用量可以写成
\begin{equation}
I=\delta_t V+\frac{\ii\Psi}{4\pi}\int_M\Tr(F\wedge F),
\label{eq:action-Psi}
\end{equation}
第一项是超对称变分，第二项是拓扑项，其系数只可能依赖一个组合。配平量纲与 $\tau$ 依赖后得到
\begin{equation}
\Psi=\frac{\tau+\bar\tau}{2}+\frac{\tau-\bar\tau}{2}\cdot\frac{t-t^{-1}}{t+t^{-1}}
=\operatorname{Re}\tau+\ii\operatorname{Im}\tau\,\frac{t-t^{-1}}{t+t^{-1}}.
\label{eq:Psi}
\end{equation}
\src{B pp.35--36} 把 \eqref{eq:Psi} 右边通分，可以写成等价的有理形式
\begin{equation}
\Psi=\frac{\tau t^2+\bar\tau}{1+t^2}.
\label{eq:Psi-rational}
\end{equation}
这个形式在讨论特殊点时最方便。它的逆（\src{B p.36 eq (3.52)}）是
\begin{equation}
t^2=\frac{\bar\tau-\Psi}{\Psi-\tau},
\label{eq:t2-inverse}
\end{equation}
注意这里没有写错：分子分母都含 $\tau$ 的共轭差，而不是把 $\Psi-\tau$ 倒过来。我们核对过原文页图像；某些转录版本把共轭符号丢了，变成错误的倒数式，读者若按错误式子验算会得到相反的 $t$ 依赖。

由 \eqref{eq:Psi-rational} 立刻读出几个特殊点。$t=\pm\ii$ 时分母 $1+t^2=0$，而分子 $\tau t^2+\bar\tau=\bar\tau-\tau\neq0$（只要 $\operatorname{Im}\tau\neq0$），所以
\begin{equation}
t=\pm\ii\quad\Longrightarrow\quad \Psi=\infty,
\label{eq:Psi-inf}
\end{equation}
并且这个 $\infty$ 与 $\tau$ 的具体值无关。注意论证用的是有理形式：若用逆式 \eqref{eq:t2-inverse}，在 $t=\pm\ii$ 处会遇上 $0/0$，什么也读不出来。$t=\pm1$ 时 $(t-t^{-1})/(t+t^{-1})=0$，故 $\Psi=\operatorname{Re}\tau$；$|t|=1$ 时 $t=\ee^{\ii\theta}$，$(t-t^{-1})/(t+t^{-1})=\ii\tan\theta$，于是 $\Psi=\operatorname{Re}\tau-\operatorname{Im}\tau\tan\theta$ 为实数；弱耦合 $\operatorname{Im}\tau\to\infty$ 而固定 $\Psi,\operatorname{Re}\tau$ 时，由逆式得 $t^2\to1$，即 $t\to\pm1$。\src{B pp.36--37}

为什么说理论只依赖 $\Psi$，即为什么 $t$ 与 $\tau$ 的其它组合都不进入物理？原文给了两步论证（\src{B pp.36--37}）。第一步，作用量 \eqref{eq:action-Psi} 里只有拓扑项的系数含 $\tau$ 与 $t$，而所有可观测的 $\tau,t$ 依赖都必须通过这个系数，也就是通过 $\Psi$。第二步，对 $t\neq\pm\ii$ 定义
\begin{equation}
\delta'_t=\frac{\delta_t}{\sqrt{1+t^2}},
\label{eq:delta-prime}
\end{equation}
则 $(\delta'_t)^2=-\ii L_\sigma$ 与 $t$ 无关；再用小的场重定义把作用量整理成 Cartan 模型的标准形式，$t$ 的剩余依赖也被吸收掉。两步合起来：$(t,\tau)$ 只有通过 $\Psi$ 才有物理意义。顺带解释了 $t\sim-t$：$\Psi$ 只含 $t^2$，所以 $t$ 与 $-t$ 给出同一个 $\Psi$。

在 $\Psi$ 的语言里，S-对偶与 T-变换的作用极其简单：
\begin{equation}
S:\ \Psi\mapsto-\frac{1}{n_g\Psi},\qquad
T:\ \Psi\mapsto \Psi+1,
\label{eq:S-T}
\end{equation}
一般群元素 $\gamma=\begin{pmatrix}a&b\\c&d\end{pmatrix}\in SL(2,\ZZ)$ 给出 $\Psi\mapsto(a\Psi+b)/(c\Psi+d)$。\src{B p.37} 这里 $n_g$ 是\term{lacing number}{花边数}：单边（ADE）时 $n_g=1$，$B,C,F_4$ 时 $n_g=2$，$G_2$ 时 $n_g=3$。它来自 $\Ldual$ 的根系与原根系的对偶，$S$ 的作用因此带一个分数系数；$n_g=1$ 时退化为熟悉的 $\Psi\mapsto-1/\Psi$。$S$ 把 $\Psi=\infty$ 送到 $\Psi=0$，这一条第 4 讲要用。

参数空间的形状见图 \ref{fig:cp1}。
\begin{figure}[htbp]
\centering
\begin{tikzpicture}[scale=1.0]
  \draw[thick] (0,0) circle (1.8);
  \node[below] at (0,-1.98) {$t=0$};
  \node[above] at (0,1.98) {$t=\infty$};
  \node[right] at (1.88,0) {$t=1$};
  \node[left] at (-1.88,0) {$t=-1$};
  \fill (0,0.55) circle (2.4pt);
  \node[right] at (0.12,0.62) {$t=\ii$};
  \draw (0,-0.55) circle (2.4pt);
  \draw[thick] (-0.07,-0.62)--(0.07,-0.48);
  \draw[thick] (-0.07,-0.48)--(0.07,-0.62);
  \node[right] at (0.12,-0.68) {$t=-\ii$};
  \node[align=left,anchor=west,text width=5.7cm] at (2.3,1.15)
    {$\ii,-\ii$：$\Psi=\infty$，B-模型\\ 与 $\tau$ 无关};
  \node[align=left,anchor=west,text width=5.7cm] at (2.3,-1.25)
    {实轴 $\pm1$：$\Psi=\operatorname{Re}\tau$；\\ 其余实 $t$：$\Psi$ 为实数，A-模型};
\end{tikzpicture}
\caption{参数空间 $\PP^1_t$ 的示意图：圆是 $\PP^1$ 的轮廓，实轴画成轮廓上的四个标记点，$\pm\ii$ 画在圆内（分别用实心点与叉点表示）。$t$ 与 $-t$ 给出同一个 $\Psi$，$t=\pm\ii$ 对应 $\Psi=\infty$。}
\label{fig:cp1}
\end{figure}

最后是一个必须说在前面的警告。这套构造只在 \emph{Euclidean 号差}下成立。在 Lorentz 号差下，扭化要用一个紧致群 $\mathrm{Spin}(6)$ 去「吃掉」非紧的 $\mathrm{Spin}(1,3)$，找不到这样的单射；若改用 $\mathrm{Spin}(1,5)$，得到的哈密顿量不是厄米的，能量无下界。原文在 \src{B p.56 脚注 19} 把这一点讲得很直白：扭化破坏幺正性，只在 Euclidean 号差下有意义。这不是技术细节，而是整个几何 Langlands 故事必须在复几何（Riemann 面、复结构）里讲的根本原因；读原文时凡遇到「流形」「度量」都要默认其欧几里得。另外，\src{B p.16 脚注 5} 指出：当 $\theta=0$ 且 $\Psi=0,\infty$ 时，连不可定向的流形也可以允许。还要提醒：本文的推导是物理推导而非数学证明，且只处理非分歧情形；2024 年前后 Gaitsgory 等人的代数几何工作才给出几何 Langlands 的严格证明。

\begin{keybox}
扭化不是改变物理，而是换一个在弯曲流形上也存在的标量超荷：GL 扭化取 $\mathrm{Spin}(4)\subset\mathrm{Spin}(4)\times\mathrm{Spin}(6)$ 的对角子群，使六个标量重排为一个一形式与两个标量，并留下一族以 $t\in\PP^1$（等价地以 $\Psi$）参数化的拓扑场论。
\end{keybox}

\begin{follow}
\begin{itemize}
\item \textbf{B §3.1, pp.17--24}：三种扭化、$\kappa$ 的选取、$Q^2=0$ 的论证、场的重排与 $\delta_T\sigma=0$（重点读 pp.17--19 与 p.24）。
\item \textbf{B §3.2, pp.25--28}：扭化后的方程 \eqref{eq:bps}、$t=\pm\ii$ 与平坦 $G_{\mathbb C}$ 联络、实 $t$ 的椭圆族。
\item \textbf{B §3.5, pp.35--38}：$\Psi$ 的定义式、t--$\tau$ 合并、$t=\pm\ii$ 的 $\Psi=\infty$、$S$ 与 $T$ 的作用。
\end{itemize}
\end{follow}

\begin{checkbox}
\begin{itemize}
\item 为什么平坦空间上的 16 个超荷不能直接搬到一般弯曲流形上？扭化究竟改变了什么、又没有改变什么？
\item 写出 $\mathbf 4$ 与 $\mathbf 6_v$ 在 $\mathrm{Spin}'(4)$ 下的分解，并解释 $\sigma$ 为何满足 $\delta_T\sigma=0$。
\item 用 \eqref{eq:Psi-rational} 验证 $t=\pm\ii$ 时 $\Psi=\infty$ 且与 $\tau$ 无关；说明为什么不能改用逆式 \eqref{eq:t2-inverse} 来验证。
\item 若两个不变超荷的 $K$ 荷不同，$(u,v)$ 的参数空间会变成什么？为什么？
\item 说明为什么可观测物理只依赖 $\Psi$，以及 $t$ 与 $-t$ 为何给出同一个理论。
\item 为什么这套扭化只在 Euclidean 号差下成立？Lorentz 号差下具体卡在哪一步？
\end{itemize}
\end{checkbox}

% ======================================================================
%  第 4 讲  紧化：四维的桥塌成二维的镜像
%  页锚以 B 版原文（\cite{kw06}）的印刷页为准
% ======================================================================
\section{紧化：四维的桥塌成二维的镜像}
\label{sec:compact}

\begin{goals}
\begin{enumerate}
\item 说清楚把理论放在 $\Sigma\times C$ 上、让 $C$ 变小时，低能理论为什么是以 $\MH$ 为靶空间的二维 sigma 模型。
\item 说清楚这个二维理论在参数特殊值处是 A-模型还是 B-模型，以及 $\Psi$ 如何统一描述它们。
\item 说清楚四维的 $S$-对偶为什么变成二维的镜像对称，并把核心交换图的两条路径对上。
\end{enumerate}
\end{goals}

\subsection{把世界做成乘积 $\Sigma\times C$}

本讲对应 \cite{kw06} 的 §4--§5。几何 Langlands 的语言是二维的：曲线 $C$ 上的层、Hecke 算子、平坦联络、$\Dmod$-模。四维的 Kapustin--Witten 理论要「长出」一个二维世界，第一步就是紧化。原文的选择是把四维理论放在
\begin{equation}
M=\Sigma\times C
\label{eq:M-prod}
\end{equation}
上，其中 $C$ 是做几何 Langlands 的那条紧 Riemann 面（亏格 $g>1$，保证一般平坦联络不可约），$\Sigma$ 是另一个二维流形，充当低能二维理论的「时空」。\src{B p.38} 尺寸的层级是全部推导的前提：$C$ 取得很小，$\Sigma$ 取得很大（$\Sigma$ 甚至可以是不紧的，例如 $\RR^2$）。因为扭化后的理论是拓扑场论，对度量没有依赖，所以把 $M$ 取成乘积、把两个因子的尺度拉开，都不损失一般性。

为什么这个层级正好给出二维理论？原因在 Kaluza--Klein 谱。把四维场按 $C$ 上的本征模式展开，第 $\mathbf n$ 个模式的质量平方约为 $\lambda_{\mathbf n}/R_C^2$；当能量远小于 $1/R_C$ 时，所有非零模式都被积掉，只剩零模。另一方面，零模沿 $\Sigma$ 可以缓变：场的 $\Sigma$ 方向导数每多一阶，作用量就多一个 $R_C^2$ 的因子，$\Sigma$ 很大时这一项可以任意小。于是「$C$ 小」保证只有零模参与低能物理，「$\Sigma$ 大」保证零模在 $\Sigma$ 上形成真正的二维场。两个条件缺一不可：若 $\Sigma$ 也小，就只剩量子力学；若 $C$ 不小，非零模式掺进来，二维描述失效。\src{B p.39}

还有一个超对称的账要算：四维的 16 个超荷在紧化后保留 8 个，它们分成左右各 4 个，即在 $\Sigma$ 上形成二维 $(4,4)$ 超对称。\src{B p.46} 这 8 个超荷正是后面二维 A-模型与 B-模型的超对称来源；若保留得更少，二维理论就不会是拓扑的。

\subsection{靶空间为什么是 $\MH$}

下一步要认出二维理论的靶空间。这要分两步，原文把它们压在一句话里，我们拆开看。

\textbf{第一步：认出超对称位形。}由第 3 讲，$t=\pm\ii$ 时四维的超对称位形是平坦 $G_{\mathbb C}$ 联络；更一般地，$t\neq0,\infty$ 时任何超对称位形都与 $t$ 无关，并且自动是平坦 $G_{\mathbb C}$ 联络（原文的定理 2）。\src{B pp.28--30} 所以在 $M=\Sigma\times C$ 上要研究平坦 $G_{\mathbb C}$ 丛 $E$。平坦丛在乘积流形上可以按两个因子分解：
\begin{equation}
E=\bigoplus_i E'_i\boxtimes E''_i,
\label{eq:E-decomp}
\end{equation}
其中 $E'_i$ 是 $\Sigma$ 上的平坦丛，$E''_i$ 是 $C$ 上的平坦丛。固定一点 $x\in\Sigma$，把 $E$ 限制到 $C_x=\{x\}\times C$ 上，得到
\begin{equation}
E_x=\bigoplus_i E'_i|_x\otimes E''_i,
\label{eq:E-x}
\end{equation}
它是 $C$ 上的一个平坦 $G_{\mathbb C}$ 丛；其同构类正是 $\MH(G,C)$ 的一个点（Hitchin 模空间与平坦丛模空间的对应，见第 1 讲）。让 $x$ 在 $\Sigma$ 上变动，就得到映射
\begin{equation}
\Phi:\ \Sigma\longrightarrow \MH(G,C),\qquad x\longmapsto [E_x].
\label{eq:Phi}
\end{equation}
这就是二维理论的场。原文在这里有一个技术约定：若某个 $E_x$ 是可约的平坦丛，它对应 $\MH$ 的奇异点；避开奇异点，即要求 $E$ 的限制不可约。\src{B p.30}

\textbf{第二步：认出能量。}光有映射还不够，还要说明低能作用量就是 sigma 模型的作用量。第 3 讲 §3.4 把四维玻色作用量写成了平方和加拓扑项（\src{B p.34 eq (3.47)--(3.48)}）：
\begin{equation}
S=\int_M\bigl(\text{平方和}\bigr)+\frac{\ii\Psi}{4\pi}\int_M\Tr(F\wedge F),
\label{eq:action-sum-sq}
\end{equation}
平方和为零正好是 Hitchin 方程
\begin{equation}
F-\varphi\wedge\varphi=0,\qquad D\varphi=0,\qquad D{*}\varphi=0,
\label{eq:hitchin}
\end{equation}
即 $A+\ii\varphi$ 平坦且 $D{*}\varphi=0$。在 $\Sigma$ 方向缓变的位形上，平方和项展开后正是映射 $\Phi$ 的动能（其二次项给出 $\MH$ 上的度量），拓扑项则给出二维的 $B$-场与 theta 角。于是低能有效理论就是以 $\MH(G,C)$ 为靶空间的二维 sigma 模型。原文 p.38 说这一步「由 §3.4 的结果得到」；把「认出靶空间」与「认出能量」分开，读者才不会误以为靶空间是从作用量直接读出来的。\src{B p.34}\src{B p.38}

顺带记下靶空间的几何：$\MH(G,C)$ 是超凯勒流形，复维数 $2(g-1)\dim G$，带一整族复结构（twistor 球面 $\PP^1_h$）。这个超凯勒结构是下一小节 A-模型/B-模型分类的全部来源。

\subsection{二维理论：A-模型、B-模型与 $\Psi$}

二维理论有 $(4,4)$ 超对称，靶空间 $\MH$ 又是超凯勒的，因此超对称变分可以按左右两个「扭化」独立选择复结构。设左右各选一个复结构 $J_+,J_-$（都来自 $\PP^1_h$），BPS 方程是
\begin{equation}
(1-\ii J_+)\,\bar\partial\Phi=0,\qquad
(1-\ii J_-)\,\partial\Phi=0.
\label{eq:bps2d}
\end{equation}
\src{B pp.59--61} 三种情形由此分出：

\begin{itemize}
\item $J_+=J_-=J$：两个方程合起来只依赖复结构 $J$，这是 \term{B-model}{B-模型}，其关联函数只依赖靶空间的复结构。
\item $J_+=-J_-=K$：方程只依赖辛形式，这是 \term{A-model}{A-模型}，其关联函数只依赖靶空间的辛结构（复化 Kähler 类）。
\item $J_-\neq\pm J_+$：既不是 A 也不是 B；但广义复几何告诉我们，它等价于带 $B$-场的 A-模型。原文只给结论、不给证明。\src{B p.59}\src{B pp.66--68}
\end{itemize}

四维的 $t$ 如何映到这一对复结构？原文给出
\begin{equation}
(w_+,w_-)=(-t,\ t^{-1}),
\label{eq:w-t}
\end{equation}
即左右复结构分别由 twistor 球面上的点 $w_+,w_-$ 指定；真正进入二维理论的是比值
\begin{equation}
q=\frac{w_+}{w_-}=-t^2.
\label{eq:q}
\end{equation}
\src{B p.61}\src{B p.62} 于是：$t=\pm\ii$ 时 $w_+=w_-$，二维理论是 \emph{B-模型}（$t=\ii$ 给复结构 $J$，$t=-\ii$ 给 $-J$）；$t$ 为实数时 $w_-=-1/\bar w_+$，二维理论是 \emph{A-模型}（$t=\pm1$ 给辛结构 $K$；$t=0,\infty$ 是退化极限）。这正是第 3 讲所说「$t=\pm\ii$ 像 B-模型、实 $t$ 像 A-模型」的精确版本。\src{B p.61}

但二维理论真正依赖的参数是 $\Psi$，不是 $t$。二维 A-模型的参数是复化 Kähler 类 $[B+\ii\omega]$；把四维的拓扑项降维后，可以算出
\begin{equation}
[B+\ii\omega]=-[\omega_I]\,\Psi,
\label{eq:B-i-omega}
\end{equation}
即在 $t\neq\pm\ii$ 时，$\PP^1$ 上每一点都等价于同一个 A-模型，只是 $B$-场平移了 $\Psi$ 倍。\src{B pp.68--69} 这就从二维一侧解释了「只有 $\Psi$ 有物理意义」。$t=\pm\ii$ 时上述推导因极点失效，二维理论是 B-模型，而 $\Psi=\infty$；此时理论对 $\tau$ 的无关性要单独论证，结论是：$t=\pm\ii$ 给出与 $\tau$ 无关的 B-模型。\src{B p.70}

把特殊值列成表（表~\ref{tab:psi-2d}），方向务必记清：
\begin{table}[htbp]
\centering
\caption{四维参数 $t$、典范参数 $\Psi$ 与二维理论基础模型的对应。$\Psi=\infty$ 对应 B-模型（复结构 $J$），$\Psi=0$ 对应 A-模型（辛结构 $K$），方向不要记反。}
\label{tab:psi-2d}
\begin{tabular}{@{}llll@{}}
\toprule
$t$ & $\Psi$ & 二维理论 & 备注 \\
\midrule
$\pm\ii$ & $\infty$ & B-模型，复结构 $J$（$t=-\ii$ 为 $-J$） & 与 $\tau$ 无关 \\
$\pm1$（且 $\operatorname{Re}\tau=0$） & $0$ & A-模型，辛结构对应 $K$ & 镜像对称的另一端 \\
其他实 $t$ & 实数 & A-模型 & 一般 A-模型 \\
$t=0,\infty$ & 退化 & 退化极限 & 非平凡瞬子在此附近 \\
一般 $t$ & 复数 & 广义复几何（等价于带 $B$-场的 A-模型） & 只有结论 \\
\bottomrule
\end{tabular}
\end{table}
\src{B pp.61--62}\src{B pp.68--71} 这里有一个方向容易记反，特别标出：\emph{$\Psi=\infty$ 对应 B-模型，$\Psi=0$ 对应 A-模型}；原文 p.70 正是从 $\Psi=\infty$ 的 B-模型出发做 $S$-对偶的。

\subsection{$S$-对偶 = 镜像对称：本讲核心}

现在到了本讲的核心。四维理论有 $S$-对偶：$G\leftrightarrow\Ldual$，$\tau\mapsto-1/\tau$。它在扭化族上的作用是
\begin{equation}
S:\ \Psi\longmapsto-\frac{1}{n_g\Psi},\qquad
T:\ \Psi\longmapsto\Psi+1,
\label{eq:S-T-2}
\end{equation}
矩阵形式为 $S=\begin{pmatrix}0&1/\sqrt{n_g}\\ -\sqrt{n_g}&0\end{pmatrix}$、$T=\begin{pmatrix}1&1\\0&1\end{pmatrix}$；$n_g$ 是第 3 讲介绍过的花边数，单边（ADE）时 $n_g=1$，$B,C,F_4$ 时 $n_g=2$，$G_2$ 时 $n_g=3$。\src{B p.70} 注意 $S$ 交换 $\Psi=0$ 与 $\Psi=\infty$：这正好把 A-模型端与 B-模型端互换。

把 $S$-对偶用在 $\Psi=\infty$ 的 B-模型上：对群 $G$，这是 $\MH(G)$ 上复结构 $J$ 的 B-模型；做 $S$-对偶后群变成 $\Ldual$，$\Psi$ 变成 $0$，二维理论变成 $\MH(\Ldual)$ 上的 A-模型。仔细跟踪复结构，可算出它对应的是辛结构 $K$，于是
\begin{equation}
\text{B-模型}_J\bigl(\MH(G)\bigr)\;\cong\;
\text{A-模型}_K\bigl(\MH(\Ldual)\bigr).
\label{eq:S-B-A}
\end{equation}
\src{B p.70} 这已经是一个镜像型陈述，但两边用的复结构不同。再做一次超凯勒旋转
\begin{equation}
\mathcal J:\ \varphi\longmapsto \ii\varphi,
\label{eq:J-rotation}
\end{equation}
它把 $K$ 转成 $J$，于是 \eqref{eq:S-B-A} 变成
\begin{equation}
\text{B-模型}_J\bigl(\MH(G)\bigr)\;\cong\;
\text{A-模型}_J\bigl(\MH(\Ldual)\bigr).
\label{eq:mirror-J}
\end{equation}
这正是「$\MH(G)$ 与 $\MH(\Ldual)$ 在复结构 $J$ 下互为镜像」的精确含义。$\mathcal J$ 不是扭化族的一员（它把 $\PP^1$ 上的点转出族外），原文说它「不十分自然」，但作为超凯勒流形上的旋转是合法的。\src{B pp.70--71}

为什么 $S$-对偶会给出镜像对称？几何上的原因是 Hitchin 纤维化。$\MH$ 有 Hitchin 纤维化 $\MH\to\mathcal B$，底 $\mathcal B$ 参数化 Higgs 场的特征多项式，纤维是阿贝尔簇（Jacobian 型环面）。$S$-对偶在底上作用平凡（单边时甚至 $\mathcal B\cong{}^L\mathcal B$），在纤维上则是 T-对偶：纤维与其对偶纤维互换。底加纤维的 T-对偶正是 SYZ 意义下的镜像对称。原文用 B-模型的局部算符验证了这一点：局部算符只测量底上的信息，B-模型的手征环等于全纯函数环，而全纯函数在紧纤维上是常数、故来自底，$S$-对偶给出 $R(G)\cong R(\Ldual)$、$\mathcal B\cong{}^L\mathcal B$；纤维上的零膜在 $S$-对偶下变成支撑在中间维数纤维上的 $(B,A,A)$-膜，其自由度正是对偶环面。\src{B pp.71--77} 这一整套是物理推导：原文明确说「相信」每个二维拓扑场论由两个广义复结构之一决定、与另一个无关，并说完整证明尚不存在（\src{B pp.66--67}）。SYZ 图像里的映射 $\Xi:\mathcal B\to{}^L\mathcal B$ 通常也假定为平凡映射。\src{B p.77}

把上面的两步画成交换图，就是本讲的核心等式（图 \ref{fig:core}）：
\begin{equation}
\text{4d TFT}_\Psi(G)\ \text{在}\ \Sigma\times C\ \text{上}
\quad\xrightarrow{\ \text{紧化}\ }\
\text{2d }\sigma\text{-模型，靶空间 }\MH(G).
\label{eq:core-eq}
\end{equation}
在 $\Psi=\infty$ 处，右边就是复结构 $J$ 上的 B-模型。
\begin{figure}[htbp]
\centering
\begin{tikzcd}[column sep=4.8em, row sep=3.8em]
  \text{4d TFT}_{\Psi=\infty}(G)\ \text{on}\ \Sigma\times C
    \arrow[r, "\text{紧化}"]
    \arrow[d, "\text{$S$-对偶}"']
  & \text{2d }B\text{-模型 on }(\MH(G),J)
    \arrow[d, "\text{镜像}"] \\
  \text{4d TFT}_{\Psi=0}(\Ldual)\ \text{on}\ \Sigma\times C
    \arrow[r, "\text{紧化}"']
  & \text{2d }A\text{-模型 on }(\MH(\Ldual),J)
\end{tikzcd}
\caption{本讲核心交换图。上排是群 $G$、$\Psi=\infty$ 的 B-模型；下排是 $\Ldual$、$\Psi=0$ 的 A-模型。左竖边是四维 $S$-对偶（$\Psi\mapsto-1/(n_g\Psi)$），右竖边是二维镜像对称。两条路径——先紧化再镜像，或先 $S$-对偶再紧化——给出同一个二维理论；交换图的右边用 \eqref{eq:J-rotation} 的 $\mathcal J$ 把辛结构 $K$ 换成 $J$。}
\label{fig:core}
\end{figure}

为什么两条路径给出同一个二维理论？物理上，$S$-对偶是四维理论的等价，紧化只是一个操作，二者可交换；所以交换图在四维一侧自动成立。真正有内容的是右边那条竖边：它把四维的 $S$-对偶翻译成了二维的镜像对称。换句话说，不是先证明镜像对称再用它，而是把 $S$-对偶紧化，得到镜像对称。这就是 Kapustin--Witten 的路线：几何 Langlands 的对偶性被读成四维规范理论的 $S$-对偶，而后者在紧化后变成二维镜像对称。\src{B pp.70--71}

最后照例提醒：这里的每一步都是物理推导，不是数学证明；原文只做非分歧情形；2024 年前后 Gaitsgory 等人的代数几何工作才给出几何 Langlands 的严格证明。\src{B p.77}

\begin{keybox}
紧化把四维扭化理论变成以 $\MH$ 为靶空间的二维 sigma 模型：$\Psi=\infty$ 时它是复结构 $J$ 的 B-模型，$\Psi=0$ 时它是辛结构 $K$ 的 A-模型，而四维 $S$-对偶把 B-模型换成 A-模型，于是 $\MH(G)$ 与 $\MH(\Ldual)$ 在 $J$ 下互为镜像。
\end{keybox}

\begin{follow}
\begin{itemize}
\item \textbf{B §4, pp.38--45}：紧化到 $\Sigma\times C$ 与靶空间 $\MH$ 的论证（重点 pp.38--39；pp.40--45 是 Hitchin 模空间的基本性质）。
\item \textbf{B §5.1, pp.56--61}：二维扭化理论、$J_+,J_-$ 与 A-模型/B-模型的判据。
\item \textbf{B §5.2, pp.61--68}：$t$ 到 $(w_+,w_-)$ 的映射、$\Psi$ 与广义复几何（本讲只取结论的部分）。
\item \textbf{B §5.3, pp.68--71}：$\Psi=\infty$ 的特化、$S$-对偶与镜像对称（本讲核心）。
\end{itemize}
\end{follow}

\begin{checkbox}
\begin{itemize}
\item 为什么 $C$ 必须小、$\Sigma$ 必须大？若把两个尺寸互换，低能理论会变成什么？
\item 为什么平坦 $G_{\mathbb C}$ 丛在每一点 $x\in\Sigma$ 的限制给出 $\MH$ 的点？$\MH$ 的奇点对应什么？
\item 用 \eqref{eq:bps2d} 说明 $J_+=J_-$ 与 $J_+=-J_-$ 分别给出 B-模型与 A-模型。
\item 验证 $t=\pm\ii$ 对应 $\Psi=\infty$、$t=\pm1$ 且 $\operatorname{Re}\tau=0$ 对应 $\Psi=0$；说明为什么这个方向容易记反。
\item 写出 $S$ 对 $\Psi$ 的作用与 $n_g$ 的定义；单边时 $S$ 是什么？
\item 交换图四条边各是什么？两条路径为什么给出同一个二维理论？
\end{itemize}
\end{checkbox}

% ======================================================================
%  05_branes.tex —— 第 5 讲 膜、线算子、特征膜
%  依据：Kapustin--Witten 2006 \S6--\S8（印刷 pp.78--123）
% ======================================================================

\section{膜、线算子、特征膜}
\label{sec:branes}

\begin{goals}
\begin{itemize}
\item 说清 2d sigma 模型的边界条件为什么就是膜，A-膜与 B-膜各保留哪个超荷、各要什么样的几何（Lagrange/余各向同性 对 复子流形）。
\item 掌握 Wilson 算子与 't~Hooft 算子的定义差别（order 对 disorder），以及它们各自出现在哪个 $\Psi$ 值。
\item 理解特征膜的定义 $XB\cong B\otimes V$：为什么「特征值」必须换成向量空间，为什么零膜是 Wilson 的特征膜、Hitchin 纤维是 't~Hooft 的特征膜，以及全体特征膜如何合成 $C$ 上一条局部系统。
\end{itemize}
\end{goals}

前几讲已经把舞台搭好：四维 $\mathcal N=4$ 超杨--米尔斯理论紧化在曲线 $C$ 上，低能有效理论是一个以 $\MH$ 为靶空间的二维 sigma 模型，$\MH$ 超凯勒、带三个复结构 $I,J,K$（\src{B p.60}）。本讲放三样东西：膜、线算子、特征膜（内容据原文 \S6--\S8，\cite{kw06}）。先说膜分几类；再说 Wilson 与 't~Hooft 线算子如何降成二维算子；最后说算子作用在膜上何时「本征」，以及全体特征膜怎样拼出 $C$ 上的局部系统——那是几何 Langlands 的伽罗瓦一侧在物理里的第一次显形。

\subsection{边界条件就是膜：A 型与 B 型}

二维拓扑场论最容易被忽略的一点是：有边界时，理论的一半数据在边界上。世界面 $\Sigma$ 带边界时，边界条件要指定三样东西：场 $\Phi:\Sigma\to X$ 在 $\partial\Sigma$ 上落在哪个子流形 $N\subset X$ 里；边界上带什么 Chan--Paton 丛与联络；法向导数满足什么一阶条件。超对称把这些数据卡得很死：要让边界保留某一个超荷 $Q$，$(N,\text{丛})$ 就必须满足一条几何方程。选不同的 $Q$，就得到不同类型的膜（\src{B p.77}--\src{B p.78}）。

先看 A 型。A 型超荷的边界项正比于 $\int_{\partial\Sigma}\omega(\delta\Phi,\partial_t\Phi)$，要它为零，必须 $\omega|_N=0$，即 $N$ 是 Lagrange 子流形；同时边界规范场必须平坦，$F=0$。把 Lagrange 放宽成\term{coisotropic}{余各向同性}子流形（$TN^{\omega}\subset TN$）并配上横截的复结构，也能给出 A-膜；这是后文第 7 讲真正要用的情形。A 型超荷的上同调要数世界面瞬子（全纯圆盘），计算它用的是 Floer 型或量子上同调的技术。

再看 B 型。B 型超荷里含 $\bar\partial$，保它要求 $N$ 是复子流形、Chan--Paton 丛是全纯的。B 型超荷的上同调就是 Dolbeault 上同调，世界面瞬子不贡献修正。所以一句话：A 型要 Lagrange（或余各向同性）加平坦丛，保的是「辛」的那个超荷；B 型要复子流形加全纯丛，保的是「复」的那个超荷。

$\MH$ 是超凯勒的，三个复结构各有一对 A 型与 B 型超荷，因此一张膜可以「对 $I$ 是 B 型、对 $J$ 是 A 型」。真正有意思的是同时保留二维理论一半超对称的膜（这也等于保留四维理论四分之一超对称）：这样的膜只有四种基本类型，记作
\[
(B,B,B),\qquad (B,A,A),\qquad (A,B,A),\qquad (A,A,B),
\]
括号里三个字母依次表示该膜在 $I,J,K$ 中的类型（\src{B p.78}）。它们与两个常用例子一起列在表~\ref{tab:brane-types} 里。

\begin{table}[htbp]
\centering
\small
\caption{四类保留二维一半超对称的膜。括号中字母依次对应复结构 $I,J,K$；末列是典型例子（据 \src{B p.78} 与 \src{B p.72} Table~2）。}
\label{tab:brane-types}
\begin{tabular}{@{}lcccc@{}}
\toprule
类型 & $I$ & $J$ & $K$ & 典型例子 \\
\midrule
$(B,B,B)$ & B & B & B & 零膜；空间充满膜（全纯丛） \\
$(B,A,A)$ & B & A & A & Hitchin 纤维；固定 $E$ 让 $\varphi$ 变 \\
$(A,B,A)$ & A & B & A & $J$ 中的复 Lagrange 子流形 \\
$(A,A,B)$ & A & A & B & $K$ 中的复 Lagrange 子流形 \\
\bottomrule
\end{tabular}
\end{table}

两个例子要记牢。第一个是\term{zerobrane}{零膜} $B_x$：支撑在点 $x\in\MH$ 上，Chan--Paton 丛是支撑在 $x$ 的摩天楼层；点在任何复结构下都是复子流形，丛在点上全纯，故零膜是 $(B,B,B)$ 型。第二个是 Hitchin 纤维 $F$：它在复结构 $I$ 中是复子流形，并且对全纯辛形式 $\Omega_I$ 是 Lagrange 的（因而对 $\omega_J,\omega_K$ 也是），配上平凡丛，故是 $(B,A,A)$ 型（\src{B p.78}--\src{B p.79}）。

另外两条。其一，A-膜的支撑维数至少是靶空间的中维：Lagrange 最常见，余各向同性的支撑可以更高——第 7 讲构造 D-模时用的正是这个极端情形。其二，S-对偶把 $(B,B,B)$ 换成 $(B,A,A)$，把 $(A,B,A)$ 与 $(A,A,B)$ 各自映到自身（\src{B p.72}）；这个交换本讲末尾会再次出现。

\subsection{线算子：Wilson 与 't~Hooft}

四维理论里，算子可以支撑在任意子流形上。支撑是一维时叫线算子。理论是拓扑的，把支撑连续形变不改变算子，于是线算子只依赖支撑的同伦类：支撑是闭圈时叫圈算子，支撑是一条跑到无穷远的线时叫线算子。路径积分里插入线算子，相当于在它周围挖掉一个小管、在管的边界上指定某种数据；指定方式不同，算子就不同。两种最重要的方式是 Wilson 与 't~Hooft，分别出现在对称性增强的两个特殊值 $\Psi=\infty$ 与 $\Psi=0$（\src{B p.78}--\src{B p.79}）。

Wilson 算子在 $\Psi=\infty$（即 $t=\pm\ii$）。取规范群的表示 $R$，沿圈 $S$ 取 holonomy 的迹：
\begin{equation}
W_R(S)=\Tr_R\,P\exp\Big(-\oint_S\mathcal A\Big),\qquad \mathcal A=A+\ii\varphi,
\label{eq:wilson}
\end{equation}
其中 $A$ 是规范场，$\varphi$ 是复标量。为什么只用这个组合：在 $t=\pm\ii$ 处，拓扑扭化的变换把 $A$ 与 $\varphi$ 一起旋转，唯一保持不变的复化联络就是 $A+\ii\varphi$。为什么它是拓扑的：在 $t=\pm\ii$，BPS 方程蕴含 $\mathcal A$ 平坦，于是 holonomy 只依赖 $S$ 的同伦类，把 $S$ 挪一挪不改变 $W_R(S)$。开线也不浪费：从 $p$ 到 $q$ 的开线给出一个平行移动矩阵，可以在两端之间搬运表示（\src{B p.79}--\src{B p.80}）。

't~Hooft 算子在 $\Psi=0$（即 $t=1$、$\theta=0$）。它不是某个联络的迹，而是一个 disorder 算子：在支撑 $S$ 附近指定经典场的奇点，然后对「带这个奇点的场」做路径积分——算子本身就是这个受限制的路径积分。$G=\mathrm{U}(1)$ 时，磁荷 $m$ 的奇点就是 Dirac 单极，绕 $p$ 的小球面上 $\oint_{S^2}c_1(L)=m$（\src{B p.83}）。一般 $G$：奇点由同态 $\rho:\mathrm{U}(1)\to G$ 的共轭类指定；GNO 对偶说，这样的共轭类等价于 $\Ldual$ 的一个最高权，也等价于 $\Ldual$ 的一个不可约表示 ${}^{L}R$，于是把 't~Hooft 算子记作 $T({}^{L}R,S)$。这里还有一层离散数据：$\xi({}^{L}R)\in H^2(S^2,\pi_1(G))$ 度量算子的扭曲程度，它非平凡当且仅当 $\Ldual$ 的中心在 ${}^{L}R$ 上作用非平凡（\src{B p.84}）。

两个算子的对比值得背下来：Wilson 是 order 算子，它读出联络；'t~Hooft 是 disorder 算子，它指定奇点，算子的值是奇点存在下的路径积分。S-对偶 $\Psi\mapsto-1/(n_g\Psi)$ 把 $\Psi=\infty$ 与 $\Psi=0$ 互换，于是把 Wilson 与 't~Hooft 互换，并把表示 $R$ 换成 ${}^{L}R$（\src{B p.98}）。这个互换是本讲后半段与下一讲的引擎。

\subsection{线算子作用在膜上：特征膜}

现在让线算子作用在膜上。做法是「折叠」：线算子支撑在 $\Sigma$ 上的一条线（再乘上 $C$ 上的点 $p$）；把 $\Sigma$ 沿这条线剪开再翻折，得到一个界面——它是两套二维理论的公共边界条件，当然也是膜，准确说是一个 $(Y,Z)$-膜。把线算子 $X$ 变成界面，这一步就是把算子「打开」成一个边界条件。膜 $B$ 与界面复合，结果仍是膜，记作 $XB$：线算子把膜映到膜（\src{B p.90}）。

用态空间说更准确：设 $H_{B_1,B_2}$ 是两端带边界条件 $B_1,B_2$ 的区间上的态空间（边界算子的空间），则 $XB$ 的作用由 $H_{\tilde B,XB}$ 刻画；$H_{B,B}$ 是结合代数（边界算子的 OPE），$H_{B_1,B_2}$ 是双模（\src{B p.91}--\src{B p.92}）。

接下来定义本讲的主角。普通的特征方程 $O\psi=\lambda\psi$ 在这里没有意义：算子作用的对象是膜（一个范畴的对象），不是向量，没有「本征值」可谈。要类比，就得把 $\lambda$ 换成同样能「张量上去」的东西——一个向量空间。先定义膜与向量空间的张量积 $B\otimes V$：几何上就是把 Chan--Paton 丛 $W$ 换成 $W\otimes V$（$V=\CC^k$ 时就是取 $k$ 份膜）；抽象地由
\begin{equation}
H_{\tilde B,B\otimes V}=H_{\tilde B,B}\otimes V
\label{eq:tensor}
\end{equation}
刻画。若
\begin{equation}
XB\cong B\otimes V,
\label{eq:eigenbrane}
\end{equation}
就说 $B$ 是 $X$ 的\term{eigenbrane}{特征膜}，$V$ 是它的「特征值」（\src{B p.95}）。

这个定义有两点要紧。第一，特征值不是数而是向量空间；一旦有了向量空间，「特征值随参数变化」才有意义。第二，式~\eqref{eq:eigenbrane} 正是 Beilinson--Drinfeld 在几何 Langlands 里定义的 eigensheaf 的物理版；这不是巧合，而是第 6 讲要还的债。

第一个例子（电）：Wilson 算子的特征膜是零膜。Wilson 算子 $W_q(R)$ 作用在膜上的方式，是把它的 Chan--Paton 丛 $U$ 换成 $U\otimes E_q(R)$，其中 $E_q(R)$ 是从 $C$ 上的通用丛诱导、在表示 $R$ 下取纤维的丛。零膜 $B_x$ 的 Chan--Paton 丛是支撑在 $x$ 的摩天楼层 $U_x$；任何丛限制到一个点都平凡，于是
\begin{equation}
U_x\otimes E_q(R)=U_x\otimes E_q(R)|_x,
\label{eq:zerobrane}
\end{equation}
右边只是 $x$ 处的一个固定向量空间。故 $W_q(R)B_x\cong B_x\otimes E_q(R)|_x$：零膜是电特征膜，特征值是「表示空间在 $x$ 处的纤维」（\src{B p.121}）。

第二个例子（磁）：'t~Hooft 算子的特征膜是 Hitchin 纤维，准确说是「卷在纤维上、带秩 1 平坦 Chan--Paton 丛」的膜，原文叫 type~F 膜。它和零膜的关系正是 S-对偶：$\MH(\Ldual,C)$ 里的零膜，经 S-对偶变成 $\MH(G,C)$ 里 Hitchin 纤维上的 type~F 膜。为什么非要从点换成纤维？因为电与磁「局域化」的方向相反：零膜把膜局域在一个点，读出电表示；磁荷的信息沿 Hitchin 纤维的环面方向铺开，type~F 膜正好卷在整条纤维上，它的平坦 Chan--Paton 丛给出对偶环面上的一个点，这个点就是 't~Hooft 算子的「特征值」。原文还用 OPE 印证这件事：Wilson 算子的 OPE 生成 $G$ 的表示环（原文 (6.27)），而 't~Hooft 算子的 OPE 应该生成 $\Ldual$ 的表示环（原文 (6.29)；原文在这里写的是 “we expect”）（\src{B p.98}--\src{B p.99}）。

还有一条结构性事实：从四维来的线算子自动两两交换。物理理由很直接：它们支撑在 $C$ 上的不同点，四维里根本不相交，可以随意前后移动；$C$ 是二维的，两个不同点在空间中的位置足够挪开而不产生单值性。既然算子交换，就可以同时找它们的特征膜——这是下一讲把整个算子族放在一起讨论的前提（\src{B p.97}）。

\subsection{全体特征膜 = 曲线上的局部系统}

最后把所有特征膜的「特征值」沿 $C$ 拼起来。线算子可以放在任意 $p\in C$，记 $X_p$。若 $B$ 对某个 $p$ 是特征膜，则对所有 $p$ 都是：
\begin{equation}
X_pB\cong B\otimes V_p,\qquad p\in C .
\label{eq:family}
\end{equation}
单看一个 $p$，$V_p$ 是一个向量空间；但 $p$ 可以连续移动。移动 $p$ 相当于改变 $C$ 的度规，作用量只变一个 $Q$-恰当项，故可在族 $V_p$ 上装一个平坦联络，使 $V_p$ 局部与 $p$ 无关；整体上可以有非平凡单值性。用几何语言：$p\mapsto V_p$ 是 $C$ 上一个平坦向量丛（\src{B p.96}）。

不同表示的数据还要合起来。OPE 给出
\begin{equation}
W(R)W(\tilde R)=\sum_\alpha n_\alpha\,W(R_\alpha),\qquad R\otimes\tilde R=\bigoplus_\alpha n_\alpha R_\alpha,
\label{eq:wilson-ope}
\end{equation}
特征值必须满足同样的分解，于是 $V_p(R)\otimes V_p(\tilde R)\cong V_p(R\otimes\tilde R)$；用范畴的话说，特征膜给出一个张量函子 $\mathrm{Rep}(G)\to\mathrm{Vect}$。Tannaka 对偶说，这样的张量函子等价于一个局部系统。电侧（零膜）得到的是取值在 $G_{\CC}$ 里的局部系统；磁侧（type~F 膜）得到的是取值在 $\LdualC$ 里的局部系统
\begin{equation}
\rho:\pi_1(C)\longrightarrow \LdualC .
\label{eq:local-system}
\end{equation}
这就是「全体特征膜合起来等于 $C$ 上一条局部系统」的准确含义（\src{B p.96}--\src{B p.98}）。

为什么说这已经像几何 Langlands？几何 Langlands 的标准故事正好从这里开始。给定 $\Ldual$ 与同态 $\vartheta:\pi_1(C)\to\LdualC$（并要求它半稳定），得到 $\MH(\Ldual,C)$ 中一点 $x(\vartheta)$；$\vartheta$ 不可约时该点光滑，其零膜是电特征膜；S-对偶把它变成 $\MH(G,C)$ 中 Hitchin 纤维上的 type~F 磁特征膜（\src{B p.122}）。几何 Langlands 断言，这个 $\vartheta$ 应该对应 $\MC$ 上一个 Hecke 特征层——既是 Hecke 算子的特征对象，又是 D-模。物理侧现在有了「磁特征膜」这一候选对象，但还没有证明它就是 Hecke 特征层；这一步正是第 6 讲的任务。原文在这里的措辞很克制：它只说我们「至少可以看到一线微光」（at least a glimmering）（\src{B p.122}--\src{B p.123}）。

所以本讲的收获是：线算子把膜映到膜，只乘一个向量空间的膜是特征膜，两类特征膜的特征值沿 $C$ 拼成局部系统。语言一旦立住，剩下的问题就只有一个：'t~Hooft 算子到底是不是数学上的 Hecke 算子？

\begin{keybox}
线算子把膜映到膜；特征膜是满足 $XB\cong B\otimes V$ 的膜，「特征值」是向量空间而非数；电侧特征膜是零膜，磁侧是 Hitchin 纤维上的 type~F 膜，其特征值沿 $C$ 拼成局部系统——几何 Langlands 的伽罗瓦一侧在物理里第一次显形。
\end{keybox}

\begin{follow}
\begin{itemize}
\item B \S6, pp.78--102：线算子的定义（(6.4)--(6.5)）、't~Hooft 的奇点与 GNO 对偶（(6.9) 附近）、折叠与特征膜、OPE。第一遍抓住 (6.22) 与 (6.27) 即可。
\item B \S7, pp.103--115：中心、离散 flux 与 twisted 丛，可跳读，只记住「电/磁 flux 成对、S 交换它们」。
\item B \S8, pp.116--123：零膜 $=$ 电特征膜（(8.10)）、type~F 膜 $=$ 磁特征膜、原文对几何 Langlands 的「一线微光」；这是第 6 讲的直接入口。
\end{itemize}
\end{follow}

\begin{checkbox}
\begin{enumerate}
\item 为什么 A-膜要求 $\omega|_N=0$（或余各向同性）而 B-膜要求 $N$ 是复子流形？它们各自保留哪个超荷，对应的上同调分别是什么？
\item Wilson 算子为什么只依赖圈的「同伦类」？'t~Hooft 算子为什么不是某个联络的 holonomy，而是一个 disorder 算子？
\item 特征膜的定义里为什么必须用向量空间代替数？写出 $XB\cong B\otimes V$ 的态空间刻画，并说明 $B\otimes V$ 在几何上到底是什么操作。
\item 零膜为什么是 Wilson 算子的特征膜？把「任何丛限制到点都平凡」这一步补全。
\item Hitchin 纤维为什么是 't~Hooft 算子的特征膜？它和零膜通过什么操作联系起来？为什么磁侧非要从点换成纤维？
\item 电侧、磁侧两组特征膜分别给出什么群里的局部系统？几何 Langlands 用的是哪一侧，为什么？
\end{enumerate}
\end{checkbox}

% ======================================================================
%  06_hecke.tex —— 第 6 讲 核心等式：'t Hooft 算子 = Hecke 算子
%  依据：Kapustin--Witten 2006 \S9--\S10（印刷 pp.123--184）
% ======================================================================

\section{核心等式：'t~Hooft 算子 = Hecke 算子}
\label{sec:hecke}

上一讲结尾留下一个问题：'t~Hooft 算子的特征膜（Hitchin 纤维上的 type~F 膜）与几何 Langlands 的 Hecke 特征层，是不是同一个东西？本讲回答它。原文第 9、10 节分两步走：先把 't~Hooft 算子与数学上的 Hecke 算子对上（\src{B p.123}--\src{B p.148}），再把 Hecke 修改的空间实现为 Bogomolny 方程的解空间（\src{B p.148}--\src{B p.184}）。两步合起来，就是全文的核心等式（\cite{kw06}）。

\begin{goals}
\begin{itemize}
\item 掌握 Hecke 修改的定义，并知道它的几何化身是仿射 Grassmannian 的 Schubert 胞腔。
\item 理解 't~Hooft 算子的奇点数据与 Hecke 修改的权是同一组整数（U(1) 情形就是同一个 twist $m$），以及 Bogomolny 方程如何把「修改空间」实现为「解空间」。
\item 知道单极子冒泡、Hecke 特征层的标准陈述，以及这条推导里哪些步骤原文只写「应该」。
\end{itemize}
\end{goals}

\subsection{Hecke 修改：在一点把丛改一下}

先讲数学侧。设 $C$ 是紧黎曼面，$E_-$ 是 $C$ 上 $G$-丛，$p_0\in C$。一个\term{Hecke modification}{Hecke 修改}是一对 $(E,\sigma)$：$E$ 是 $C$ 上 $G$-丛，$\sigma:E|_{C\setminus\{p_0\}}\simeq E_-|_{C\setminus\{p_0\}}$ 是同构，而且 $\sigma$ 一般不能延拓到 $p_0$。直观地说，就是「在 $p_0$ 处把丛改一下」：$p_0$ 之外两个丛一样，$p_0$ 处允许跳变。这正是「点状缺陷」的数学语言，也是 't~Hooft 算子所指的东西（\src{B p.145}）。

局部描述最清楚。取 $C=\PP^1$、$p_0=0$，用两个开集覆盖：$U_\infty=\PP^1\setminus\{0\}$ 与 $U_0=\PP^1\setminus\{\infty\}$，交叠是 $\CC^*$。$E_-$ 在 $U_\infty$ 上给平凡化 $\sigma$；$E$ 在 $U_0$ 上给平凡化 $\sigma'$。在交叠上两者相差 $g(z)=\sigma'\sigma^{-1}$，它是 $\CC^*$ 上全纯、并在 $\PP^1$ 上可亚纯延拓的 $\mathrm{GL}(N,\CC)$-值函数。换平凡化 $\sigma'$ 相当于右乘 $h_0\in\mathrm{GL}(N,\mathcal O)$，其中 $\mathcal O$ 是 $U_0\simeq\CC$ 上全纯函数的环。于是带固定平凡化的修改后丛的参数空间是
\begin{equation}
\mathrm{GL}(N,\mathcal X)/\mathrm{GL}(N,\mathcal O)=\mathrm{Gr}_N,
\label{eq:grn}
\end{equation}
其中 $\mathcal X$ 是 $\CC^*$ 上全纯、在 $\PP^1$ 上可亚纯延拓的函数环。这个空间叫\term{affine Grassmannian}{仿射 Grassmannian}，也等于形式版本 $\mathrm{GL}(N,\CC((z)))/\mathrm{GL}(N,\CC[[z]])$（\src{B p.146}）。

权的定义如下。若修改后的丛 $E_+$ 由截面 $s_j=z^{-m_j}f_j$ 生成（$f_j$ 是 $E_-$ 的全纯截面基，$m_j\in\ZZ$），则交叠上的转移矩阵是 $g_{ji}(z)=z^{-m_j}f_{ji}(z)$。最简单的情形 $f_{ji}=\delta_{ji}$，于是 $g(z)=\operatorname{diag}(z^{-m_1},\dots,z^{-m_N})$。这组整数 ${}^{L}w=(m_1,\dots,m_N)$（模置换）就叫权。关键观察：$g(z)$ 限制在 $\CC^*$ 上恰是一个同态 $\CC^*\to\mathrm{GL}(N,\CC)$，它正是
\begin{equation}
\rho:\mathrm{U}(1)\to\mathrm{U}(N),\qquad \ee^{\ii\alpha}\mapsto\operatorname{diag}(\ee^{\ii m_1\alpha},\dots,\ee^{\ii m_N\alpha})
\label{eq:rho}
\end{equation}
的复化（\src{B p.147}）。记住这个 $\rho$，下一小节它就变成 't~Hooft 算子的磁荷。

固定权 ${}^{L}w$ 后，所有这种形状的点构成一个轨道 $Y({}^{L}w;p_0)$，它是 $\mathrm{GL}(N,\mathcal O)$ 在 $\mathrm{Gr}_N$ 中的轨道；这些轨道把 $\mathrm{Gr}_N$ 分层，叫 Schubert 胞腔。紧化取轨道闭包，叫 Schubert cycle；它一般奇异，由 $Y({}^{L}w)$ 与若干低维轨道组成。为什么轨道有限维：$\mathrm{Gr}_N(i)=\{g:\text{在 }0\text{ 处极点阶}\le i\}$ 有限维、紧、$\mathrm{GL}(N,\mathcal O)$-不变，每个轨道都落在某个 $\mathrm{Gr}_N(i)$ 里。一般 $G$：$\mathrm{Gr}_G=G_{\CC}(\mathcal X)/G_{\CC}(\mathcal O)$，$Y(\rho)$ 是 $\rho_{\CC}:\CC^*\to G_{\CC}$ 所在的轨道；更一般地，'t~Hooft 算子的权是 $\Ldual$ 的权。修改空间的维数是 $\Delta_{{}^{L}w}=\sum_{i<j}(a_i-a_j)$，其中 $a_i$ 是权的分量（\src{B p.148}）。

例子：${}^{L}w=(1,0)$ 给 $\PP^1$；${}^{L}w=(2,0)$ 给加权射影空间 $\PP(1,1,2)$（它是 $A_1$ 奇点 $\CC^2/\ZZ_2$ 的解消），自然紧化须再加入 $(1,1)$ 型极限点；minuscule 权 ${}^{L}w(k)=(1,\dots,1,0,\dots,0)$ 给 Grassmannian $\mathrm{Gr}(k,N)$，光滑紧（\src{B p.138}--\src{B p.144}）。

\subsection{'t~Hooft 的奇点结构就是 Hecke 的权}

现在看物理侧。$\Psi=0$，把四维理论放在 $M=\RR\times W$ 上，$W=I\times C$；静态 't~Hooft 算子支撑在 $\RR\times p_0$（\src{B p.123}）。按上一讲，'t~Hooft 算子是「在 $p_0$ 附近指定经典场的奇点、再对带奇点的场做路径积分」。奇点的类型：$\mathrm{U}(1)$ 是磁荷 $m$；$\mathrm{U}(N)$ 是同态 $\rho:\mathrm{U}(1)\to\mathrm{U}(N)$，$\ee^{\ii\alpha}\mapsto\operatorname{diag}(\ee^{\ii m_1\alpha},\dots,\ee^{\ii m_N\alpha})$，即权 ${}^{L}w=(m_1,\dots,m_N)$（模置换）；一般 $G$ 是 $\rho$ 的共轭类，GNO 对偶又把它等同于 $\Ldual$ 的不可约表示 ${}^{L}R$，故记 $T({}^{L}R)$（\src{B p.131}）。

注意：这与上一小节的式~\eqref{eq:rho} 是同一个 $\rho$。Hecke 修改的转移矩阵 $\operatorname{diag}(z^{-m_j})$ 正是这个 $\rho$ 的复化。于是「't~Hooft 的磁荷」$=$「Hecke 修改的权」$=$「Schubert 胞腔的标号」：三者是同一组整数 $m_1,\dots,m_N$。这就是原文所说的基本联系（\src{B p.132}）。

为什么必须是同一组整数？物理上有一条硬论证。Bogomolny 方程的解在无算子的区间里沿 $y$ 方向不变：取规范 $A_y=0$，有 $\partial_yA_{\bar z}=-\ii D_{\bar z}\varphi_0$；于是 $E_y$ 的全纯型在相邻算子之间不变，只有跨过 't~Hooft 算子时才跳变。$\mathrm{U}(1)$ 情形，取小圆盘包住 $p_0$，用 $c_1$ 的积分得 $0=q_+-m-q_-$，故 $L_+\cong L_-\otimes\mathcal O(p_0)^m$：奇点的磁荷 $m$ 恰好等于输出丛相对输入丛的 twist 阶，拓扑守恒同时定住了两边（\src{B p.130}--\src{B p.131}）。非阿贝尔情形把 $E_y$ 在 $p_0$ 附近对角化 $E_y=\bigoplus L_i$，每个 $L_i$ 各自跳变 $L_i\mapsto L_i\otimes\mathcal O(p_0)^{m_i}$；$c_1$ 的变化是 $\sum m_i$，故保拓扑要求 $\sum m_i=0$。一般 $G$：$E_{R,w}\mapsto E_{R,w}\otimes\mathcal O(p_0)^{\langle{}^{L}w,w\rangle}$，其中 $\langle\cdot,\cdot\rangle$ 是权与余权的配对（\src{B p.133}--\src{B p.134}）。

把这些对应列成表~\ref{tab:correspondence}。

\begin{table}[htbp]
\centering
\small
\caption{'t~Hooft 算子与 Hecke 修改的对应：两侧由同一组整数标号（据 \src{B p.130}--\src{B p.134}）。}
\label{tab:correspondence}
\begin{tabular}{@{}lll@{}}
\toprule
物理：'t~Hooft 算子 & 数学：Hecke 修改 & 共同的整数 \\
\midrule
$\mathrm{U}(1)$：磁荷 $m$ & $L\mapsto L\otimes\mathcal O(p_0)^m$ & $m\in\ZZ$ \\
$\mathrm{U}(N)$：$\rho=\operatorname{diag}(\ee^{\ii m_i\alpha})$ & $L_i\mapsto L_i\otimes\mathcal O(p_0)^{m_i}$ & ${}^{L}w=(m_1\ge\dots\ge m_N)$ \\
一般 $G$：$\rho:\mathrm{U}(1)\to G$ 的共轭类 & 轨道 $Y(\rho)=$ Schubert 胞腔 & $\Ldual$ 的权 ${}^{L}w$（${}^{L}R$ 的最高权） \\
\bottomrule
\end{tabular}
\end{table}

\subsection{Bogomolny 方程：修改空间就是解空间}

标号对上了，剩下的是「作用」是否也对上。原文的答案是：把 Hecke 修改空间实现为 Bogomolny 方程的解空间。先写边界条件。$W=I\times C$，'t~Hooft 算子插在 $p_i=y_i\times z_i$。左端 $C_-$：固定 $E_-$（即指定 $A_-$，Dirichlet 边界条件），$\varphi_0$ 自由，规范变换在 $C_-$ 上取平凡——于是输入的是丛 $E_-$ 本身，而不是它的同构类。右端 $C_+$：不指定 $A$，而令 $\varphi_0=0$；Bogomolny 方程随即要求 $A$ 满足协变 Neumann 条件 $\iota_nF=0$。规范变换在 $C_+$ 上允许非平凡，故输出的是 $E_+$ 的同构类。这组边界条件是椭圆的；没有算子时，能量恒等式强制 $F=D\varphi_0=0$，唯一解是平凡解，解空间是一个点（\src{B p.150}--\src{B p.155}）。

有算子时，解空间 $Z$ 有两种描述。其一（辛商）：场空间 $(A,\varphi_0)$ 上有复结构 $I$ 与辛形式 $\omega$；矩映射 $\mu=0$ 给出 Bogomolny 方程，故 $Z=\mu^{-1}(0)/G$ 是辛商。其二（复商）：复 Bogomolny 方程的复商 $Z'=C_0/G_{\CC}$ 给出的正是 Hecke 修改空间——因为 $E_y$ 的全纯型沿 $y$ 不变、只在算子处跳变，跨算子时恰好做一次 Hecke 修改。两者共享同一个复结构：修改空间 $=$ 解空间（\src{B p.162}）。检查维数：椭圆复形的指标等于 $\dim Z$，算出来正是 $\Delta_{{}^{L}w}=\sum_{i<j}(a_i-a_j)$，与上一小节的公式一致，且对任意紧 $G$ 成立（Pauly--Kronheimer）（\src{B p.155}）。原文由此说，这给了 't~Hooft 与 Hecke 的关系「一个精确的框架」。图~\ref{fig:chain} 把这条链画成方块。

\begin{figure}[htbp]
\centering
\begin{tikzpicture}[font=\small,>=Stealth,
  bx/.style={draw=kwblue!70,fill=kwblue!4,rounded corners=2pt,align=center,
             text width=3.45cm,inner sep=5pt}]
\node[bx] (th) {'t~Hooft 算子\\ 磁荷 ${}^{L}w$，奇点 $\rho:\mathrm{U}(1)\to G$};
\node[bx,right=2.1cm of th] (hk) {Hecke 修改\\ 权 ${}^{L}w$，Schubert 胞腔 $Y({}^{L}w)$};
\node[bx,below=1.8cm of th] (bg) {Bogomolny 解空间\\ $Z=\mu^{-1}(0)/G$};
\node[bx,below=1.8cm of hk] (md) {Hecke 修改空间\\ $Z'=C_0/G_{\CC}$};
\draw[<->,thick,kwblue] (th) -- (hk) node[midway,above]{\scriptsize 核心等式：作用 $=$ 修改};
\draw[<->,thick,kwgray] (bg) -- (md) node[midway,below]{\scriptsize 同一个复流形};
\draw[<->,thick,kwgray] (th) -- (bg) node[midway,left]{\scriptsize 奇点$\mapsto$边界条件};
\draw[<->,thick,kwgray] (hk) -- (md) node[midway,right]{\scriptsize 修改空间$=$解空间};
\end{tikzpicture}
\caption{核心等式的三方块图。上边横边是核心等式：'t~Hooft 算子的作用 $=$ Hecke 修改；下边横边是它们共同的几何载体（同一个复流形）。左边竖边把算子的奇点数据翻成 Bogomolny 方程的边界条件，右边竖边把修改空间翻成解空间（\src{B p.162}）。}
\label{fig:chain}
\end{figure}

还有一层结构要说清：沿 $I$ 与沿 $C$ 移动的角色不同。把联络改成 $A'=A-\ii\varphi_0\dd y$，则 $A'_z$、$A'_y$ 在复结构 $I$ 中全纯；保纤维化 $W\to C$ 的微分同胚群保持 $I$、保持复 Bogomolny 方程与 $C_0$，因此 $Z$ 的复结构不随 $p_i$ 变；但辛形式与矩映射不做此不变，所以沿 $y$ 方向滑动 $p_i$ 只改变 $Z$ 的辛结构、不改变复结构，而改变 $z_i$ 只改变复结构、不改变辛结构（\src{B p.163}--\src{B p.164}）。这就是为什么「修改空间 $=$ 解空间」只在复结构的意义上成立——下一讲的 D-模正需要这个复结构。

接下来是冒泡。若两个 't~Hooft 算子的奇点重合，Bogomolny 解会出现\term{monopole bubbling}{单极子冒泡}：把 $\CC^2$ 看作 $\RR^4$，$\mathrm{U}(1)$ 作用在 $\CC^2$ 上，$\CC^2/\mathrm{U}(1)\simeq\RR^3$；等变瞬子对应 $\RR^3$ 上的 Bogomolny 解。瞬子有尺度，尺度趋于零时缩成一个点——磁荷集中到奇点并消失，这就是 bubbling。ADHM 复形给出纤维的特征标公式 $\chi(E_0')=\chi(E_0)+\chi(H_1)(2-\ee^{\ii\theta}-\ee^{-\ii\theta})$，对应的权变化是 $(m+1,m-1)\to(m,m)$（原文 (10.13)）（\src{B p.157}--\src{B p.158}）。这恰好是上一小节紧化 $Y({}^{L}w)$ 时必须加入的 associated weights：物理上的冒泡 $=$ 数学上的轨道闭包。它也解释了为什么 Schubert cycle 的奇异结构正好是超凯勒的：它来自 $\RR^4$ 上瞬子模空间的超凯勒结构（\src{B p.159}）。

最后是 OPE，只给结论不证。两个 't~Hooft 算子靠近时，物理上可作算子乘积展开；几何上是两个修改空间的卷积。原文的结论是：Hecke 算子的 OPE 代数就是 $\Ldual$ 的表示环，${}^{L}R\otimes{}^{L}R'=\bigoplus_\alpha N_\alpha\otimes{}^{L}R_\alpha$（原文 (10.21)）；解空间的（交叠）上同调恰好给出表示本身：$H({}^{L}R)\cong{}^{L}R$（原文 (10.22) 与 \src{B p.173}）；卷积的精确形式是 IC 层的分解
\begin{equation}
\pi_*\big(IC({}^{L}w)\star IC({}^{L}w')\big)=IC({}^{L}w+{}^{L}w')\oplus\bigoplus_\alpha \tilde N_\alpha\otimes IC({}^{L}w_\alpha),
\label{eq:convolution}
\end{equation}
其中系数 $\tilde N_\alpha$ 与张量积系数 $N_\alpha$ 相等，靠的是 Lusztig--Ginzburg--Mirkovi\'c--Vilonen 的数学定理（\src{B p.178}）。物理侧只得到首项；系数的证明在数学文献里（\src{B p.169}）。

\subsection{Hecke 特征层与几何 Langlands}

现在把整条链翻译成层论。数学上，Hecke 算子是作用在 $\mathrm{Bun}_G$ 上的层（或 D-模）范畴的函子：对 $x\in C$ 与 $\Ldual$ 的表示 $V$，Hecke 函子 $H_x^V$ 由 Hecke 对应 $\mathrm{Bun}_G\leftarrow \mathcal H_x\to\mathrm{Bun}_G$ 定义（$\mathcal H_x$ 参数化「在 $x$ 处相差一个 Hecke 修改」的丛对）：先沿对应拉回到中间空间、张量上「通用修改」给出的丛 $V$，再推回 $\mathrm{Bun}_G$。\term{Hecke eigensheaf}{Hecke 特征层}就是对它的特征对象：
\begin{equation}
H_x^V(\mathcal F)\cong V\otimes\mathcal F
\quad\text{（在 } C\times\mathrm{Bun}_G\text{ 上；限制到 }\{x\}\times\mathrm{Bun}_G\text{ 就是 }V\otimes\mathcal F\text{）}.
\label{eq:eigensheaf}
\end{equation}
用上一讲的语言：把「Hecke 算子 $\leftrightarrow$ 't~Hooft 算子」、「$\mathrm{Bun}_G$ 上的层 $\leftrightarrow$ 膜」对上，式~\eqref{eq:eigensheaf} 就是特征膜定义 $XB\cong B\otimes V$ 的数学版——「特征值」从向量空间变成局部系统的纤维。

几何 Langlands 的标准陈述是：对每个满足半稳定条件的 $\Ldual$-局部系统 $\rho:\pi_1(C)\to\LdualC$，存在 $\mathrm{Bun}_G$ 上一个非零 Hecke 特征层 $\mathcal F_\rho$，使得对每个 $x$ 与 $V$，Hecke 算子的作用等于张量上 $\rho$ 的纤维；并且 $\mathcal F_\rho$ 是 D-模（holonomic）。物理侧的对应物就是上一讲的图景：$\rho$ 给出 $\MH(\Ldual,C)$ 中的点 $x(\rho)$，其零膜是电特征膜；S-对偶把它变成 $\MH(G,C)$ 中 Hitchin 纤维上的 type~F 膜，即磁特征膜。于是「't~Hooft 算子的特征膜 $=$ Hecke 特征层」——本讲的标题等式，在层论层面就是几何 Langlands 的「伽罗瓦 $\to$ 自守」方向（\src{B p.122}）。

最后是诚实提醒，共四条。第一，本讲的每一步几乎都是物理推导，不是数学证明。原文自己大量使用 should、we expect：「't~Hooft OPE 生成 $\Ldual$ 表示环」是 we expect（\src{B p.99}）；「$H({}^{L}R)\cong{}^{L}R$ 的物理解释从物理观点看是神秘的」（\src{B p.173}）；extended Bogomolny 与任意 A-膜的作用只是 strongly suggests（\src{B p.181}），完整证明在别的文献（\src{B p.183}）。第二，原文只处理非分歧、半单、generic 的情形；分歧情形属于别的工作。第三，2024 年前后 Gaitsgory 等人的纯代数几何工作才把几何 Langlands 真正证出来；K--W 提供的是「为什么应该如此」的物理图像。第四，就连本讲最核心的「修改空间 $=$ 解空间」也只是复结构层面的等同（辛结构会随 $p_i$ 移动而变），原文对此很谨慎。所以读完这一讲，你应该能画出图~\ref{fig:chain}，并指着每一条边说出它为什么成立、哪一边原文没有证。

\begin{keybox}
同一组整数既是 't~Hooft 算子的磁荷（权 ${}^{L}w$），又是 Hecke 修改的权；Bogomolny 方程让 Hecke 修改空间成为解空间；于是 't~Hooft 算子的特征膜就是 Hecke 特征层——几何 Langlands 的物理版核心等式。
\end{keybox}

\begin{follow}
\begin{itemize}
\item B \S9.1, pp.123--135：'t~Hooft 与 Hecke 的初步比较；$\mathrm{U}(1)$、$\mathrm{U}(N)$、一般 $G$ 的权与作用。先读这一段，抓住 $L_i\mapsto L_i\otimes\mathcal O(p_0)^{m_i}$。
\item B \S9.2, pp.135--144：修改空间 $Y({}^{L}w)$ 的显式构造、紧化与 associated weights、维数与例子。
\item B \S9.3, pp.145--148：仿射 Grassmannian 与 Schubert 胞腔；原文自己说「本文并不真正使用这个框架」，可快读。
\item B \S10, pp.148--184：Bogomolny 的边界条件与指标（\S10.1）、冒泡（\S10.2）、K\"ahler 商（\S10.3）、OPE（\S10.4）；\S10.5 的 extended Bogomolny 可选，第 7 讲会再用。
\end{itemize}
\end{follow}

\begin{checkbox}
\begin{enumerate}
\item 写出 Hecke 修改的定义，并说明为什么带固定平凡化的修改后丛的参数空间是 $\mathrm{GL}(N,\mathcal X)/\mathrm{GL}(N,\mathcal O)$。
\item 权 ${}^{L}w=(m_1,\dots,m_N)$ 的修改为什么恰好对应同态 $\rho:\mathrm{U}(1)\to\mathrm{U}(N)$？这个 $\rho$ 在 't~Hooft 一侧是什么数据？
\item $\mathrm{U}(1)$ 情形用 $c_1$ 的积分证明 $L_+\cong L_-\otimes\mathcal O(p_0)^m$；指出哪一步是纯拓扑论证，并说明为什么非阿贝尔情形不能照搬。
\item 为什么 $\sum m_i=0$ 对应 't~Hooft 算子的规范群是 $G$ 而不是 $\mathrm{U}(N)$？
\item 说明 Bogomolny 方程的两种商（辛商与复商）各是什么，为什么说「修改空间 $=$ 解空间」只在复结构意义上成立。
\item 什么是单极子冒泡？它对应修改空间紧化的哪一步？Hecke 特征层的定义与几何 Langlands 的标准陈述各是什么？
\end{enumerate}
\end{checkbox}

% ======================================================================
%  07_dmod.tex —— 第 7 讲 翻译器：从 A-膜到 D-模
%  依据：Kapustin--Witten 2006 \S11（印刷 pp.184--206）
% ======================================================================

\section{翻译器：从 A-膜到 D-模}
\label{sec:dmod}

\begin{goals}
\begin{itemize}
\item 说清几何 Langlands 为什么需要一台「从 A-膜到 D-模」的翻译器，以及它必须满足什么条件。
\item 掌握典范余各向同性 A-膜的定义与唯一性，并理解为什么它的开弦代数恰好是微分算子代数。
\item 记住 $f(\Psi)=\Psi$ 的结论、来源与不确定度，知道「普通 D-模」落在 $\Psi=h$ 而不是别处。
\end{itemize}
\end{goals}

几何 Langlands 的标准陈述是：曲线 $C$ 上每个局部系统 $\rho:\pi_1(C)\to\LdualC$，对应 $\MC$ 上一个既被 Hecke 算子作用、又是其本征对象的层——一个 Hecke 特征层，而 Hecke 特征层首先是一个 D-模。前几讲已经把左边翻译成了物理对象：局部系统是 magnetic eigenbrane，Hecke 算子是 't~Hooft 线算子（\cite{kw06}；\src{B p.184}）。但右边在物理侧完全缺席：我们手里只有 $\MH$ 上的膜，而 D-模住在另一个模空间上。本讲要补的就是这台翻译器，原文 \S11 把它称为「缺的一环」（\src{B p.184}）。

\subsection{缺口的形状：为什么翻译器必须「空间充满」}

要理解 \S11 的策略，先把它想成一句范畴论常识：如果手头有一张膜 $B$，那么对任意另一张膜 $B'$，$(B,B')$ 弦构成 $(B,B)$ 弦代数上的一个模；而 $(B,B)$ 弦就是这个膜作为范畴对象的自同态代数（\src{B p.196}）。于是只要把某一张膜的\textbf{自同态代数}算清楚，所有膜就自动成为这个代数的模——翻译器不是一台把膜映成 D-模的机器，而是一张\textbf{万有探针}，它的开弦代数就是答案本身。

剩下的问题是：哪张膜能当探针？普通的 A-膜支撑在 Lagrange 子流形上，维数只有靶空间的一半，它撑不满空间，也就管不到所有膜；而「所有膜」里必须包括那些支撑在别处的膜。所以探针必须满足两条：第一，它是 A-膜；第二，它撑满整个 $\MH$。前者保证它属于我们正在研究的 A-模型，后者保证它的自同态代数足够大，大到能容纳所有膜。

这里要立刻打一剂预防针。原文的承诺是「我们将会说明」每个复结构 $K$ 中的 A-膜都自动给出 $\MC$ 上的一个 D-模，措辞是 will show 加 roughly speaking（\src{B p.184}）；「范畴等价」只是转引更早的一般结论（\src{B p.184}），本情形是该结论的一个特例，原文并没有在本节证明任何范畴层面的等价。本讲把话说准：我们讲的是\textbf{单向构造}，不是函子性定理。

\subsection{典范余各向同性膜：曲率由几何唯一决定}

\term{coisotropic submanifold}{余各向同性子流形}是辛流形 $X$ 中局部由一族 Poisson 交换的函数 $f_1,\dots,f_d$ 的零点定义的子流形；这类函数最多有 $\tfrac12\dim X$ 个独立，所以 $\dim Y\ge\tfrac12\dim X$，取等号时 $Y$ 才是 Lagrange 的（\src{B p.185}）。\S11 只需要最简单的极端情形：$Y=X$、rank one，即膜是空间充满的，Chan--Paton 丛是一条 $U(1)$ 线丛 $L$，规范场局部是实一形式 $A$，曲率 $F=\dd A$。这样的膜要是 A-膜，曲率必须满足
\begin{equation}
(\omega^{-1}F)^2=-1 .
\label{eq:cc-condition}
\end{equation}
这条方程的几何含义是：$\hat N=\omega^{-1}F$ 是一个近复结构（\src{B p.185}）。由于 $F$ 与 $\omega$ 都闭，(11.1) 还保证 $\hat N$ 可积；而 $\omega$ 与 $F$ 对 $\hat N$ 都是 $(2,0)\oplus(0,2)$ 型，故 $\omega-\ii F$ 是 $\hat N$ 下的全纯辛形式。原文由此得出一个限制：这类膜只在 $\omega$ 是全纯辛形式的实部时才可能存在（\src{B p.185}）。

在 $\MH$ 上这件事恰好成立。取 $\omega=(\Im\tau)\omega_K$，则 $\hat N=\omega_K^{-1}\omega_J=I$，而 $\omega_J$ 恰当，故存在势 $\alpha$ 使 $F=(\Im\tau)\omega_J=\dd\alpha$，它是平凡线丛上的一个固定联络（\src{B p.186}）。一般解是 $F=(\Im\tau)(a\omega_J+b\omega_I)$、$a^2+b^2=1$；但当 $b\neq0$ 时该形式不恰当、周期不是 $2\pi$ 的整数倍，根本定义不了膜，所以可用的只剩 $b=0$ 与整体变号。这就是「典范」二字的含义：$F$ 不是选出来的，而是由 $\MH$ 的全纯辛几何唯一决定的；$G$ 非单连通时还可以扭一个平坦线丛，而平坦扭不改变开弦代数层（\src{B p.186}；\src{B p.193}）。

再看它的类型。$\hat N=I$ 说明它在复结构 $I$ 中是 A-膜；在 $K$ 中 $(A,B,A)$ 的判定用 $(\omega_I^{-1}\omega_J)^2=-1$ 同样成立；而在 $J$ 中 $F\propto\omega_J$ 是 $(1,1)$ 型，故它是 B-膜。三句话合起来：c.c.\ 膜是 $(A,B,A)$ 型（\src{B p.187}）。另外，它没有世界面瞬子修正：边界条件与 $K$ 中全纯性联立给出 $(\omega-\ii F)\partial_0\Phi=0$，$\omega-\ii F$ 非退化迫使 $\Phi$ 在边界上为常数，再由解析性整体为常数（\src{B p.187}）。于是计算开弦代数只需要微扰论。

开弦代数为什么是微分算子？局部看，限制到 $\MH$ 中形如 $T^*V$ 的开集（$V\subset\MC$ 小开集）后，边界态的乘法由 $\omega_J=\Re\Omega_I$ 给出的 Poisson 括号控制：$[\kappa,f]=-\ii(\Im\tau)^{-1}\kappa(f)$，而 $\hat p_\alpha=-\frac{\ii}{\Im\tau}\frac{\partial}{\partial q^\alpha}$（\src{B p.191}）。也就是说，$\varphi_z$ 方向的坐标在量子化后变成\textbf{微分}，而 $A_z$ 方向的坐标变成\textbf{乘算子}——这正是全纯微分算子的局部模型。整体上要粘：转移函数的三重上链若属于 $2\pi\ZZ$ 就得到线丛上的 $\Dmod_L$，否则得到复幂 $\Dmod_{L^\gamma}$（\src{B p.192}）。最后一步是时间反演：$t=1$ 处 $T$ 是精确对称，它把开弦代数送到 opposite 代数，而 $\Dmod_L$ 的 opposite 是 $\Dmod_{L^{-1}\otimes K_M}$；自对偶要求 $L^2=K_M$，即 $L=K_M^{1/2}$（$\MC$ 是 spin 流形，平坦扭不改层）（\src{B pp.193--195}）。这就是 c.c.\ 膜的开弦代数层 $\Dmod_{K_{\MC}^{1/2}}$，也就是几何 Langlands 里那个「所有 D-模的老家」。

\subsection{翻译如何发生：开弦代数、模与磁特征膜}

把上一小节压成一句可用的命题：任意 $B'$ 在 $K$ 中是 A-膜，则 $(B_{c.c.},B')$ 弦层构成 $\Dmod_{K_{\MC}^{1/2}}$ 的一个模——这就叫「$B'$ 给出了一个 D-模」（\src{B p.196}）。原文对此用 we claim、we argue，态度是诚实的：这是物理上自然的构造，不是数学证明。

两个例子能让人相信它。第一个例子是校准用的：取 $B'$ 为 $\varphi=0$ 处的平坦膜，它的支撑是 $\MC$ 在 $\MH$ 中的拷贝，属于 $(B,A,A)$ 型。法向费米子没有零模，切向零模量子化给出 $\MC$ 上的旋量，即 $K_{\MC}^{1/2}\otimes(\bigoplus_j\Omega^{0,j})$；拓扑超荷 $Q$ 化为 $\bar\partial$，于是物理态恰是 $K_{\MC}^{1/2}$ 的全纯截面——这正好是 $\Dmod_{K_{\MC}^{1/2}}$ 的定义模（\src{B pp.197--198}）。第二个例子才是几何 Langlands 要用的：type F 膜，即卷在 Hitchin 纤维上、带平坦 Chan--Paton 线丛的膜。它是 $(A,B,A)$ 型，在 $K$ 中是 A-膜，所以按上面的命题自动给出 D-模；而它同时又是 zerobrane 的 S-对偶像，即 magnetic eigenbrane（\src{B p.198}）。

「magnetic eigenbrane」与「D-模」的对应，值得单独说清楚。特征膜的定义是：某条线算子作用在它上面时，只把它乘以一个固定的向量空间（或局部系统），而不是把它送到别的膜。在线算子一侧，这体现为「本征值」；翻译到代数一侧，同一个作用必须实现为该模层的一个自同态——因为 $(B_{c.c.},-)$ 把膜的作用映成模的作用。于是物理上的「本征值」就变成层论上的「Hecke 算子的本征作用」，而 magnetic eigenbrane 经翻译得到的那层，就是 Hecke 特征层。这一步与第 6 讲「't~Hooft 算子 $=$ Hecke 算子」是同一件事的两半：算子对应建立了，作用对象才第一次有意义（\src{B pp.184--185}）。

这里有一个必须承认的技术缝隙：层的粘接只对 $T^*V$ 型开集严格成立，精确理论中的完全层化涉及 Moyal 乘积的收敛问题，原文自己也没有做完（\src{B p.190}）。所以「每个 A-膜给出一个 D-模」在数学上仍是一条待修补的构造，而不是定理。

还要把话说白：原文这一步的口吻是「物理上应该」——c.c.\ 膜、$(B_{c.c.},B')$ 弦、层化，每一步都是物理上自然的构造，但没有一条定理保证它真的造出一个 D-模；数学上严格的对应要等后来纯代数几何的工作，下一讲再说。

\begin{figure}[htbp]
\centering
\begin{tikzpicture}[font=\small,
  bx/.style={draw=kwblue!70,fill=kwblue!4,rounded corners=1.5pt,align=center,
             inner sep=4pt,text width=3.15cm,font=\footnotesize},
  arr/.style={-{Stealth[length=2mm]},thick}]
\node[bx,draw=kwred!70!black,fill=kwred!5] (cc) at (0,2.3) {c.c.\ 膜 $B_{c.c.}$\\space-filling，$(A,B,A)$};
\node[bx] (alg) at (5.4,2.3) {$(B_{c.c.},B_{c.c.})$ 弦\\$=\Dmod_{K_{\MC}^{1/2}}$};
\node[bx] (a) at (0,0) {任意 A-膜 $B'$\\（复结构 $K$ 中）};
\node[bx] (mod) at (5.4,0) {$(B_{c.c.},B')$ 弦\\$=$ 上述代数上的模};
\node[bx,draw=kwgreen!65!black,fill=kwgreen!5] (dm) at (10.9,0) {$\MC$ 上 D-模\\$\Psi=h$ 时为普通 D-模};
\draw[arr] (cc) -- (alg) node[midway,above,font=\scriptsize]{自同态};
\draw[arr] (a) -- (mod) node[midway,above,font=\scriptsize]{张成模};
\draw[arr] (alg) -- (mod);
\draw[arr] (mod) -- (dm) node[midway,above,font=\scriptsize]{层化};
\draw[arr,dashed,draw=kwred!75!black] (dm.south) to[bend left=22]
  node[below,font=\scriptsize,text=kwred!75!black]{$\Psi\neq0$：把 $K_M^{1/2}$ 换成 $K_M^{1/2}\otimes L^{f(\Psi)}$} (mod.south);
\end{tikzpicture}
\caption{翻译器的工作方式。c.c.\ 膜的自同态代数给出微分算子层；任意 A-膜通过 $(B_{c.c.},B')$ 弦自动成为该代数的模，于是得到一个 D-模。$\Psi\neq0$ 时线丛指数要按 $f(\Psi)$ 调整（\src{B p.196}；\src{B p.207}）。}
\label{fig:translator}
\end{figure}

\subsection{量子化参数 $f(\Psi)$：普通 D-模与扭化 D-模}

为什么不能只在 $\Psi=0$ 处做翻译？因为几何 Langlands 对偶本身是交换 $\Psi=0$ 与 $\Psi=\infty$ 的 S 变换，$S$ 作用在参数上是 $(G,\Psi)\mapsto(\Ldual G,-\Psi/n_g)$，$T$ 作用是 $\Psi\mapsto\Psi+1$；全 S-对偶群把所有有理 $\Psi$ 串起来（\src{B p.201}）。若翻译器只在一个点上工作，就无法陈述对偶本身。于是 \S11.3 必须把 c.c.\ 膜推广到任意 $\Psi$。

一上 $\Psi\neq0$ 就遇到新障碍：A-模型的 B-场变成 $B=-\omega_I\Re\tau$，它沿 Hitchin 纤维不为零也不平坦，而带 B-场的 A-膜条件 $F+B|_Y=0$ 要求 $[B|_Y/2\pi]$ 是整上同调类；$\theta\neq2\pi n$ 时 Hitchin 纤维上根本没有合法的 rank-one A-膜（\src{B pp.202--203}）。原文的办法是改造探针：空间充满膜要解
\begin{equation}
\bigl(\omega^{-1}(F+B)\bigr)^2=-1,\qquad \omega=(\Im\tau)\omega_K ,
\label{eq:cc-psi}
\end{equation}
其解为 $F=\Im\tau\cos q\,\omega_J$、$\sin q=-\Re\tau/\Im\tau$（\src{B p.204}）。此时 $\MH$ 上的全纯结构不再是 $I$，而是 $I(\Psi)=I\cos q-J\sin q$；纤维方向的结构也不再是普通联络，而是 $\lambda$-连接，$\lambda=\ii\tan(q/2)$（\src{B p.204}）。换句话说，$\Psi$ 不只改变一个系数，它改变了「哪个复结构是拓扑的」。

代数结构随之改变：c.c.\ 弦层从 $\Dmod_{K_M^{1/2}}$ 变成作用在 $K_M^{1/2}\otimes L^{\Psi}$ 上的微分算子层（\src{B p.203}）。原文接着用整数点把它钉死：$\Psi\to\Psi+n$ 使 $B/2\pi$ 平移 $c_1(L^{-n})$，于是在 $\Psi=n$ 处可以取支撑在 $\varphi=0$、Chan--Paton 丛为 $L^n$ 的膜 $B_n$，量子化 $(B_{c.c.}^n,B_n)$ 弦得到的模恰是 $\Dmod_{K^{1/2}\otimes L^n}$ 的模（\src{B p.205}）。从这里出发，原文先假设 $f$ 线性，由整数点 $f(n)=n$ 得 $f(\Psi)=\Psi$；再用欧氏作用量证明线性：把它写成「$T^*\MC$ 上恰当项 $+$ 从 $\MC$ 拉回的闭二形式」，后者的系数正比于 $\Re\tau$，而粘接修正只需算到 $1/(\Im\tau)^2$ 阶，在这一阶上严格线性（\src{B pp.206--207}）。最后一个常数来自 $\MC$ 的典范线丛同构于 $L^{-2h}$（$h$ 为 $G$ 的对偶 Coxeter 数），于是任意 $\Psi$ 处的代数层是 $L^{\Psi-h}$ 上的全纯微分算子层；取 $\Psi=h$ 就得到\textbf{普通}微分算子层，即作用在函数上的 D-模（\src{B p.207}）。

所以本讲给出的具体式子是
\begin{equation}
f(\Psi)=\Psi ,
\label{eq:f-Psi}
\end{equation}
来源是原文 \S11.3 的两步论证（整数点定值 $+$ 微扰线性），把典范丛的常数 $\!-h$ 吸收进去之后，普通 D-模落在 $\Psi=h$。这里的不确定度要如实标注：$f(\Psi)=\Psi$ 的推导是微扰的，只算到 $1/(\Im\tau)^2$ 阶，线性性依赖「$\Psi$ 是唯一参数」这一物理假设，不是非微扰构造；整数 $\Psi$ 处的值是硬锚点，非整数 $\Psi$ 处的结论则软得多（\src{B p.205}；\src{B p.207}）。

最后纠一个常见的记混。$\Psi=\ii$ 确实是一个特殊值，但那是 \S5 镜像对称的特殊值：在 $\Psi=\ii$ 上，对 $J$ 的 B-模型等价于对 $K$ 的 A-模型（\src{B p.71}）。它管的是「哪个模型描述同一个物理」，与 D-模的「通常／扭化」不是一回事。在 \S11.3 的归一化下，普通 D-模出现在 $\Psi=h$；若换用某些 D-模文献把普通 D-模记为扭化参数为零的约定，则此处的扭化参数应读作 $\Psi-h$。把这个约定记住，就不容易在 $f(\Psi)$ 上翻车。

\begin{keybox}
c.c.\ 膜之所以能当翻译器，是因为它空间充满、且曲率由 $\omega_J$ 典范地定死，于是它的开弦代数恰是微分算子层 $\Dmod_{K_M^{1/2}}$，任何 $K$ 中的 A-膜自动成为它的模；把 $\Psi$ 从 $0$ 调开，只是把这个层扭成 $L^{\Psi-h}$，$f(\Psi)=\Psi$，在 $\Psi=h$ 处回到普通 D-模。
\end{keybox}

\begin{follow}
\begin{itemize}
\item 主线是 \S11.1--\S11.3（\src{B pp.185--207}）。先读 \S11.1 开头的余各向同性膜定义与 (11.1)，确认 $F$ 为什么由 $\omega_J$ 定死、$(a,b)$ 族里为什么只剩 $b=0$（\src{B pp.185--187}）。
\item 再读 \S11.2 开头两段与 $U(1)$ 例子（\src{B pp.196--201}），重点只有一句：为什么任意 A-膜都是 c.c.\ 代数的模。
\item \S11.3 可跳读：只抓 (11.30)--(11.34) 与最后三页（\src{B pp.204--207}），$\lambda$-连接的细节不影响对 $f(\Psi)$ 的理解。
\item 时间紧就先背 $\Psi$ 轴上的三个点：$\Psi=0$ 给 $\Dmod_{K^{1/2}}$，$\Psi=n$ 给 $\Dmod_{K^{1/2}\otimes L^n}$，$\Psi=h$ 给普通 D-模（\src{B pp.205--207}）。
\end{itemize}
\end{follow}

\begin{checkbox}
\begin{itemize}
\item 为什么普通 Lagrangian A-膜不能当翻译器？「空间充满」这一条在机制里具体用在哪一步？
\item 在 $\MH$ 上，$(11.1)$ 的一般解有一整族，为什么只有 $b=0$ 那一支能定义膜？「典范」的「典」在哪里？
\item 从开弦代数走到 $\Dmod_{K_{\MC}^{1/2}}$，哪一步用到时间反演？它把线丛指数 $L$ 钉成了什么？
\item $f(\Psi)=\Psi$ 的论证用了哪两个输入？哪一个是硬锚点，哪一个是微扰假设？不确定度在哪一页？
\item $\Psi=h$ 与 $\Psi=\ii$ 各自特殊在哪里？为什么不能把 $\Psi=\ii$ 说成「普通 D-模所在的位置」？
\end{itemize}
\end{checkbox}

% ======================================================================
%  08_outlook.tex —— 第 8 讲 收获与边界：证明了什么，没证明什么
%  依据：Kapustin--Witten 2006 \S12（印刷 pp.207--221）与全文边界
% ======================================================================

\section{收获与边界：证明了什么，没证明什么}
\label{sec:outlook}

\begin{goals}
\begin{itemize}
\item 看懂 \S12 的检查表：三类膜 $(B,B,B)$、$(B,A,A)$、$(A,B,A)$ 如何在四维规范理论里被边界条件造出来，为什么每一类都保两个拓扑超荷。
\item 能把这份推导的收获与缺口分开说：哪些是算出来的，哪些是「物理上应该」，哪些根本没做。
\item 知道 2024 年前后纯代数几何的证明与这条物理路线是什么关系，以及接下来该读什么。
\end{itemize}
\end{goals}

前十讲我们把一条链从四维拉到了二维：$\mathcal N=4$ 超对称 Yang--Mills 扭化、紧化到 $C$ 上、得到以 $\MH$ 为靶的 sigma 模型，再把 Wilson/'t~Hooft 线算子翻译成膜上的算子，把 't~Hooft 算子认成 Hecke 算子，最后用 c.c.\ 膜把 A-膜翻成 D-模（\cite{kw06}）。这一讲做三件事：把膜放回四维（\S12），清点这份推导的边界，再给出继续读下去的路线。

\subsection{回到四维：\S12 是一张检查表}

为什么最后还要回四维？因为几何 Langlands 的 Hecke 算子定义在 $\mathrm{Bun}_G(C)$ 上——所有 $G$-丛构成的栈，而不只是稳定丛的模空间 $\MC$；二维 sigma 模型的靶只有 $\MH$ 的光滑部分，在那里定义不了整体对象（\src{B p.184}）。原文把困难转成规范理论的类比：栈的规范理论对应物是所有联络的空间，于是「膜在二维定义不了的地方」，可以回到四维用边界条件定义。\S12 的目标原话就是：说明重要的膜确实可以在四维规范理论中定义（\src{B p.207}）。

四维的设定是 $M=\RR^2_+\times C$，边界在 $x^1=0$；一张膜就是边界上的一个超对称边界条件。做法是固定的四步：选一个保超对称的对合 $T$，猜费米子边界条件 $T\lambda|=\pm\lambda|$，用超对称变分 $\delta\lambda=\tfrac12\Gamma^{IJ}F_{IJ}\epsilon$ 推出玻色子边界条件，最后验证超流法向分量消失、并检查存在超对称态（\src{B pp.208--210}）。每一步都是线性条件，可以逐条核对，所以 \S12 读起来像清单而不像新理论——它不产出新结论，只验证前面用到的膜确有四维来源。

三类膜是清单的主角。$(B,B,B)$ 型是 $\MH(\Ldual G,C)$ 上的 zerobrane：一个点，在三个复结构里都是 B-膜。它的镜像是 Hitchin 纤维——对 $I$ 复、对 $J$ 与 $K$ Lagrange，故是 $(B,A,A)$ 型。c.c.\ 膜是 $(A,B,A)$ 型，它的镜像也是（\src{B p.207}）。这里的 $I,J,K$ 就是第 1 讲（原文 §4.1）介绍的三个复结构：一张膜在某个复结构中算 A 还是 B，看它保住的是那个结构的 A-型还是 B-型拓扑超荷；同一张膜相对三个结构可以有不同的类型，这正是三元组记号的读法。

为什么每一类都恰好保两个拓扑超荷？扭化后的理论里有两个标量超荷 $Q_l,Q_r$（它们的任意非零线性组合 $Q=uQ_l+vQ_r$ 都可以充当拓扑超荷），而边界条件一旦保住 $Q$，也自动保住它的共轭 $Q^\dagger$；两个拓扑超荷加上各自的共轭，共四个超对称——是 $\mathcal N=4$ 的 16 个中的四分之一，所以这类膜叫 $1/4$ BPS 膜；从只有一半超对称的二维 sigma 模型看，它们反而是 $1/2$ BPS（\src{B p.208}）。三种类型之所以都保同样两个超荷，是因为它们在不同复结构中的 A/B 类型，最终都化为 $Q_l,Q_r$ 的线性组合。

三类膜的四维实现可以各记一个标志性条件。$(B,B,B)$：取 $T=\Gamma^{04}$，玻色条件是 $A_0+A_4=0$、其余 $A_I$ 取 Dirichlet，而超对称态条件化成 Hitchin 方程，于是边界值给出 $\MH$ 上的一个点（\src{B pp.211--213}）。$(B,A,A)$：对 $(B,B,B)$ 的条件作 S-对偶，$T$ 变成 $\Gamma^1\Gamma^2\Gamma^3\Gamma^4$，剩下的条件等价于复 Hitchin 方程，膜就卷在 Hitchin 纤维上（\src{B p.216}）。$(A,B,A)$：混合边界条件一部分 Dirichlet、一部分 Neumann，把降维后的条件与 sigma 模型的混合边界条件对比，可读出目标空间规范场恰是 $f=\omega_J$——正是 \S11 里 c.c.\ 膜的曲率（\src{B p.220}）。这一步是全讲的闭环：\S11 用 c.c.\ 膜造出 D-模，\S12 又从四维边界条件把同一张膜造回来。

\begin{table}[htbp]
\centering
\footnotesize
\begin{tabular}{@{}l p{0.19\textwidth} c p{0.23\textwidth} p{0.20\textwidth}@{}}
\toprule
膜类型 & 在 $\MH$ 上的几何 & $(I,J,K)$ & 保哪两个超荷 & 四维实现要点 \\
\midrule
$(B,B,B)$ & zerobrane：支撑在一个点 & $B,B,B$ &
两个拓扑超荷 $Q_l,Q_r$；三个结构中的 B-型 BRST 都是它们的线性组合 &
$T=\Gamma^{04}$；$A_0+A_4=0$，其余 Dirichlet \\
$(B,A,A)$ & Hitchin 纤维：固定 $\varphi_z$ 的特征多项式 & $B,A,A$ &
同样是 $Q_l,Q_r$；$I$ 中取 B-型、$J,K$ 中取 A-型的组合 &
$T'=\Gamma^{1}\Gamma^{2}\Gamma^{3}\Gamma^{4}$；复 Hitchin 方程 \\
$(A,B,A)$ & space-filling：全 $\MH$，带曲率 $F\propto\omega_J$ & $A,B,A$ &
同样是 $Q_l,Q_r$；$I,K$ 中取 A-型、$J$ 中取 B-型的组合 &
$\hat T$ 的混合边界条件；$f=\omega_J$ 认出 c.c.\ 膜 \\
\bottomrule
\end{tabular}
\caption{三类关键膜的四维对照。三元组记录的是同一张膜相对三个复结构 $I,J,K$ 的 A/B 类型，不是三种不同的膜理论；三者都保两个拓扑超荷（连共轭共四个超对称），所以都是 $1/4$ BPS（\src{B p.207}；\src{B p.208}）。}
\label{tab:three-branes}
\end{table}

清单也有它自己的诚实之处：提升不唯一。原文明确说，某些情形会给出多个不同的四维边界条件，它们在降到二维后、至少在奇点之外看起来等价；$(B,B,B)$ 的第二种构造在平坦时空中保八个超荷而不是四个，两种构造只在奇点附近才可能显出差别（\src{B pp.213--214}）。因此 \S12 检查的是「存在性」，不是「唯一性」。

\subsection{这份推导给了什么：一个图像、一套对应、大量证据}

把整条链压成一句话：\kw{} 用四维 $\mathcal N=4$ 理论的 S-对偶实现了几何 Langlands 的两边——左边是局部系统（magnetic eigenbrane），右边是 D-模（经 c.c.\ 膜翻译），中间的桥是 S-对偶在二维变成镜像对称（\cite{kw06}）。它给出的收获分三层。第一层是\textbf{物理图像}：几何 Langlands 不是孤立猜想，而是电磁对偶在特定紧化下的表现。第二层是\textbf{对应关系}：对象层（eigenbrane $\leftrightarrow$ Hecke 特征层）与算子层（'t~Hooft $\leftrightarrow$ Hecke）各有一套明确的字典。第三层是\textbf{证据}：$U(1)$ 情形可以完整算完，三类膜的类型表自洽，\S12 用四维边界条件独立造出 \S11 用到的膜，边界算子的 Q-上同调又与开弦代数对上。

但它是推导，不是证明。原文用 will show、roughly speaking、we claim、we argue 这些词，\S11 的翻译是单向构造，「范畴等价」只是转引（\src{B p.184}）。三处最明显的软肋是：其一，翻译与 Hecke 作用的\textbf{交换性}没有逐条证明——若「先翻译再作用」与「先作用再翻译」不一致，得到的 D-模未必是 Hecke 特征层；其二，层的完全粘接只在 $T^*V$ 型开集上成立，精确理论里有收敛问题（\src{B p.190}）；其三，非整数 $\Psi$ 处的 $f(\Psi)$ 只由微扰论证确定（\src{B p.207}）。

所以我们建议这样记：\kw{} 给的是\textbf{为什么}，不是\textbf{是什么的证明}。物理推导的可信度来自独立路径给同一答案，而不是来自某一步的形式化；读它时要能指出每一步的证据强度，这正是 \S12 存在的意义——它是这份推导的诚实度指标。

\subsection{边界在哪里：非分歧、二维 CFT、以及 2024 年的证明}

原文自己标出了两条边界。第一条：只做非分歧（unramified）情形，即只考虑 $C$ 上没有分歧点的局部系统，分歧与野分歧完全没碰（\src{B p.5}）；分歧情形要在 $\mathrm{Bun}_{G,y}$ 上加 level structure、引入仿射 Hecke 代数，是 Gukov--Witten 的主题（\cite{gw06}）。第二条：没有说清二维共形场论在几何 Langlands 里到底起什么作用（\src{B p.5}）。这不是小遗漏：Frenkel 那一派恰恰以二维 CFT 为主工具——顶点代数、共形块、临界 level $k=-h^\vee$ 的仿射代数——从另一头构造 Hecke 特征层（\cite{fr09}）。两条路线对象相同、入口相反：\kw{} 解释这条对应为什么存在，Frenkel 的路线构造它是什么。

还有两个技术边界值得记住：D-模本身必须做在 $\mathrm{Bun}_G(C)$ 这个栈上，只在稳定丛上做不闭合（\src{B p.185}）；某些陈述对 $C$ 的 spin 结构敏感，$\Psi=h$ 的普通 D-模陈述也要据此修正（\src{B p.207}）。

那么数学上到底证明了没有？2006 年时没有。2024 年前后，Gaitsgory 与合作者用五篇系列预印本给出了特征 $0$ 代数闭域上光滑射影曲线的几何 Langlands 的证明（D-模、非分歧情形）。要把这件事的位置说准：它是\textbf{数学定理}，走的是纯代数几何与范畴化的路线，既不用 S-对偶，也不用膜；它不验证 \kw{} 的物理论证，\kw{} 也没有预言它的方法。两者的关系是分工：物理路线提供直觉、猜想与「为什么应该对」，数学路线提供定理与「确实对」。有意思的是，定理的适用范围（非分歧、光滑射影曲线）与 \S12 检查表的适用范围恰好一致——这不是巧合，而是 \kw{} 一开始就把自己限制在能做干净的那一块。

\subsection{接下来读什么：文献与三条建议}

如果只再读一份，读 Witten 2008 的十五页短文（\cite{w08}）：它把「紧化 $+$ 镜像对称」这一段重写成 Hitchin 纤维化上的镜像对称，是整条链最短的物理重述。想补二维 CFT 那一侧，读 Frenkel 的短综述（\cite{fr09}），它会告诉你 Hecke 特征层在 CFT 里是怎么被造出来的。想要更全的对照，可用 Ben-Zvi 的综述（\cite{bz04}）与 Hitchin 系统的材料（\cite{ht03}）。要往分歧情形走，读 Gukov--Witten（\cite{gw06}）；想把框架推到三维与面算子范畴，读 Setter 的博士论文（\cite{set13}）——注意它本身既不证明几何 Langlands，也不定义 Hecke 作用，读它是学技术而不是找结论。

给读者三条建议。第一，\textbf{先读第 5、6 讲再回原文 \S11}：c.c.\ 膜如果没有 eigenbrane 与 Hecke 算子的动机，看起来只是一张奇怪的膜；有了动机，才知道它为什么必须空间充满。第二，\textbf{\S12 可以当检查表跳读}：只需记住三类膜、各自的一句话来源，以及「提升不唯一」这一条；四步法的细节等要用时再回查。第三，\textbf{卡住时用 $U(1)$ 例子校准}：\S11.2 末的 $U(1)$ 计算是全文唯一能从头算到尾的例子，把它走一遍，$\lambda$-连接、平坦联络、Poincaré 对偶这些词就都有了落点（\src{B pp.199--201}）；再卡，就查解经本对应页段的骨血层。

\begin{keybox}
\kw{} 交给几何 Langlands 的是一份物理推导：一个 S-对偶图像、一套膜与算子的字典、以及 \S11 与 \S12 两方向互相对上的大量证据；它没给证明，只做非分歧情形，也不解释二维 CFT 的作用——真正的定理要等 2024 年前后纯代数几何的路线。
\end{keybox}

\begin{follow}
\begin{itemize}
\item 先读 \S12 的开头三页（\src{B pp.207--210}）：目标、三类膜、为什么保两个拓扑超荷、四步法，全在这里。
\item 再挑一类细读：$(B,B,B)$ 在 \src{B pp.211--215}，$(A,B,A)$ 与 c.c.\ 膜的识别在 \src{B pp.217--221}；后者与 \S11 的接口最值得看。
\item 边界部分读原文 \S1 的结尾（\src{B p.5}）：作者自陈的两个 gap（非分歧、二维 CFT），是理解这份推导定位的关键。
\item 后续文献按顺序：\cite{w08} 十五页 $>$ \cite{fr09} 三十二页 $>$ \cite{bz04} 综述 $>$ \cite{gw06} 分歧 $>$ \cite{set13} 博士论文。
\end{itemize}
\end{follow}

\begin{checkbox}
\begin{itemize}
\item 三类膜各对应 $(I,J,K)$ 中的哪种 A/B 类型？为什么同一张膜可以在三个复结构里有不同身份？
\item 「保两个拓扑超荷」为什么等于「共四个超对称」？$1/4$ BPS 与 $1/2$ BPS 两个说法分别从哪个理论的角度看？
\item \S12 为什么叫检查表而不是新理论？举一个它检查出「提升不唯一」的例子。
\item 原文自陈的两个 gap 是什么？Frenkel 的路线如何从另一头绕开第二个 gap？
\item 2024 年的证明与 \kw{} 的物理路线是什么关系？为什么不能说「物理证明了几何 Langlands」？
\end{itemize}
\end{checkbox}


\appendix
\appendixstyle
\section{术语与符号表}

本表按主题分七组列出正文用到的术语。每条先给英文与中文，再给一句话解释；标「待核」者表示笔者对经典文献中的精确表述尚无十足把握，读者引用前请核对原文。符号的基本约定见第 1 讲，此处不重复。

\begingroup
\small
\setlength{\tabcolsep}{4pt}
\begin{longtable}{@{}>{\raggedright\arraybackslash}p{0.285\textwidth}>{\raggedright\arraybackslash}p{0.66\textwidth}@{}}
\toprule
\textbf{英文（中文）} & \textbf{一句话解释} \\
\midrule
\endfirsthead
\toprule
\textbf{英文（中文）} & \textbf{一句话解释} \\
\midrule
\endhead
\bottomrule
\endlastfoot

\multicolumn{2}{@{}l}{\textbf{一、对偶、群论与耦合常数}}\\
\addlinespace[2pt]
\term{electric--magnetic duality}{电磁对偶} & 把电荷与磁荷、局域算子与 disorder 算子互换的对称性，同时把耦合常数取倒数。\\
\term{Montonen--Olive duality}{Montonen--Olive 对偶} & 1977 年提出的猜想：$N=4$ 超对称规范理论在强弱耦合互换下自对偶。\\
\term{S-duality}{S-对偶} & $\tau\mapsto-1/(n_g\tau)$ 的强弱对偶变换，交换 Wilson 算子与 't Hooft 算子。\\
\term{T-duality}{T-对偶} & 环面上半径 $R\mapsto 1/R$ 的等价性；在 Hitchin 纤维上表现为对偶环面。\\
\term{GNO duality}{GNO 对偶} & Goddard--Nuyts--Olive 的结论：电、磁荷取值格的对偶确定对偶群。\\
\term{Langlands dual group}{Langlands 对偶群 $\Ldual G$} & 把根系与余根系互换所得的复群，例如 $SU(N)$ 与 $PSU(N)$ 互换。\\
\term{lacing number}{织数 $n_g$} & 描述 Dynkin 图边长不对称的整数：单边型为 $1$，$B_n$、$C_n$、$F_4$ 为 $2$，$G_2$ 为 $3$。\\
\term{Hecke group}{Hecke 群} & $G_2$、$F_4$ 时 $S$ 生成的 $SL(2,\RR)$ 的无限离散子群，代替 $SL(2,\ZZ)$。\\
\term{duality group}{对偶群} & 由 $S$ 与 $T$ 生成的离散群，单边型时为 $SL(2,\ZZ)$。\\
\term{complexified coupling}{复耦合 $\tau$} & $\tau=\theta/2\pi+4\pi\ii/e^2$，实部与虚部分别编码 $\theta$ 与 $1/e^2$。\\
\term{theta angle}{$\theta$ 角} & 瞬子数项的系数除以 $2\pi$；$\theta\mapsto\theta+2\pi$ 是规范等价，故 $\Re\tau$ 周期为 $1$。\\
\term{instanton number}{瞬子数} & $(1/8\pi^2)\int\Tr F\wedge F$ 的整数值，标记规范场的拓扑类。\\
\term{weight lattice}{权格 $P$} & 表示权所在的格；$SU(2)$ 的权格为 $(1/2)\ZZ$，$SO(3)$ 的为 $\ZZ$。\\
\term{coweight lattice}{余权格 $P^\vee$} & 磁荷所在的格，与权格互为对偶，GNO 对偶据此定义。\\

\addlinespace[3pt]
\multicolumn{2}{@{}l}{\textbf{二、超对称与扭化}}\\
\addlinespace[2pt]
\term{$\mathcal N=4$ super Yang--Mills}{$N=4$ 超对称 Yang--Mills} & 四维最大超对称规范理论，含 $16$ 个超荷，共形且具有 S-对偶。\\
\term{supercharge}{超荷} & 超对称变换的生成元，把玻色子变为费米子、费米子变为玻色子。\\
\term{topological supercharge}{拓扑超荷} & 扭化后成为 BRST 算子的标量超荷，满足 $Q^2=0$。\\
\term{twisting}{扭化} & 把 R-对称与时空旋转群对角化，从超对称理论得到拓扑场论的操作。\\
\term{GL twist}{GL 扭化} & Kapustin--Witten 所用的特定扭化，因与几何 Langlands 纲领相关而得名。\\
\term{canonical parameter}{典范参数 $\Psi$} & 扭化参数 $t$ 与 $\tau$ 的组合 $\Psi=(\tau t^2+\bar\tau)/(1+t^2)$，是 S-对偶下的自然耦合。\\
\term{twisting parameter}{扭化参数 $t$} & 扭化族的 $\PP^1$ 参数；$t$ 与 $-t$ 等价，$t=\pm\ii$ 对应 $\Psi=\infty$。\\
\term{BRST operator}{BRST 算子} & 扭化理论中幂零的标量超荷 $Q$，物理态由其上同调给出。\\
\term{$Q$-cohomology}{$Q$-上同调} & 以 BRST 微分取上同调，得到拓扑不变的态空间与算子代数。\\
\term{$R$-symmetry}{R-对称} & 在超代数中旋转超荷的对称群；$N=4$ 时为 $SU(4)\cong Spin(6)$。\\

\addlinespace[3pt]
\multicolumn{2}{@{}l}{\textbf{三、Hitchin 模空间与超凯勒几何}}\\
\addlinespace[2pt]
\term{Riemann surface}{Riemann 面 $C$} & 紧复曲线；四维理论沿 $C$ 紧化，$C$ 是二维理论的底空间。\\
\term{genus}{亏格 $g$} & 复曲线的洞数；$\MH$ 的维数与 Hitchin 基的维数都随 $g$ 变化。\\
\term{$G$-bundle}{$G$-丛} & 以群 $G$ 为结构群的纤维丛；其模空间记为 $\mathrm{Bun}_G$。\\
\term{principal bundle}{主丛} & 纤维与 $G$ 同构且带 $G$ 右作用的丛，是 Higgs 丛的底。\\
\term{Higgs bundle}{Higgs 丛 $(E,\varphi)$} & 全纯 $G$-丛 $E$ 配一个全纯的 $\mathrm{ad}$-值 $(1,0)$ 形式 $\varphi$。\\
\term{Higgs field}{Higgs 场 $\varphi$} & $\varphi\in H^0(C,\mathrm{ad}(E)\otimes K_C)$，即 $\mathrm{ad}(E)$ 值全纯 $(1,0)$ 形式。\\
\term{Hitchin equations}{Hitchin 方程} & $F_A+[\varphi,\varphi^*]=0$ 与 $\bar\partial_A\varphi=0$，是 $\MH$ 上的自对偶方程。\\
\term{Hitchin moduli space}{Hitchin 模空间 $\MH$} & Higgs 丛（等价地 Hitchin 方程解）的模空间，是超凯勒流形。\\
\term{stability}{稳定性} & 排除破坏模空间良好性的子丛，Higgs 丛的稳定性使模空间为流形。\\
\term{hyperkähler}{超凯勒} & 带有满足四元数关系的三个可积复结构的度量几何；$\MH$ 是超凯勒流形。\\
\term{hyperkähler quotient}{超凯勒商} & 三个矩映射条件联立后商去群作用，$\MH$ 即无穷维超凯勒商。\\
\term{complex structures $I,J,K$}{复结构 $I,J,K$} & 三个满足 $I^2=J^2=K^2=IJK=-1$ 的可积复结构，膜类型按它们区分。\\
\term{symplectic form}{辛形式} & 闭非退化 $2$-形式；$\omega_J$、$\omega_K$ 恰当而 $\omega_I$ 给出拓扑类。\\
\term{moment map}{矩映射 $\mu$} & 哈密顿群作用的生成函数；辛商取 $\mu^{-1}(0)$ 再商群作用。\\
\term{Hitchin fibration}{Hitchin 纤维化} & 由 Casimir 不变量给出的映射 $\MH\to B$，其纤维是复环面。\\
\term{completely integrable}{完全可积} & Hitchin 系统：存在 $\dim B$ 个泊松对合的函数，其共同水平集为环面。\\
\term{complex torus}{复环面} & 形如 $\CC^n/\Lambda$ 的紧复流形；$\MH$ 的 generic 纤维就是复环面。\\
\term{spectral curve}{谱曲线} & 由 $\det(\lambda-\varphi)=0$ 定义的曲线，用来描述 Hitchin 纤维。\\

\addlinespace[3pt]
\multicolumn{2}{@{}l}{\textbf{四、sigma 模型、广义复几何与膜}}\\
\addlinespace[2pt]
\term{sigma model}{sigma 模型} & 二维场论，基本场是从世界面到靶空间的映射，作用量含靶空间度量。\\
\term{A-model}{A-模型} & 依赖辛结构、只保留 $(1,1)$ 类超荷的拓扑 sigma 模型。\\
\term{B-model}{B-模型} & 依赖复结构、保留全纯超荷的拓扑 sigma 模型。\\
\term{mirror symmetry}{镜像对称} & 两个不同靶空间上 A-模型与 B-模型等价；本文中靶空间为同一 $\MH$ 的不同复结构。\\
\term{generalized complex geometry}{广义复几何} & 统一复结构与辛结构的框架，用 $T\oplus T^*$ 上的结构描述膜的类型。\\
\term{brane}{膜} & sigma 模型中的边界条件；在拓扑场论中也指可被线算子作用的对象。\\
\term{A-brane}{A-膜} & A-模型允许的边界条件：Lagrange 或余各向同性子流形加平坦丛。\\
\term{B-brane}{B-膜} & B-模型允许的边界条件：复子流形加全纯 Chan--Paton 丛。\\
\term{Lagrangian submanifold}{Lagrange 子流形} & 辛形式在其上限制为零且维数为靶空间一半的子流形。\\
\term{coisotropic submanifold}{余各向同性子流形} & 法向方向上辛形式退化（含 Lagrange 情形为特例）的子流形。\\
\term{Chan--Paton bundle}{Chan--Paton 丛} & 膜上的向量丛，编码开弦端点的内部自由度。\\
\term{SYZ conjecture}{SYZ 猜想} & 镜像流形由互为对偶的环面纤维化相联系（Strominger--Yau--Zaslow）。\\

\addlinespace[3pt]
\multicolumn{2}{@{}l}{\textbf{五、线算子与特征膜}}\\
\addlinespace[2pt]
\term{Wilson operator}{Wilson 算子} & 沿闭曲线的和乐之迹 $W_R(S)=\Tr_R P\exp(-\oint A)$，是电型线算子。\\
\term{'t Hooft operator}{'t Hooft 算子} & 先规定沿曲线的奇异性再作路径积分的 disorder 算子，带磁荷标签。\\
\term{line operator}{线算子} & 支撑在一维曲线上的算子；在二维理论中作用在膜上。\\
\term{loop operator}{圈算子} & 支撑在闭曲线上的线算子，其周边界与膜相交。\\
\term{order operator}{有序算子} & 由局域场（如联络）构造的算子，Wilson 算子是其典型代表。\\
\term{disorder operator}{无序算子} & 由奇异性规定而非局域场构造的算子，'t Hooft 算子是其代表。\\
\term{eigenbrane}{特征膜} & 满足 $XB=B\otimes V$ 的膜，即线算子的本征对象；对应几何 Langlands 的特征层。\\
\term{electric eigenbrane}{电特征膜} & Wilson 算子族的联合特征膜，自动带一个平坦向量丛。\\
\term{magnetic eigenbrane}{磁特征膜} & 't Hooft 算子族的联合特征膜，即 Hitchin 纤维上的 type $F$ 膜。\\
\term{zerobrane}{零膜（点膜）} & 支撑在单点上、Chan--Paton 为天穹层的 $(B,B,B)$ 型膜。\\
\term{local system}{局部系统} & 基本群 $\pi_1(C)$ 的表示，等价于局部常值层或平坦联络。\\
\term{fundamental group}{基本群 $\pi_1(C)$} & 曲线上闭路同伦类的群；$\Ldual G$-局部系统即其到 $\Ldual G_\CC$ 的同态。\\
\term{brane type}{膜类型} & 对 $I,J,K$ 三个复结构分别标 $A$ 或 $B$ 所得的三字母标签。\\

\addlinespace[3pt]
\multicolumn{2}{@{}l}{\textbf{六、Hecke 修改、单极子与仿射 Grassmannian}}\\
\addlinespace[2pt]
\term{Hecke operator}{Hecke 算子} & 几何 Langlands 中把丛在指定点处改写后拉回、再推前的对应算子。\\
\term{Hecke modification}{Hecke 修改} & 在点 $p_0$ 外保持同构、在 $p_0$ 处允许极点的丛的重写。\\
\term{Hecke correspondence}{Hecke 对应} & $\mathrm{Bun}\times\mathrm{Bun}$ 中记录修改关系的子簇，由此定义 Hecke 算子。\\
\term{Hecke eigensheaf}{Hecke 特征层} & 在 Hecke 算子下被送到自身张量一个向量空间的层，对应特征膜。\\
\term{Bogomolny equations}{Bogomolny 方程} & 四维自对偶方程降维所得的单极子方程，形如 $F=\ii\star D\varphi$。\\
\term{monopole}{单极子} & 三维空间带磁荷的有限能量解；'t Hooft 算子的奇异性对应它。\\
\term{monopole bubbling}{单极子冒泡} & 单极子奇点处瞬子尺度退化为零、并入单极子的经典现象。\\
\term{affine Grassmannian}{仿射 Grassmannian} & 圈群商 $G_\CC(X)/G_\CC(O)$，参数化在一点处的 Hecke 修改。\\
\term{Schubert cell}{Schubert 胞腔} & 仿射 Grassmannian 中按支配余权的 $G_\CC(O)$-轨道分层所得胞腔。\\
\term{Schubert cycle}{Schubert 循环} & Schubert 胞腔的闭包，出现在 Hecke 对应的纤维中。\\

\addlinespace[3pt]
\multicolumn{2}{@{}l}{\textbf{七、几何 Langlands、$D$-模与量子化}}\\
\addlinespace[2pt]
\term{geometric Langlands program}{几何 Langlands 纲领} & $\Ldual G$-局部系统与 $D$-模（或特征层）之间的对应，是本文的主题。\\
\term{Langlands program}{Langlands 纲领} & 数论与表示论中的巨大对偶网；几何版本用曲线上的丛替换数域。\\
\term{eigensheaf}{特征层} & Hecke 算子的本征对象；KW 在拓扑场论一侧把对应物称为特征膜。\\
\term{$D$-module}{$D$-模} & 微分算子层 $D$ 上的模，几何 Langlands 中对应局部系统的一侧。\\
\term{sheaf of differential operators}{微分算子层 $D$} & 由向量场及其积生成的层；c.c.\ 开弦代数取值于其上。\\
\term{twisted $D$-module}{扭化 $D$-模} & 系数扭化或带中心特征的 $D$-模（待核：不同文献的扭化方式略有差异）。\\
\term{c.c.\ A-brane}{c.c.\ 膜} & 秩 $1$、满足 $(\omega^{-1}F)^2=-1$ 的典范余各向同性 A-膜，型为 $(A,B,A)$。\\
\term{open string BRST}{开弦 BRST} & 边界条件所保留的超荷在边界上成为 BRST 微分，其幂零性给出弦上同调。\\
\term{quantization parameter}{量子化参数 $f(\Psi)$} & c.c.\ 弦代数是作用在 $K_M^{1/2}\otimes L^{f(\Psi)}$ 上的微分算子层，$f(\Psi)=\Psi$。\\
\term{three brane types}{三类膜} & 实际出现的 $(B,B,B)$、$(B,A,A)$、$(A,B,A)$ 三种膜；$(B,B,A)$ 与 $(A,A,A)$ 不出现。\\
\term{ramification}{分歧} & 局部系统在一点处有非平凡单值性；几何 Langlands 向分歧情形推广。\\
\term{unramified}{非分歧} & 局部系统在每点局部平凡，即经典（非分歧）几何 Langlands 的情形。\\
\term{wild ramification}{野分歧} & 局部单值性含幂零部分、需额外数据刻画的情形（待核）。\\
\end{longtable}
\endgroup

\section{习题}

以下三十题按八讲分组，题号连续。每题先给题目，再给一句提示；提示只指出第一步，完整论证请自己补全。标 $\dagger$ 的题需要回到原文相应小节查阅细节。

\subsubsection*{第 1 讲\quad 舞台}

\begin{enumerate}[label=\arabic*.,start=1,leftmargin=2.4em,itemsep=6pt]
\item 设 $C$ 是亏格 $g\ge2$ 的紧 Riemann 面，$G$ 是单连通的单群。证明 Hitchin 模空间 $\MH$ 的复维数为 $(2g-2)\dim G$。
\par\smallskip\noindent\emph{提示：}$\dim\mathrm{Bun}_G=(g-1)\dim G$；$\varphi\in H^0(C,\mathrm{ad}E\otimes K_C)$，对稳定丛用 Riemann--Roch 与 Serre 对偶把 $h^0(\mathrm{ad}E\otimes K_C)$ 算成 $(g-1)\dim G$。

\item 取 $G=U(1)$。证明 $\MH\cong T^*\mathrm{Jac}(C)$，写出 Hitchin 纤维化及其纤维，并说明为什么上一题的维数公式对含 $U(1)$ 因子的群不适用。
\par\smallskip\noindent\emph{提示：}$\mathrm{Bun}_{U(1)}=\mathrm{Jac}(C)$，而 $H^0(C,K_C)\cong\CC^g$；公式失效的原因是 $H^0(C,\mathcal O)\neq0$ 使 $\chi(\mathrm{ad}E)$ 多出中心项。

\item 说明 Casimir 多项式的系数给出 Hitchin 纤维化 $\MH\to B=\bigoplus_iH^0(C,K_C^{d_i})$，并证明 $\dim B=(g-1)\dim G$。
\par\smallskip\noindent\emph{提示：}$\varphi\mapsto(\Tr\varphi^2,\Tr\varphi^3,\dots)$ 在纤维上取常值；再用 $\dim H^0(C,K_C^{d})=(2d-1)(g-1)$ 与 $\sum_i(2d_i-1)=\dim G$。

\item 把 $\MH$ 写成无穷维超凯勒商：写出平坦联络与 Higgs 场空间上的三个矩映射，并验证三个水平条件合起来就是 Hitchin 方程。
\par\smallskip\noindent\emph{提示：}三个条件分别是 $F_A+[\varphi,\varphi^*]=0$、$\bar\partial_A\varphi=0$ 以及后者的共轭；商去的群是规范群。
\end{enumerate}

\subsubsection*{第 2 讲\quad 4d 主角}

\begin{enumerate}[label=\arabic*.,start=5,leftmargin=2.4em,itemsep=6pt]
\item 证明 $\Ldual{SU(N)}=PSU(N)$，且 $\Ldual{PSU(N)}=SU(N)$。
\par\smallskip\noindent\emph{提示：}对偶群把根系与余根系互换；$SU(N)$ 的中心 $\ZZ_N$ 在对偶群里变成基本群 $\ZZ_N$。

\item 写出 $SU(2)$ 与 $SO(3)$ 的权格与余权格，并验证 GNO 对偶把二者互换。
\par\smallskip\noindent\emph{提示：}$SU(2)$ 的权格是 $(1/2)\ZZ$，$SO(3)$ 的权格是 $\ZZ$；磁荷所取的余权格恰是权格的对偶。

\item 由 $\tau=\theta/2\pi+4\pi\ii/e^2$ 读出 $\Re\tau=\theta/2\pi$、$\Im\tau=4\pi/e^2$，说明 $T:\tau\mapsto\tau+1$ 等价于 $\theta\mapsto\theta+2\pi$，并解释为什么 $\theta$ 的周期必须是 $2\pi$。
\par\smallskip\noindent\emph{提示：}瞬子数取整数值，$\theta$ 的平移只改变 $e^{\ii\theta n}$ 的相位；$T$ 变换就是这一平移。

\item 用 $S=\begin{pmatrix}0&1/\sqrt{n_g}\\-\sqrt{n_g}&0\end{pmatrix}$ 验证 $S:\tau\mapsto-1/(n_g\tau)$，并说明为什么 $n_g=2,3$ 时 $S$ 与 $T$ 生成的是 $SL(2,\RR)$ 的离散子群而不是整个 $SL(2,\ZZ)$。
\par\smallskip\noindent\emph{提示：}直接做矩阵乘法；$n_g$ 不是平方数时 $1/\sqrt{n_g}$ 无理，故 $S\notin SL(2,\ZZ)$。
\end{enumerate}

\subsubsection*{第 3 讲\quad 扭化}

\begin{enumerate}[label=\arabic*.,start=9,leftmargin=2.4em,itemsep=6pt]
\item 由 $\Psi=(\tau t^2+\bar\tau)/(1+t^2)$ 验证：$t=\pm1$ 时 $\Psi=\Re\tau$ 为实；$t=\pm\ii$ 时 $\Psi=\infty$ 且与 $\tau$ 无关；$|t|=1$ 时 $\Psi$ 为实。
\par\smallskip\noindent\emph{提示：}$t=\pm\ii$ 时分母为零；$|t|=1$ 时 $1/t=\bar t$，故 $(t-1/t)/(t+1/t)$ 为纯虚。

\item 设 $S$ 对 $(\tau,t)$ 的作用为 $\tau'=-1/(n_g\tau)$、$t'^2=t^2\tau/\bar\tau$。由 $\Psi$ 的定义验证 $\Psi'=-1/(n_g\Psi)$。
\par\smallskip\noindent\emph{提示：}把 $\Psi'$ 通分后代入 $\tau\bar\tau=|\tau|^2$，并注意 $t'^2\tau'$ 与 $t'^2\bar\tau'$ 的结构。

\item 固定 $\Psi$ 与 $\Re\tau$，令 $\Im\tau\to\infty$，证明 $t^2\to1$，即弱耦合极限对应 $t\to\pm1$。
\par\smallskip\noindent\emph{提示：}$t^2=(\bar\tau-\Psi)/(\Psi-\tau)$；写 $\tau=\Re\tau+\ii\Im\tau$，比较分子分母的最高阶项。

\item 说明为什么 $16$ 个超荷在 $Spin(4)_{space}\times Spin(6)_R$ 下分解、取对角 $Spin(4)$ 后恰有两个不变超荷，并解释为什么扭化族的参数是 $\PP^1$。
\par\smallskip\noindent\emph{提示：}$16=(2,1,4)+(1,2,\bar4)$；把 $Spin(6)$ 的 $4$ 限制到 $Spin(4)$ 后两条链各留下一个单态，两个超荷的复线性组合只差一个整体标度。
\end{enumerate}

\subsubsection*{第 4 讲\quad 紧化与镜像}

\begin{enumerate}[label=\arabic*.,start=13,leftmargin=2.4em,itemsep=6pt]
\item 解释为什么四维理论沿 $C$ 紧化、只保留零模后，二维理论是靶空间为 $\MH$ 的 sigma 模型。
\par\smallskip\noindent\emph{提示：}零模由平坦联络加全纯 Higgs 场刻画，而 BPS 方程正是 Hitchin 方程；sigma 模型的靶空间就是这些模构成的模空间。

\item 在 $\tau$ 为纯虚（$\theta=0$）时，验证 $t=\pm\ii$ 对应 $(w_+,w_-)=(-\ii,-\ii)$、因而是 $\MH(G,C)$ 上复结构 $J$ 的 B-模型；再施加 $S$ 对偶得到 $t=\pm1$ 与 $\Psi=0$，对应 $\MH(\Ldual G,C)$ 上复结构 $K$ 的 A-模型。
\par\smallskip\noindent\emph{提示：}判据是 $J_+=J_-$ 给 B-模型、$J_+=-J_-$ 给 A-模型；$S$ 把 $t=\pm\ii$ 送到 $t=\pm1$。

\item 验证 $S$-对偶把支撑在单点的 zerobrane（型 $(B,B,B)$）送到型 $(B,A,A)$ 的膜，并说明为什么后者的支撑必须是某条 Hitchin 纤维。
\par\smallskip\noindent\emph{提示：}查 Table 2 的变换 $(B,B,B)\mapsto(B,A,A)$；A-膜至少是中维子流形，而 zerobrane 的像维数不超过中维，故恰好落在 Hitchin 纤维上。

\item 说明 Hitchin 纤维 $F$ 与对偶群一侧的 Hitchin 纤维互为对偶环面，并解释 $S$-对偶在纤维上为什么表现为 T-对偶。
\par\smallskip\noindent\emph{提示：}环面上的平坦 $U(1)$ 丛（带平坦联络）与对偶环面的点一一对应；KW 用 SYZ 图景作启发，并未把它当作证明。
\end{enumerate}

\subsubsection*{第 5 讲\quad 膜与线算子}

\begin{enumerate}[label=\arabic*.,start=17,leftmargin=2.4em,itemsep=6pt]
\item 判断下列子流形在 $I,J,K$ 下是复子流形还是 Lagrange 子流形，并据此写出膜类型：单点、整个 $\MH$、Hitchin 纤维 $F$、固定 $E$ 而让 $\varphi$ 变化的族、典范余各向同性膜的支撑 $W$。
\par\smallskip\noindent\emph{提示：}点与整个 $\MH$ 在三个复结构下都复（型 $(B,B,B)$）；Hitchin 纤维与固定 $E$ 的族是 $\Omega_I$-Lagrangian 复子流形（型 $(B,A,A)$）；$W$ 在 $J$ 中复、在 $I,K$ 中余各向同性（型 $(A,B,A)$）。

\item 设 $B_p$ 是支撑在 $p\in\MH$ 的 zerobrane，其 Chan--Paton 对象是点 $p$ 上的天穹层。利用特征膜的定义 $XB=B\otimes V$，说明 Wilson 线算子作用在 $B_p$ 上只能把它变成 $B_p\otimes R$，从而 $B_p$ 是电特征膜的原型。
\par\smallskip\noindent\emph{提示：}线的效应只会改变 Chan--Paton 数据；张量一个向量空间不改变膜的支撑，所以特征值不是数而是向量空间。

\item 说明 't Hooft 算子所带的磁荷为什么取值在余权格中，并解释其奇异性如何由三维单极子解实现。
\par\smallskip\noindent\emph{提示：}GNO 量子化条件把磁荷格定为权格的对偶；万有覆盖上的单极子解在远处给出 Dirac 型奇异性。

\item 阅读 KW 第 6.4 节末段，说明为什么电特征膜与磁特征膜都自动带 $C$ 上的平坦向量丛；再解释这个平坦向量丛为什么被解释为 $\Ldual G$-局部系统而不是 $G$-局部系统。
\par\smallskip\noindent\emph{提示：}$\pi_1(C)$ 的表示给出局部系统；Wilson 算子的标签是 $G$ 的表示 $R$，而特征值随 $R$ 变化的规律由对偶群的局部系统编码。
\end{enumerate}

\subsubsection*{第 6 讲\quad 't Hooft 与 Hecke}

\begin{enumerate}[label=\arabic*.,start=21,leftmargin=2.4em,itemsep=6pt]
\item 对 $G=GL(1)$ 写出点 $x$ 处的 Hecke 修改，并证明只有 $L\mapsto L(x)$ 与 $L\mapsto L(-x)$ 两种。
\par\smallskip\noindent\emph{提示：}在 $C\setminus x$ 上同构、在 $x$ 处允许一阶零点或一阶极点；秩 $1$ 时局部数据只由 $\ZZ$ 的一个整数标记。

\item 写出 Hecke 特征层的定义式，并说明它与特征膜定义 $XB=B\otimes V$ 为什么形式相同；解释这个形式上的对应在 KW 词典中意味着什么。
\par\smallskip\noindent\emph{提示：}Hecke 算子把层送到自身张量局部系统的纤维，正如线算子把膜送到自身张量一个向量空间；两者都是「本征值不是数而是空间」的例证。

\item 把四维自对偶方程沿一个平移方向约化，得到三维的 Bogomolny 方程，并说明为什么 't Hooft 算子沿 $C$ 的奇异性对应三维单极子。
\par\smallskip\noindent\emph{提示：}把被约化方向的分量解释为标量 $\varphi$，自对偶条件化为 $F=\ii\star D\varphi$；有限能量解在空间无穷远处的渐近行为就是单极子奇异性。

\item 说明为什么点 $x$ 处的 Hecke 修改由仿射 Grassmannian $G_\CC((z))/G_\CC[[z]]$ 参数化，并说明其 $G_\CC[[z]]$-轨道由支配余权标记。
\par\smallskip\noindent\emph{提示：}在 $x$ 附近取形式坐标，修改的数据是形式圆盘上的一个转移函数；按极点阶数做 Birkhoff 分解即得轨道分解。
\end{enumerate}

\subsubsection*{第 7 讲\quad 翻译器}

\begin{enumerate}[label=\arabic*.,start=25,leftmargin=2.4em,itemsep=6pt]
\item 从典范余各向同性膜（型 $(A,B,A)$）的边界条件出发，解释为什么其开弦代数（$(B_{\mathrm{c.c.}},B_{\mathrm{c.c.}})$ 弦）是作用在 $K_M^{1/2}\otimes L$ 上的微分算子层。
\par\smallskip\noindent\emph{提示：}$J$ 中的 B-型条件给出全纯性，$I,K$ 中的 A-型条件把一半的振荡模消掉；量子化后剩下的态是乘子与微分的组合，即微分算子，$K^{1/2}$ 因子来自测度与反常。

\item 已知 $f$ 是线性的，且对整数 $\Psi=n$ 有 $f(n)=n$。由此证明 $f(\Psi)=\Psi$。
\par\smallskip\noindent\emph{提示：}设 $f(\Psi)=a\Psi+b$，对所有整数 $n$ 代入 $f(n)=n$，比较两个整数点即可。

\item 用 $f(\Psi)=\Psi$ 解释 c.c.\ 膜的弦代数如何随 $\Psi$ 改变；说明为什么 $\Psi=0$ 时它退化为 $K_M^{1/2}$ 上通常 $D$-模的扭化版本。
\par\smallskip\noindent\emph{提示：}$f(0)=0$ 时 $L^{f(\Psi)}=\mathcal O$，只剩 $K_M^{1/2}$ 因子；$\Psi\neq0$ 时多出线丛 $L^\Psi$ 的扭化。
\end{enumerate}

\subsubsection*{第 8 讲\quad 收获与边界}

\begin{enumerate}[label=\arabic*.,start=28,leftmargin=2.4em,itemsep=6pt]
\item 列出 KW 论证链中依赖以下物理假设的步骤：(1) 扭化后的拓扑场论存在且在 $S$-对偶下等价；(2) 膜与 Hecke 特征层的词典；(3) 强耦合与弱耦合之间的解析延拓。说明哪一条最难完全数学化。
\par\smallskip\noindent\emph{提示：}见第 8 讲的总结表；难点集中在把 $S$-对偶写成数学上严格的等价，以及膜范畴与 $D$-模范畴的比较。

\item$^\dagger$ 2024 年 Gaitsgory 等人的工作是纯代数几何的证明。说明它证明了什么、没有证明什么（是否包含 $S$-对偶？是否覆盖分歧情形？），并解释为什么这不妨碍 KW 的物理论证仍然有价值。
\par\smallskip\noindent\emph{提示：}该证明处理特征 $0$、非分歧的 de Rham 情形，证明 Hecke 特征层存在；它不构造物理的 $S$-对偶，也不覆盖分歧与量子化参数的全部范围。

\item 若只考虑非分歧情形，KW 的构造需要哪些额外限制（如局部系统在每点局部平凡、't Hooft 奇异性无冒泡）？要推广到分歧与野分歧，需要补充什么数据？
\par\smallskip\noindent\emph{提示：}分歧情形要在点上加入规范场奇异性的数据（对应抛物结构或 Level 结构），野分歧还需要高阶极点/幂零单值性的信息。
\end{enumerate}

\section{参考文献与阅读路线}

建议先用本讲义第 1--4 讲把舞台、主角、扭化与紧化这条主线走一遍，建立 $\tau,\Psi,t$ 与三种复结构之间的对应；再用第 5--7 讲把膜、线算子与 $D$-模的词典串起来。之后回到原始文献：先读 \cite{kw06} 的 \S1--\S5（约前三分之一），对照讲义逐个核对符号；再读 \S6--\S12，留意 't Hooft 算子的冒泡与分歧情形。物理背景可补读 \cite{w08} 的镜像与 Hitchin 视角，数学动机读 \cite{fr09} 的 Bourbaki 报告，分歧情形读 \cite{gw06}。最后可浏览 \cite{set13} 的博士论文与 \cite{bz04} 的课程笔记作为二次文献。

\begin{thebibliography}{9}
\bibitem{kw06} A.~Kapustin and E.~Witten, \emph{Electric--Magnetic Duality and the Geometric Langlands Program}, Commun.\ Number Theory Phys.\ 1 (2007) 1--236; arXiv:hep-th/0604151.
\bibitem{w08} E.~Witten, \emph{Mirror Symmetry, Hitchin's Equations, and Langlands Duality}, arXiv:0802.0999.
\bibitem{fr09} E.~Frenkel, \emph{Gauge Theory and Langlands Duality}, S\'eminaire Bourbaki 1010, Ast\'erisque 332 (2010) 369--403; arXiv:0906.2747.
\bibitem{gw06} S.~Gukov and E.~Witten, \emph{Gauge Theory, Ramification, and the Geometric Langlands Program}, arXiv:hep-th/0612073.
\bibitem{bz04} D.~Ben-Zvi, \emph{Between Electric--Magnetic Duality and the Langlands Program}, 课程笔记（arXiv 与作者主页可查）.
\bibitem{set13} N.~Setter, \emph{Topological Quantum Field Theory and the Geometric Langlands Program}, 博士论文.
\bibitem{ht03} N.~Hitchin, \emph{The Self-Duality Equations on a Riemann Surface}, Proc.\ London Math.\ Soc.\ 55 (1987) 59--126.
\end{thebibliography}


\end{document}
